% SPDX-FileCopyrightText: © 2026 Idiap Research Institute <contact@idiap.ch>
%
% SPDX-FileContributor: Andre Freitas <andre.freitas@idiap.ch>
% Neuro-symbolic AI Group
%
% SPDX-FileContributor: Danilo Gusicuma <danilo.gusicuma@idiap.ch>
% Neuro-symbolic AI Group
%
% SPDX-License-Identifier: MIT
\documentclass[11pt]{article}

\usepackage[letterpaper,margin=1in]{geometry}
\usepackage[T1]{fontenc}
\usepackage{lmodern}
\usepackage{microtype}
\usepackage{amsmath,amssymb,amsthm,mathtools}
\usepackage{booktabs,array,tabularx,longtable}
\usepackage{float}
\usepackage{enumitem}
\usepackage{xcolor}
\usepackage[colorlinks=true,linkcolor=blue!55!black,citecolor=blue!55!black,urlcolor=blue!55!black]{hyperref}
\hypersetup{pdftitle={The Formula-Footprint Barrier for Label-Induced Hardness Predicates},pdfauthor={Anonymous}}
\usepackage[nameinlink,capitalise,noabbrev]{cleveref}

\setlist[itemize]{leftmargin=*,itemsep=0.35em,topsep=0.35em}
\setlist[enumerate]{leftmargin=*,itemsep=0.35em,topsep=0.35em}
\allowdisplaybreaks

\newtheorem{theorem}{Theorem}[section]
\newtheorem{lemma}[theorem]{Lemma}
\newtheorem{proposition}[theorem]{Proposition}
\newtheorem{corollary}[theorem]{Corollary}
\theoremstyle{definition}
\newtheorem{definition}[theorem]{Definition}
\newtheorem{example}[theorem]{Example}
\newtheorem{remark}[theorem]{Remark}
\newtheorem{openproblem}[theorem]{Open Problem}

\newcommand{\bits}{\{0,1\}}
\newcommand{\F}{\mathbb{F}_2}
\newcommand{\SIZE}{\operatorname{SIZE}}
\newcommand{\Ckt}{\operatorname{C}}
\newcommand{\SAT}{\mathrm{SAT}}
\newcommand{\Ppoly}{\mathrm{P/poly}}
\newcommand{\poly}{\operatorname{poly}}
\newcommand{\polylog}{\operatorname{polylog}}
\newcommand{\dist}{\operatorname{dist}}
\newcommand{\supp}{\operatorname{supp}}
\newcommand{\wt}{\operatorname{wt}}
\newcommand{\Eq}{\operatorname{EQ}}
\newcommand{\Lookup}{\operatorname{Lookup}}
\newcommand{\Force}{\operatorname{force}}
\newcommand{\eps}{\varepsilon}
\newcommand{\elln}{\ell_n}
\newcommand{\Mn}{M_n}
\newcommand{\Gn}{\mathbf{G}_n}
\newcommand{\Un}{U_n}
\newcommand{\Om}{\Omega_n}
\newcommand{\HardF}{\mathsf{Hard}^{\mathrm{form}}}
\newcommand{\HardO}{\mathsf{Hard}^{\mathrm{occ}}}
\newcommand{\HardQ}{\mathsf{Hard}^{\mathrm{quot}}}
\newcommand{\EasyF}{\mathsf{Easy}^{\mathrm{form}}}
\newcommand{\pospart}[1]{\left(#1\right)_{+}}

\title{The Formula-Footprint Barrier\\
\large for Label-Induced Hardness Predicates}
\author{Anonymous}
\date{}

\begin{document}
\maketitle

\begin{abstract}
Let a batched generator map each source input $x\in\bits^n$ to $M_n$ short strings, and let an external Boolean labeling of those strings induce a source Boolean function through a fixed row aggregator.  The formula-footprint barrier says that any completion-sound certificate that this induced function has circuit complexity greater than \(s_n\) must expose many generated labels of both polarities.  We study this barrier in the standard nondeterministic decision-tree sense of inputwise \(1\)-certificates.  The input variables may be labels of distinct generated strings, independent labels of generated occurrences, or quotient labels after sound equality identifications.

The main result is a formula-label footprint lower bound.  A sorting-network dictionary evaluates a batch of $M$ variable keys against $K$ hardwired exceptional keys in size
\[
  \widetilde O\bigl((M+K)(\ell+\log(M+K))\bigr).
\]
Consequently, every nondeterministic decision-tree \(1\)-certificate for hardness above $s_n$ must contain, separately, a large set of generated labels fixed to $0$ and a large set fixed to $1$.  In the regime $s_n=n^{1+\eps-o(1)}$, $M_n=n^{q+o(1)}$, and generator size $n^{\rho+o(1)}$ with $\max\{q,\rho\}<1+\eps$, each polarity has size $n^{1+\eps-o(1)}$.

This lower bound is close to best possible in the standard identity truth-table carrier.  A circuit-counting test-set argument shows that a generic function of circuit complexity above \(s\) has a completion-sound certificate of size \(O(s\log(s+n))\), with both polarities represented.  Thus the formula-footprint theorem is not merely a necessary condition; in the baseline model it captures the right certificate scale up to polylogarithmic factors.

Two supporting structural lemmas record what changes under different label granularities.  In the private occurrence model, sparse-row patching shows that a sound certificate must force $\Omega(s_n/n)$ rows to each output value.  If $\Force_b(A_M)$ is the minimum number of coordinates needed by any partial assignment forcing an aggregator $A_M$ to $b$, then every certificate queries
\[
  \Omega\!\left((\Force_0(A_{M_n})+\Force_1(A_{M_n}))\frac{s_n}{n}\right)
\]
occurrences.  For parity, $\Force_0=\Force_1=M_n$, so every forced row is complete and the query lower bound is $\Omega(M_ns_n/n)$.

For parity after sound equality identifications, the residual completion space is an affine image, and soundness requires every completion to stay at Hamming distance $\Omega(s_n/n)$ from both constant functions.  For nonlinear semantic relations, support locality forces an ordinary relation-membership certificate on each constrained tuple; in particular, an unsatisfiable SAT self-reduction split \(y=y_0\vee y_1\) must expose parent and both children.  For an acyclic OR self-reduction tree after internal nodes are eliminated, the corresponding frontier relation has an exact global certificate: root \(0\) forces the root and all frontier leaves, while root \(1\) needs the root and one true leaf.  The same exact certificate holds for finite acyclic OR DAG frontiers after repeated reachable frontier formulas are soundly quotiented: sharing changes raw leaf occurrences into distinct reachable frontier classes, but not the residual OR certificate.  More generally, any exposed finite acyclic Boolean frontier circuit has residual support cost equal to the root coordinate plus ordinary pointwise certificate complexity for the induced frontier function.  Finally, bounded-error randomized verifiers inherit the same lower bounds for the union of labels queried on accepting executions: that union is itself a deterministic certificate whenever the acceptance probability on the hard labeling exceeds the soundness bound for easy completions.  A displayed-atom extraction theorem applies the formula-footprint bound to any proof interface whose accepted proof explicitly displays a faithful, completion-sound set of formula labels.  Restricted resolution theorems apply this to explicit label assumptions, to the used unit core of full-assignment refutations, to the support of explicit local label clauses, and to auxiliary encodings that are support-local under all completions.  For extension proofs, the paper gives both a declared dependency-cone normalizer and a canonical syntactic-cone normalizer extracting the same support directly from the extension-definition graph; the final backward-slice corollary records the corresponding proof-object budget, a slice-explicit extension-resolution corollary turns that budget into a restricted ER-style size lower bound, a displayed acyclic ER corollary shows that ordinary displayed extension clauses already charge the same boundary, an expansion-bounded compressed ER transfer states the exact decompression budget needed for compressed representations, a macro-compressed ER corollary turns ordinary macro reuse into an explicit reuse-product lower bound, a support-declared macro ER corollary charges any sound declared support budget, and a support-audited LRAT corollary packages the same obligation for checker-visible clausal traces.  A provenance-typed LRAT corollary shows that a trusted variable-role map plus block-local clause provenance supplies such a support audit, and a namespace-typed LRAT corollary shows that, with reserved formula-label variables, the LRAT antecedent closure of the empty clause already identifies the used label support when the used auxiliary clauses have a hard witness.  A strict-prefix Qrep replay proposition adds the quotient-compression version of the same rule: representative reuse is transparent only when an already charged same-theory certificate forest can replay it.  A trace-complete DPLL decision-log corollary gives the first DPLL-style proof-object instantiation of the hidden-adaptive audit, and an antecedent-complete CDCL corollary adds learned clauses when their acyclic antecedent DAGs expose all local support used by the replay.  A CDCL antecedent-erasure trichotomy records the matching boundary: if learned-clause antecedents and role provenance are erased, a support extractor must add equivalent provenance, pay conservative support, or be unsound on a same-skeleton lift.  A RUP-complete CDCL corollary gives the positive non-plain bridge: checker-visible RUP propagation certificates over a known live clause database, together with trusted label namespace or role parsing and retained deleted-clause support, recover the support audit without requiring an explicit learned-clause antecedent DAG.  A canonical-namespace RUP-CDCL corollary then makes the namespace route checkable from a deterministic encoding manifest with a reserved label-variable interval and live-database reconstruction.  A support-audit necessity separator then shows why the LRAT audit is a real invariant rather than a harmless annotation: checked clausal unsatisfiability alone can coexist with an unsound small support declaration.  A support-erasure separator shows the provenance version of the same boundary: the same untyped LRAT skeleton may have an auxiliary lift with empty formula-label support and a label lift requiring all labels.  A plain LRAT extraction trichotomy packages the raw-trace route: add provenance, pay role-agnostic support, or be unsound on a same-skeleton semantic lift.  A support-local normalization transfer theorem then gives the general proof-size consequence for any proof system that admits such a normalizer with controlled support budget.  These are certificate-complexity statements.  They do not prove $P\ne NP$, do not prove $\SAT\notin\Ppoly$, do not construct an explicit hard SAT-labeled carrier, and do not imply a lower bound for Frege, Extended Frege, IPS, raw RUP/DRUP without a manifest or support parsing, arbitrary CDCL without provenance or a certified RUP support audit, or arbitrary compressed proof systems without supplying the required normalization theorem.
The carrier-independent audit upgrade criterion isolates the remaining top-paper
proof-system target: a proof-system-intrinsic extractor, not a carrier-specific
manifest or hidden role annotation, would instantiate the transfer theorem and
yield the same footprint lower bound with its audit budget.
A self-auditing RUP-CDCL instantiation shows that this target is nonvacuous:
if the proof syntax itself contains the checked label decoder, live-database
reconstruction, RUP propagation witnesses, and retained deleted-clause support,
then the audit is computed from the proof artifact itself.  An append-only
self-auditing variant removes the retained deleted-clause support field by
making the clause database monotone.  A prefix-typed append-only variant removes
the proof-internal decoder as well when formula-label variables are a public
syntactic class of the proof system.  These remain typed certificate-format
theorems, not claims about raw DIMACS/RUP/DRUP
syntax.
\end{abstract}

\section{Introduction}

Hardness magnification asks whether a lower bound only slightly beyond an accessible frontier can imply a major complexity separation.  The transfer is useful only when its quantitative losses are controlled.  The proof of the source lower bound then becomes the central issue: which kinds of local evidence could certify that the source function is hard, and what information must such evidence expose?

This paper isolates a certificate-complexity question that arises naturally in this setting.  A batched generator produces a list of short objects from a source input.  An external Boolean labeling assigns a bit to each object, and a fixed aggregator combines the labels in each source row.  The labeling therefore induces an $n$-input Boolean function.  We ask for nondeterministic decision-tree certificates of the predicate
\[
  \text{``the induced source function has circuit complexity greater than $s_n$.''}
\]
The key point is that this is an ordinary Boolean property of the labels.  A sound local transcript is exactly a $1$-certificate: every completion of the unexposed labels must still induce a hard function.

The contribution is deliberately narrower than an unrestricted proof-system lower bound or an explicit SAT lower bound.  The main technical point is the \emph{formula-footprint barrier}: distinct generated formulas share a single static dictionary, so a small two-polarity formula footprint gives a cheap completion circuit.  The row-forcing and quotient-distance results are structural tests for local proof interfaces; they are not claimed as lower bounds for arbitrary proof systems.  The SAT parity carrier is used only as a conditional calibration: if an explicit source lower bound were supplied, local completion-sound certificates would have to pay the costs identified here.

We analyze three granularities of label variables.

\begin{itemize}
  \item In the \emph{formula-label model}, equal generated strings share one variable.  A certificate may exploit all collisions in the generator range.
  \item In the \emph{occurrence-label model}, every generated position $(x,j)$ has an independent variable.  This relaxation is stronger adversarially and makes row structure explicit.
  \item In the \emph{quotient-label model} for parity, sound equality links identify labels before the row parities are formed.  The residual completion space is an affine image under a binary incidence map.
\end{itemize}

Across these models, the same proof strategy appears in different forms:

\begin{quote}
A sound certificate cannot admit a completion whose induced source function has a cheap circuit.
\end{quote}

For formula labels, a cheap completion is nearly constant on the generated formula universe.  For occurrence labels, it is constant outside a small set of forced rows.  For equality quotients, it is an affine completion close to a constant source truth table.

\subsection{Conceptual spine: one obstruction in three interfaces}

The results are meant to be read as one cohered theorem package, not as a fragmented ensemble of nearby observations or disconnected payoffs.  The organizing obstruction is \emph{completion soundness}: after a local proof, verifier, or certificate exposes some labels, all unexposed labels remain adversarially completable.  A certificate is strong only if no such completion makes the induced source function easy.

The sharp result carrying the paper is the formula-footprint theorem.  Before this theorem, a local hardness certificate for a generator-induced labeling could be viewed as an arbitrary sparse partial truth table.  After the theorem, a small two-polarity footprint has a concrete consequence: the unexposed labels can be set to a constant, the exposed exceptions can be evaluated by a shared batched dictionary, and the resulting source function has a circuit of size nearly linear in the batch size plus the footprint.  Thus, near the magnification threshold, every completion-sound formula-local certificate, equivalently every nondeterministic decision-tree \(1\)-certificate for this hardness predicate, must expose \(n^{1+\eps-o(1)}\) generated formulas of each polarity.

The row-forcing and quotient-distance theorems are not competing main results.  They explain how the same obstruction changes when the local interface changes.  Private occurrence access turns completion soundness into many forced source rows; parity then makes each forced row complete.  Equality links turn completion soundness into an affine residual-distance condition.  Nonlinear semantic relations, such as one-step SAT self-reduction, are not freely available under arbitrary completion: \Cref{prop:semantic-relation-support} shows that a support-local audit must contain an ordinary membership certificate for each relation instance, \Cref{prop:acyclic-or-tree-residual} gives the exact residual certificate for acyclic OR self-reduction trees after eliminating internal nodes, \Cref{prop:or-dag-frontier-residual} extends that exact certificate to finite acyclic OR DAG frontiers after sound quotienting of repeated reachable frontier formulas, and \Cref{prop:acyclic-boolean-frontier-residual} gives the general exposed acyclic Boolean frontier version by reducing support-locality to ordinary pointwise certificate complexity of the induced frontier function.  The identity-carrier sharpness theorem shows that the formula footprint scale is essentially optimal in the baseline truth-table model.  The displayed-atom theorem, the restricted resolution theorems, the support-local extension-resolution instantiation, the dependency-cone normalizer, canonical syntactic-cone normalizer, final backward-slice lower bound, slice-explicit extension-resolution lower bound, displayed acyclic ER length lower bound, expansion-bounded compressed ER transfer, macro-compressed ER reuse tradeoff, support-declared macro ER normalizer, support-audited LRAT trace normalizer, provenance-typed LRAT support audit, namespace-typed LRAT antecedent-closure normalizer, strict-prefix Qrep replay normalizer, trace-complete DPLL decision-log normalizer, antecedent-complete CDCL learned-clause normalizer, CDCL antecedent-erasure trichotomy, RUP-complete CDCL support normalizer, canonical-namespace RUP-CDCL standardizer, self-auditing RUP-CDCL instantiation, append-only self-auditing RUP-CDCL instantiation, prefix-typed append-only RUP-CDCL instantiation, support-audit necessity separator, support-erasure indistinguishability separator, plain LRAT extraction trichotomy, cone-budget corollary, bounded-depth tradeoff, and the transfer theorem in \Cref{thm:displayed-atom-extraction,thm:assumption-resolution,cor:resolution-unit-core,thm:resolution-local-clause-support,thm:support-local-aux-resolution,thm:support-local-extension-resolution,thm:dependency-cone-extension-normalizer,cor:canonical-syntactic-cone,cor:final-backward-slice,cor:slice-explicit-er,cor:displayed-er-length,cor:compressed-er-expansion,cor:macro-compressed-er,cor:support-declared-macro-er,cor:support-audited-lrat,cor:provenance-typed-lrat,cor:namespace-typed-lrat,prop:strict-prefix-qrep,cor:trace-complete-dpll,cor:antecedent-complete-cdcl,prop:cdcl-antecedent-erasure-trichotomy,cor:rup-complete-cdcl,cor:canonical-namespace-rup-cdcl,cor:self-auditing-rup-cdcl,cor:append-only-self-auditing-rup-cdcl,cor:prefix-typed-append-only-rup-cdcl,prop:support-audit-necessity,prop:support-erasure-indistinguishability,prop:plain-lrat-extraction-trichotomy,cor:cone-budget-extension,cor:bounded-depth-extension,thm:support-local-normalization-transfer} are the proof-complexity bridge for this package: a proof-system application must extract, with controlled overhead, one of these completion-sound local objects or expose a syntactic dependency-cone, typed provenance or namespace package, strict-prefix replay dictionary, trace-complete DPLL controller log, antecedent-complete CDCL learned-clause DAG, RUP-complete propagation-trace support audit, canonical label-interval encoding manifest, self-auditing label decoder, append-only RUP audit, prefix-typed label namespace, non-plain CDCL semantic invariant, or sound support-declaration budget for its auxiliary encoding.

The carrier-independent audit upgrade criterion is the same proof obligation in
admission-rule form.  It says exactly what a stronger recognizable proof-system
payoff must contribute: the support audit must be forced by the proof-system
artifact itself, not supplied by a carrier-specific manifest or hidden role
advice.  This keeps the proof-system route attached to the flagship obstruction
rather than becoming a fragmented list of proof-format examples.
\Cref{cor:self-auditing-rup-cdcl} supplies the corresponding positive
instantiation for an enriched but checkable RUP-CDCL certificate format: the
label decoder and support audit are part of the proof object, so the audit is
intrinsic to the proof system rather than imported from a side manifest.
\Cref{cor:append-only-self-auditing-rup-cdcl} removes one annotation layer in
the deletion-free case: no retained deleted-clause support records are needed
when the live database is append-only.  \Cref{cor:prefix-typed-append-only-rup-cdcl}
removes another: no proof-internal decoder is needed when the proof system has a
public formula-label variable namespace.

This changes the proof obligation.  The paper does not say that a major lower bound follows from locality alone.  It says that a local route to such a lower bound must either pay the formula-footprint, row-forcing, or quotient-distance cost, or explain why the proof interface escapes those completion models.

\subsection{The formula-footprint barrier theorem}

Let
\[
  \Gn(x)=\bigl(G_{n,1}(x),\ldots,G_{n,M_n}(x)\bigr)
\]
be a single multi-output circuit of size $g_n$, with each output string of length $\ell_n$.  Let $U_n$ be the set of distinct generated strings.  For a label table $\lambda:U_n\to\bits$ and an aggregator $A_{M_n}$, define
\[
  H_{n,A}^{\lambda}(x)
  =A_{M_n}\bigl(\lambda(G_{n,1}(x)),\ldots,
                    \lambda(G_{n,M_n}(x))\bigr).
\]
The formula-label hardness predicate is $1$ precisely when $\Ckt(H_{n,A}^{\lambda})>s_n$.

Suppose a certificate fixes $K_b$ distinct generated formulas to the value $b$.  Complete every other formula to the opposite value $1-b$.  The resulting label table differs from the constant $1-b$ on exactly those $K_b$ formulas.  The naive circuit compares every generated query against every exception and costs $\Theta(M_nK_b\ell_n)$.  The useful observation is that the dictionary is shared by the entire batch.  Sorting the $M_n$ variable query records together with the $K_b$ hardwired records gives a circuit of size
\[
  O\!\left(
    (M_n+K_b)(\ell_n+\log(M_n+K_b))\log^2(M_n+K_b)
  \right).
\]
Thus a sound certificate must satisfy, separately for $b=0$ and $b=1$,
\[
  g_n+\widetilde O\bigl((M_n+K_b)\ell_n\bigr)>s_n.
\]
In particular, if $s_n=n^{a-o(1)}$, $M_n=n^{q+o(1)}$, $g_n=n^{\rho+o(1)}$, $\ell_n=n^{o(1)}$, and $a>\max\{q,\rho\}$, then
\[
  K_0,K_1\ge n^{a-o(1)}.
\]
This is an inputwise $1$-certificate lower bound, not an average-case or naturalness statement.  It respects every equality among repeated generated formulas.

\subsection{The row-forcing theorem}

Now give each occurrence $(x,j)$ an independent label $z(x,j)$.  A partial assignment on row $x$ \emph{forces} output $b$ if every completion of the unqueried positions has aggregator value $b$.  Let $R_b$ be the set of rows forced to $b$.

If a row is not forced to $b$, it has a completion to $1-b$.  Complete every such row independently in this way.  The resulting source function differs from the constant $1-b$ only on $R_b$, and equality tests patch those rows with $O(n|R_b|)$ gates.  Therefore every hardness certificate above $s_n$ must satisfy
\[
  |R_0|,|R_1|=\Omega(s_n/n).
\]

Let $\Force_b(A_M)$ be the minimum size of a partial assignment that forces $A_M=b$.  Every row in $R_b$ contains at least $\Force_b(A_M)$ queried coordinates, giving
\[
  |Q|=\Omega\!\left(
      (\Force_0(A_{M_n})+\Force_1(A_{M_n}))\frac{s_n}{n}
    \right).
\]
For parity, both forcing numbers equal $M_n$.  Hence all forced rows are complete and
\[
  |Q|=\Omega(M_ns_n/n).
\]
At $M_n=n^{q+o(1)}$ and $s_n=n^{1+\eps-o(1)}$, this becomes $n^{q+\eps-o(1)}$ queried occurrences.

The theorem also explains the role of the aggregator.  OR has a one-bit certificate for output $1$ but needs a full row for output $0$; parity has full-row certificate complexity on both sides.  The theorem records the two forcing costs separately rather than hiding this asymmetry.

\subsection{Equality quotients}

Occurrence independence is invalid once a certificate uses sound equality links.  If two formulas are known to be equisatisfiable, their labels agree; in a parity row, two occurrences from the same class may cancel without either value being fixed.  The correct equality-aware model quotients the generated formula universe by a label-respecting equivalence relation $E_n$.

For parity, let $L_{E_n}$ be the binary incidence map from quotient labels to source truth tables.  Fixing some quotient labels leaves an affine set $\mathcal A$ of possible source truth tables.  We prove that completion soundness above $s_n$ implies
\[
  \min_{h\in\mathcal A}\min_{b\in\bits}
  \dist(h,b\mathbf 1)
  =\Omega(s_n/n).
\]
Otherwise one affine completion is close to a constant function and can be patched by a small circuit.  This is the equality-aware replacement for the private-occurrence completion argument.  It does not normalize nonlinear semantic constraints such as SAT self-reduction.

\subsection{Application to polylogarithmic SAT parity carriers}

The motivating application is a generator of SAT instances of length
\[
  \ell_n=\lceil(\log n)^d\rceil,
  \qquad d>1,
\]
with $M_n=n^{q+o(1)}$ outputs and total generation cost $g_n=n^{\rho+o(1)}$.  Define
\[
  H_n^{\SAT}(x)
  =\bigoplus_{j=1}^{M_n}\SAT(G_{n,j}(x)).
\]
If SAT has circuits of size $m^c$, substituting one such circuit at every generated position gives
\[
  \Ckt(H_n^{\SAT})
  \le g_n+M_n\ell_n^c+O(M_n)
  =n^{\max\{q,\rho\}+o(1)}.
\]
The target exponent $c$ disappears because $\ell_n=n^{o(1)}$.  Thus an infinitely-often source lower bound
\[
  \Ckt(H_n^{\SAT})>n^{1+\eps-o(1)},
  \qquad \max\{q,\rho\}<1+\eps,
\]
would imply $\SAT\notin\Ppoly$.

Our certificate theorems do not prove this source lower bound.  They show what several local proof interfaces would have to pay in order to certify it.  At the threshold above, the costs are summarized in \Cref{tab:intro-summary}.

\begin{table}[htbp]
\centering
\small
\begin{tabularx}{\textwidth}{@{}>{\raggedright\arraybackslash}p{0.19\textwidth}
>{\raggedright\arraybackslash}p{0.34\textwidth}
>{\raggedright\arraybackslash}X@{}}
\toprule
Label granularity & Completion model & Necessary resource at $s_n=n^{1+\eps-o(1)}$ \\
\midrule
Distinct formulas & Arbitrary labels on unexposed generated strings &
$n^{1+\eps-o(1)}$ exposed generated formulas of each polarity \\
Independent occurrences & Arbitrary completion of each position $(x,j)$ &
$\Omega(M_nn^{\eps-o(1)})=n^{q+\eps-o(1)}$ probes for parity; $n^{\eps-o(1)}$ complete rows of each parity \\
Equality quotient & Arbitrary completion of unexposed quotient labels &
Residual affine source set at distance $n^{\eps-o(1)}$ from both constant truth tables \\
\bottomrule
\end{tabularx}
\caption{The three completion barriers for a near-threshold SAT parity carrier.}
\label{tab:intro-summary}
\end{table}

\subsection{Status of the theorem package}

\Cref{tab:theorem-status} records the role of each claim in the revised paper.  This table is part of the mathematical scoping, not a meta-commentary: later sections use the same boundaries.  In particular, the formula-footprint theorem is the main technical result, while the row-forcing and quotient-distance theorems are supporting structural lemmas.  The SAT carrier discussion is conditional calibration, and the displayed acyclic ER consequence remains a controlled local proof-object theorem rather than an unrestricted lower bound for stronger proof systems.

The admission rule is that every result must play one of five roles in the same arc: flagship theorem, calibration, interface variant, proof-system bridge, or nonclaim boundary.  A new lemma is not admitted merely because it is locally correct.  It must strengthen the formula-footprint spine, calibrate its sharpness, extend its proof-system reach with controlled overhead, or clarify why a stronger proof-system or SAT-family claim has not yet been proved.  This is the paper's cohered-results criterion: the theorem package should not become a fragmented ensemble of results, and technically correct side payoffs are quarantined unless they serve the same before/after understanding shift.

{\scriptsize
\begin{longtable}{@{}p{0.24\textwidth}p{0.27\textwidth}p{0.43\textwidth}@{}}
\caption{Claim status after the referee-objection closure pass.}
\label{tab:theorem-status}\\
\toprule
Claim & Status in this paper & Boundary \\
\midrule
\endfirsthead
\toprule
Claim & Status in this paper & Boundary \\
\midrule
\endhead
Batched dictionary lemma & Patched primitive & Query records and dictionary records are tagged; dictionary records are masked out before aggregation. \\
Formula-footprint barrier theorem & Main technical result & Gives polarity-refined lower bounds for nondeterministic decision-tree certificates of the induced hardness predicate. \\
Identity-carrier generic sharpness & Sharpness theorem & Matches the formula-footprint scale up to polylogarithmic factors in the standard truth-table carrier. \\
Row-forcing theorem & Supporting structural lemma & Applies to completion-sound private occurrence certificates; not a proof-system theorem. \\
Equality quotient theorem & Supporting structural lemma & Handles sound equality identifications for parity; not nonlinear SAT semantics. \\
Semantic-relation support certificate & Interface-boundary proposition & Nonlinear relation constraints require ordinary membership certificates inside the declared support; this is not yet a global SAT residual-completion theorem. \\
Acyclic OR-tree residual certificate & Bounded semantic-frontier proposition & After eliminating internal OR nodes, root \(0\) forces the root and all frontier leaves, while root \(1\) needs the root and one true leaf; shared pure-OR frontiers are handled by the next row. \\
OR-DAG frontier residual certificate & Bounded shared-frontier proposition & For finite acyclic OR DAGs after sound frontier quotienting, root \(0\) forces the root and every distinct reachable frontier class, while root \(1\) needs the root and one true reachable frontier class; exposed non-OR frontiers are handled by the Boolean-frontier row. \\
Acyclic Boolean frontier residual certificate & Exposed semantic-frontier proposition & For any exposed finite acyclic Boolean frontier circuit, support-locality is exactly the root coordinate plus an ordinary pointwise Boolean certificate for the induced frontier function; visible adaptivity is handled by the next row. \\
Visible adaptive frontier residual certificate & Exposed adaptive-frontier proposition & For any checker-visible finite adaptive decision tree over frontier coordinates, support-locality is exactly the root coordinate plus an ordinary pointwise Boolean certificate for the induced decision-tree function; the actual transcript path is sufficient support, while shared-state adaptivity is handled by the next row. \\
Visible adaptive branching-program frontier residual certificate & Exposed shared adaptive-frontier proposition & For any checker-visible finite acyclic branching program over frontier coordinates, support-locality is the root coordinate plus an ordinary pointwise Boolean certificate for the induced branching-program function; the actual computation path is sufficient support, while hidden adaptive control remains outside scope. \\
Hidden adaptive transcript insufficiency separator & Boundary proposition & A reported adaptive transcript that is not certified complete can be compatible with both a one-variable frontier function and an all-frontier OR function at the same hard tuple; transcript-only support extraction is therefore unsound unless it exposes the controller/dependencies or pays worst-case support. \\
Controller-complete adaptive audit normalizer & Proof-system bridge proposition & A hidden adaptive controller can be normalized when the proof object exposes an acyclic controller, certifies local equivalence on the claimed completion cylinder, and pays for the replay path support; unaudited transcript summaries remain outside scope. \\
Trace-complete DPLL decision-log normalizer & Recognizable proof-object bridge & A DPLL-style proof object supplies the controller-complete audit when it logs the full finite decision controller, the replay path, and clause or propagation witnesses certifying local equivalence on the claimed completion cylinder; partial solver transcripts and unaudited learned clauses remain outside scope. \\
Antecedent-complete CDCL learned-clause normalizer & Recognizable learned-clause bridge & A CDCL-style proof object supplies the same audit when every learned clause used by the replay has an acyclic antecedent DAG down to original local clauses, decisions, propagations, or already audited learned clauses; arbitrary CDCL traces without antecedent provenance remain outside scope. \\
CDCL antecedent-erasure trichotomy & Boundary for arbitrary CDCL logs & If the learned-clause antecedents and role provenance are erased, a support extractor must either add equivalent provenance or a semantic invariant, pay the conservative role-agnostic support budget, or be unsound on a same-skeleton lift; arbitrary provenance-free CDCL syntax is not a free normalizer. \\
RUP-complete CDCL support normalizer & Positive non-plain CDCL bridge & A CDCL/DRUP-style log supplies the audit when every used learned clause has a checker-visible RUP propagation certificate over a known live clause database, with label namespace or role parsing and retained support for deleted clauses; raw RUP syntax without support parsing remains outside scope. \\
Canonical-namespace RUP-CDCL standardizer & Checkable encoding bridge & A deterministic encoding manifest can replace trusted role advice when it reserves a formula-label variable interval, maps that interval injectively to frontier coordinates, and lets the checker compute every RUP-trace support by intersecting touched live clauses with that interval; unmanifested DIMACS syntax remains outside scope. \\
Self-auditing RUP-CDCL instantiation & Carrier-independent proof-system bridge & If the trace format itself contains the checked label decoder, live-database reconstruction, RUP propagation witnesses, and retained deleted-clause support, the carrier-independent audit criterion applies with linear audit budget. \\
Append-only self-auditing RUP-CDCL & Less-annotated RUP proof bridge & If the self-auditing trace is deletion-free, the support audit is recovered from the monotone clause-addition log and RUP witnesses, so retained deleted-clause support annotations are unnecessary. \\
Prefix-typed append-only RUP-CDCL & Typed proof-system bridge & If formula-label variables are a public syntactic class of an append-only RUP-CDCL system, the audit needs neither a proof-internal label decoder nor retained deleted-clause support records; untyped raw RUP remains outside scope. \\
SAT parity application & Conditional calibration & Requires an explicit source lower bound not proved here. \\
Unrestricted standard proof-system lower bound & Not claimed & Frege, Extended Frege, IPS, compressed proofs, and unrestricted ER still require a separate extraction theorem or support budget. \\
Displayed-atom extraction theorem & Restricted bridge theorem & Applies only when accepted proofs explicitly display faithful completion-sound formula atoms. \\
Assumption-resolution theorem & Restricted named-system instantiation & Applies to resolution refutations whose measured proof object explicitly lists faithful formula-label assumptions. \\
Resolution unit-core corollary & Restricted named-system instantiation & Applies to full-assignment resolution refutations by counting the distinct faithful label-unit leaves used by the proof. \\
Resolution local-clause support theorem & Restricted named-system instantiation & Applies to resolution refutations that use explicit faithful clauses over formula-label variables, by counting their distinct label support. \\
Support-local auxiliary resolution theorem & Restricted auxiliary-resolution instantiation & Applies to mixed auxiliary encodings whose auxiliary constraints remain satisfiable under every completion preserving a fixed label support. \\
Support-local extension resolution theorem & Restricted extension-resolution instantiation & Applies to extension definitions computed from a fixed formula-label support. \\
Dependency-cone extension normalizer & Cone-annotated normalizer theorem & Applies when extension variables carry syntactic formula-label dependency cones whose union controls the final contradiction. \\
Canonical syntactic-cone normalizer & Automatic extension normalizer & Extracts dependency cones directly from the acyclic extension-definition graph; no separate cone annotations are required. \\
Final backward-slice lower bound & Proof-object budget corollary & The final extension slice must contain near-threshold many distinct formula-label inputs of each polarity. \\
Slice-explicit extension-resolution lower bound & Restricted ER-style proof-size corollary & Applies when the proof-size measure charges the distinct formula-label inputs in the final backward slice. \\
Displayed acyclic ER length lower bound & Restricted displayed-proof corollary & Applies when extension definitions and final clauses are written explicitly; ordinary displayed literal length charges the final slice boundary. \\
Expansion-bounded compressed ER transfer & Conditional compressed-proof corollary & Applies when a compressed proof expands to a displayed acyclic ER object with a controlled expansion budget. \\
Macro-compressed ER reuse tradeoff & Conditional macro-compression corollary & Applies when a displayed skeleton and macro library expand to displayed acyclic ER; forces a large skeleton-plus-reuse product. \\
Support-declared macro ER normalizer & Conditional support-budget corollary & Applies when compressed macro syntax supplies a sound declared formula-label support; forces that support budget to be near-threshold large. \\
Support-audited LRAT trace normalizer & Checkable support-budget corollary & Applies when an LRAT-style clausal trace includes a sound support audit; forces the audited support budget to be near-threshold large. \\
Provenance-typed LRAT support audit & Checkable provenance corollary & Applies when an LRAT-style trace, or DRAT trace with LRAT certificate, carries a trusted variable-role map and block-local clause provenance; the provenance package supplies the support audit. \\
Namespace-typed LRAT antecedent-closure normalizer & Checkable namespace corollary & Applies when formula-label variables occupy a trusted namespace and the LRAT antecedent closure of the empty clause exposes all used auxiliary clauses; the label variables in those used clauses form the support budget when a hard auxiliary witness exists. \\
Strict-prefix Qrep replay normalizer & Quotient-compression bridge proposition & Applies when every representative use has an earlier charged same-theory certificate forest; delayed quotient evidence is hidden support. \\
Support-audit necessity separator & Boundary proposition & Shows that a short checked clausal refutation can coexist with an unsound small support declaration; support locality must be audited separately. \\
Support-erasure indistinguishability separator & Boundary proposition & Shows that the same untyped LRAT skeleton can have an auxiliary lift with empty formula-label support and a label lift requiring all exposed labels; variable-role provenance is not cosmetic. \\
Cone-budget extension lower bound & Quantitative proof-object corollary & Forces either many final variables or a large dependency cone in cone-annotated extension proofs. \\
Bounded-depth extension tradeoff & Quantitative proof-object corollary & Forces a large final interface or logarithmic extension depth for bounded-fanin cone-annotated proofs. \\
Support-local normalization transfer theorem & Conditional proof-system transfer & Any proof system with a support-local normalizer and controlled support budget inherits the formula-footprint lower bound. \\
Carrier-independent audit upgrade criterion & Top-paper proof-system target & A system-intrinsic audit extractor, without carrier-specific namespace manifests or hidden role advice, instantiates the transfer theorem and yields a proof-size lower bound with the stated budget exponent. \\
Randomized bounded-error union footprint & Supporting theorem & Accepting-query union reduces to a deterministic certificate; expected footprint remains open. \\
\bottomrule
\end{longtable}
}

\subsection{Nonvacuity and sharpness checks}

The abstract hardness predicates are not empty.  For formula labels, take \(M_n=1\), let \(G_n\) be the identity map \(x\mapsto x\), and use the identity aggregator.  Then the formula-label carrier is the ordinary truth table of a Boolean function.  The generic sharpness theorem in \Cref{thm:identity-generic-sharpness} shows that, in this standard model, hard functions not only exist but typically have completion-sound certificates at the same scale as the lower bound up to polylogarithmic factors.  For occurrence labels, private row labels can realize any source truth table in the same way.  For quotient labels under parity, nonvacuity follows whenever the quotient incidence map \(L_{E_n}:\mathbb F_2^{\mathcal C_n}\to\mathbb F_2^{2^n}\) has an image containing a truth table outside the size-\(s\) circuit ball; full row rank is a sufficient condition.

These examples are intentionally modest.  They show that the certificate predicates are meaningful in unrestricted or rank-conditioned models.  They do not construct an explicit SAT-labeled carrier satisfying the source lower bound in the unrestricted circuit model.  Consequently, the SAT application remains conditional.

The three label granularities are genuinely different.  If a parity row contains \(M_n\) copies of one generated key \(u_x\), with \(x\mapsto u_x\) injective and \(M_n\) odd, the formula-label model needs only the shared value \(\lambda(u_x)\) for that row, whereas the private occurrence model must query all \(M_n\) positions to force the row parity.  Conversely, equality quotienting can collapse private occurrence hardness: if each parity row contains two equal occurrences and they are quotient-identified, the row contribution is always zero.  These separations justify treating the row-forcing and quotient-distance results as structural tests, while also showing that a full matching-bound theory for explicit hard SAT carriers is still future work.

\subsection{Scope and organization}

The results are deliberately scoped.  We do not prove an unrestricted circuit lower bound for an explicit SAT-labeled carrier.  We do not prove a lower bound for Frege, Extended Frege, IPS, or another standard proof system.  The occurrence theorem applies only when acceptance remains sound under independent occurrence completions; the quotient theorem applies only to equality constraints; and the formula theorem applies to functional certificates over generated formula labels.

The contribution is a compact theorem package about the nondeterministic query complexity and completion geometry of induced circuit-hardness predicates.  \Cref{sec:prelim} gives the formal model.  \Cref{sec:cheap} proves sparse-row patching and batched static lookup.  \Cref{sec:formula} proves the formula-footprint theorem.  \Cref{sec:rows} proves the row-forcing theorem.  \Cref{sec:quotient} treats equality quotients.  \Cref{sec:sat} gives the SAT magnification application and a brief naturalization observation.  \Cref{sec:audit} records the proof dependencies and nonclaim boundaries.  \Cref{sec:related} places the results in the literature and states the main open directions.

\section{Preliminaries and hardness predicates}\label{sec:prelim}

\subsection{Circuit complexity}

We use nonuniform Boolean circuits over a fixed finite fan-in-two basis.  XOR and constants have constant-size implementations and are used with constant-factor overhead.  For $f:\bits^n\to\bits$, let $\Ckt(f)$ denote its minimum circuit size.  For a size function $s$, $\SIZE(s(n))$ is the class of languages having circuits of size $O(s(n))$ at length $n$, and
\[
  \Ppoly=\bigcup_{c\ge1}\SIZE(n^c).
\]
All logarithms are base two.  Fixed constants may depend on fixed parameters such as $q,\rho,d,$ and $\eps$.

\subsection{Inputwise certificates}

Let $V$ be a finite set of Boolean variables and let $P:\bits^V\to\bits$ be a Boolean property.  A partial assignment is a map $\sigma:S\to\bits$ for some $S\subseteq V$.  A total assignment $y\in\bits^V$ extends $\sigma$ if $y|_S=\sigma$.

\begin{definition}[Inputwise $1$-certificate]\label{def:certificate}
Let $y\in\bits^V$ satisfy $P(y)=1$.  A partial assignment $\sigma:S\to\bits$ is a \emph{$1$-certificate for $P$ at $y$} if $\sigma=y|_S$ and every total assignment extending $\sigma$ also satisfies $P=1$.
The inputwise certificate complexity $C_1(P,y)$ is the minimum $|S|$ over such certificates.
\end{definition}

This is the standard nondeterministic decision-tree notion of certificate complexity; see, for example, Nisan~\cite{Nisan91} and the survey of Buhrman and de Wolf~\cite{BuhrmanDeWolf02}.  We will prove stronger, structured lower bounds that count separately which polarities or rows a certificate fixes.

\subsection{Batched generators and formula-label carriers}

\begin{definition}[Batched generator]\label{def:generator}
For each source length $n$, a batched generator is a multi-output circuit
\[
  \Gn:\bits^n\longrightarrow(\bits^{\ell_n})^{M_n},
  \qquad
  \Gn(x)=\bigl(G_{n,1}(x),\ldots,G_{n,M_n}(x)\bigr).
\]
Its size $g_n$ is the size of the single circuit producing all $M_n\ell_n$ output bits.
\end{definition}

The batched convention is important: the generator is paid once.  If $n^\rho$ were instead the cost of producing one query, the total cost could be $n^{q+\rho}$ rather than $n^{\max\{q,\rho\}}$.

Let $A_M:\bits^M\to\bits$ be a family of nonconstant aggregators having circuits of size $O(M)$.  Define the generated universe
\[
  U_n:=\{G_{n,j}(x):x\in\bits^n,\ j\in[M_n]\}.
\]
It is a set of distinct strings, so repeated generated occurrences share one formula label.

For a labeling $\lambda:U_n\to\bits$, define
\begin{equation}\label{eq:formula-carrier}
  H_{n,A}^{\lambda}(x)
  :=A_{M_n}\bigl(\lambda(G_{n,1}(x)),\ldots,
                    \lambda(G_{n,M_n}(x))\bigr).
\end{equation}

\begin{definition}[Formula-label hardness predicate]\label{def:hard-form}
For a threshold $s$, define
\[
  \HardF_{n,s,A}(\lambda)
  :=\mathbf 1[\Ckt(H_{n,A}^{\lambda})>s],
  \qquad \lambda\in\bits^{U_n}.
\]
\end{definition}

A $1$-certificate for $\HardF_{n,s,A}$ is exactly a partial formula labeling whose every completion induces a source function of circuit complexity greater than $s$.

\subsection{Occurrence-label carriers}

Let
\[
  \Omega_n:=\bits^n\times[M_n].
\]
For an occurrence assignment $z:\Omega_n\to\bits$, define
\begin{equation}\label{eq:occ-carrier}
  F_{n,A}^{z}(x)
  :=A_{M_n}\bigl(z(x,1),\ldots,z(x,M_n)\bigr).
\end{equation}

\begin{definition}[Occurrence-label hardness predicate]\label{def:hard-occ}
For a threshold $s$, define
\[
  \HardO_{n,s,A}(z)
  :=\mathbf 1[\Ckt(F_{n,A}^{z})>s],
  \qquad z\in\bits^{\Omega_n}.
\]
\end{definition}

Every formula labeling induces an occurrence assignment
\[
  z_\lambda(x,j):=\lambda(G_{n,j}(x)),
\]
but most occurrence assignments are not induced by formula labels.  The occurrence model deliberately forgets equality constraints among repeated generated strings.

\subsection{Minimum forcing numbers of an aggregator}

A partial assignment $\alpha:S\to\bits$, $S\subseteq[M]$, \emph{forces} $A_M=b$ if every total string extending $\alpha$ has aggregator value $b$.

\begin{definition}[Minimum forcing number]\label{def:forcing}
For $b\in\bits$, define
\[
  \Force_b(A_M)
  :=\min\{|S|:\text{some partial assignment on $S$ forces $A_M=b$}\}.
\]
\end{definition}

We use ``minimum forcing number'' to distinguish this quantity from the more common worst-input certificate complexity.  For the symmetric aggregators used below, the values are immediate:
\[
  \Force_0(\oplus_M)=\Force_1(\oplus_M)=M,
\]
while
\[
  \Force_0(\mathrm{OR}_M)=M,
  \qquad
  \Force_1(\mathrm{OR}_M)=1.
\]

\section{Two cheap completion lemmas}\label{sec:cheap}

The main theorems use two elementary upper bounds.  The first patches a sparse set of source rows.  The second evaluates a common hardwired dictionary on a whole batch of generated strings.

\subsection{Sparse-row patching}

\begin{lemma}[Equality patching]\label{lem:patching}
Let $f,g:\bits^n\to\bits$, and let
\[
  R:=\{x\in\bits^n:f(x)\ne g(x)\}.
\]
Then
\begin{equation}\label{eq:patching}
  \Ckt(f)\le \Ckt(g)+O((n+1)|R|).
\end{equation}
\end{lemma}

\begin{proof}
For each $r\in R$, compute the equality predicate
\[
  \Eq_r(x)=\bigwedge_{i=1}^n
  \begin{cases}
    x_i,&r_i=1,\\
    \neg x_i,&r_i=0.
  \end{cases}
\]
using $O(n)$ gates.  Let $R'\subseteq R$ be the rows on which $g$ must be flipped.  Since the predicates $\Eq_r$ are mutually exclusive,
\[
  f(x)=g(x)\oplus\bigoplus_{r\in R'}\Eq_r(x).
\]
The correction circuit has size $O((n+1)|R|)$.
\end{proof}

\begin{corollary}[Residual-error threshold]\label{cor:residual}
If $\Ckt(f)>s$ and $\Ckt(g)\le t<s$, then
\[
  \dist(f,g)=\Omega\!\left(\frac{s-t}{n+1}\right).
\]
In particular, a lower bound $\Ckt(f)>n^{1+\eps-o(1)}$ rules out every subthreshold circuit that is wrong on only $n^{\eps-o(1)}$ source rows.
\end{corollary}

\subsection{Batched static dictionary lookup}

For $M,\ell\ge1$ and $K\ge0$, define
\begin{equation}\label{eq:lookup-def}
  \Lookup(M,K,\ell)
  :=(M+K)\bigl(\ell+\log(M+K+2)\bigr)
     \log^2(M+K+2).
\end{equation}

\begin{lemma}[Batched static dictionary]\label{lem:dictionary}
Fix a dictionary $D\subseteq\bits^\ell$ of $K$ distinct hardwired keys, a hardwired bit $a(u)$ for each $u\in D$, and a default bit $b$.  There is a fan-in-two circuit that, on $M$ variable queries $q_1,\ldots,q_M\in\bits^\ell$, outputs
\[
  v_j=
  \begin{cases}
    a(q_j),&q_j\in D,\\
    b,&q_j\notin D,
  \end{cases}
  \qquad j\in[M],
\]
in size $O(\Lookup(M,K,\ell))$.
If only a symmetric aggregate of the $v_j$ is needed, the same bound holds without restoring the original query order.
\end{lemma}

\begin{proof}
Create \(M\) query records and \(K\) hardwired dictionary records.  A query record has the form
\[
  (q_i,1,i,0),
\]
where the second field is a query tag and \(i\) is the original query index.  A dictionary record for a key \(u\in D\) has the form
\[
  (u,0,0,a(u)).
\]
Sort the \(N=M+K\) records lexicographically by key and then by the tag, placing the unique dictionary record, if present, before all query records with the same key.

Batcher's sorting network has \(O(N\log^2N)\) comparators~\cite{Batcher68}.  A comparator on records of width \(O(\ell+\log N)\) has circuit size \(O(\ell+\log N)\), including comparison and conditional exchange.  After sorting, perform a left-to-right scan inside each equal-key block.  The scan records whether a dictionary record has appeared in the current block and propagates its payload value.  Each query record therefore obtains \(a(q_i)\) when \(q_i\in D\), and the default value \(b\) otherwise.  Dictionary records are masked out at the output stage.

If the final aggregator is symmetric, the circuit may aggregate the recovered values of the query records in sorted order, since only their multiset matters.  If the aggregator is not symmetric, sort all records once more by the query tag and then by the stored original index, feed the first \(M\) query records to the aggregator, and ignore the dictionary records.  The scan, masks, and optional second sort are absorbed by the same \(O(\Lookup(M,K,\ell))\) bound.  Thus original-order restoration is unnecessary only in the symmetric case; query-record separation is required in all cases.
\end{proof}

\begin{remark}[A logarithmic improvement]\label{rem:aks}
Using the $O(N\log N)$-comparator sorting networks of Ajtai, Koml\'os, and Szemer\'edi~\cite{AKS83} replaces the final $\log^2(M+K)$ factor in \eqref{eq:lookup-def} by $\log(M+K)$.  We state the transparent Batcher bound because all threshold corollaries below lose only an $n^{o(1)}$ factor.
\end{remark}

The lemma is nonuniform in exactly the intended way: the dictionary keys and values are constants in the upper-bound circuit.

For a labeling $\lambda:U_n\to\bits$ and a baseline $b\in\bits$, define the defect set
\[
  \mathrm{Def}_b(\lambda)
  :=\{u\in U_n:\lambda(u)\ne b\}.
\]

\begin{proposition}[Sparse formula-defect upper bound]\label{prop:defect}
Let $A_M$ have circuit size $O(M)$.  For every formula labeling $\lambda:U_n\to\bits$, every baseline $b$, and
\[
  K=|\mathrm{Def}_b(\lambda)|,
\]
one has
\begin{equation}\label{eq:defect-upper}
  \Ckt(H_{n,A}^{\lambda})
  \le g_n+O\bigl(\Lookup(M_n,K,\ell_n)+M_n\bigr).
\end{equation}
For a symmetric aggregator, including parity, the final reordering network is unnecessary.
\end{proposition}

\begin{proof}
Run the batched generator once.  Use the defect set as the hardwired dictionary, label every dictionary key by $1-b$, and use $b$ as the default.  The dictionary circuit recovers all $M_n$ labels, after which the size-$O(M_n)$ aggregator is applied.
\end{proof}

Repeated formulas are automatically handled by one dictionary key.  This is the quantitative reason that the formula-label barrier depends nearly linearly on $M_n+K$ rather than on their product.

\section{Formula-label certificates and near-threshold footprints}\label{sec:formula}

Fix a formula labeling $\lambda$ with $\HardF_{n,s,A}(\lambda)=1$, and let $\sigma:S\to\bits$ be a $1$-certificate at $\lambda$.  Write
\[
  S_b:=\{u\in S:\sigma(u)=b\},
  \qquad b\in\bits.
\]
For later algebra define
\begin{equation}\label{eq:lookup-envelope}
  \mathcal L_n(s)
  :=\bigl(\ell_n+\log(s+M_n+2)\bigr)
       \log^2(s+M_n+2).
\end{equation}

\begin{theorem}[Two-sided formula-footprint barrier]\label{thm:formula-footprint}
There is a constant $c\ge1$ such that every $1$-certificate $\sigma$ for $\HardF_{n,s,A}$ satisfies, for each $b\in\bits$,
\begin{equation}\label{eq:formula-implicit}
  g_n+c\Lookup(M_n,|S_b|,\ell_n)+cM_n>s.
\end{equation}
Consequently,
\begin{equation}\label{eq:formula-explicit}
  |S_b|
  \ge
  \pospart{
    \frac{s-g_n-cM_n}{c\mathcal L_n(s)}-M_n
  }.
\end{equation}
\end{theorem}

\begin{proof}
Fix $b\in\bits$.  Complete every unexposed formula label to the opposite baseline $1-b$:
\[
  \lambda^{(1-b)}(u)=
  \begin{cases}
    \sigma(u),&u\in S,\\
    1-b,&u\notin S.
  \end{cases}
\]
This total labeling extends the certificate, so
\[
  \Ckt(H_{n,A}^{\lambda^{(1-b)}})>s.
\]
Relative to the constant baseline $1-b$, its defect set is exactly $S_b$.  Applying \Cref{prop:defect} proves \eqref{eq:formula-implicit}.

For the explicit inequality, let $K=|S_b|$.  If $K>s$, then \eqref{eq:formula-explicit} is immediate after increasing $c$, since $\mathcal L_n(s)\ge1$.  If $K\le s$, then
\[
  \Lookup(M_n,K,\ell_n)
  \le(M_n+K)\mathcal L_n(s).
\]
Substitute this into \eqref{eq:formula-implicit} and rearrange.
\end{proof}

The theorem is instancewise: it applies to every hard labeling and every certificate at that labeling.  It also counts the two polarities separately.  Merely exposing many formulas is not sufficient if nearly all exposed labels have the same value.

\begin{corollary}[Nondeterministic formula-query decision-tree lower bound]\label{cor:formula-ndt}
Let \(C_1(\HardF_{n,s,A},\lambda)\) denote the standard inputwise nondeterministic decision-tree certificate complexity of the Boolean property \(\HardF_{n,s,A}\) at a hard labeling \(\lambda\).  Then
\[
  C_1(\HardF_{n,s,A},\lambda)
  \ge
  \sum_{b\in\bits}
  \pospart{
    \frac{s-g_n-cM_n}{c\mathcal L_n(s)}-M_n
  },
\]
where \(c\) is the constant from \Cref{thm:formula-footprint}.  In particular, under the hypotheses of \Cref{cor:formula-threshold},
\[
  C_1(\HardF_{n,s_n,A},\lambda)\ge 2n^{a-o(1)}.
\]
\end{corollary}

\begin{proof}
Every nondeterministic decision-tree \(1\)-certificate is exactly a partial formula-label assignment satisfying \Cref{def:certificate}.  Apply the two polarity lower bounds in \Cref{thm:formula-footprint} and sum them.
\end{proof}

\begin{corollary}[Near-threshold formula cost]\label{cor:formula-threshold}
Assume
\[
  s_n=n^{a-o(1)},\qquad
  M_n=n^{q+o(1)},\qquad
  g_n=n^{\rho+o(1)},\qquad
  \ell_n=n^{o(1)},
\]
where $a>\max\{q,\rho\}$.  Then every $1$-certificate for $\HardF_{n,s_n,A}$ satisfies
\[
  |S_0|,|S_1|\ge n^{a-o(1)}.
\]
In particular, at $a=1+\eps$ each polarity has size $n^{1+\eps-o(1)}$.
\end{corollary}

\begin{proof}
The fixed exponent gap gives $g_n+M_n=o(s_n)$, while $\mathcal L_n(s_n)=n^{o(1)}$.  Apply \Cref{thm:formula-footprint}.
\end{proof}

\begin{corollary}[Diversity and label-entropy floors]\label{cor:diversity}
Under the hypotheses of \Cref{cor:formula-threshold}, every hard labeling $\lambda$ itself has
\[
  |\lambda^{-1}(0)|,|\lambda^{-1}(1)|\ge n^{a-o(1)}.
\]
Consequently,
\[
  |U_n|\ge n^{a-o(1)}.
\]
If $U_n\subseteq\bits^{\ell_n}$, then
\[
  \ell_n\ge(a-o(1))\log n.
\]
\end{corollary}

\begin{proof}
The full assignment is always a $1$-certificate at a hard input.  Apply \Cref{cor:formula-threshold}; its polarity sets are subsets of $\lambda^{-1}(0)$ and $\lambda^{-1}(1)$.  The entropy bound follows from $|U_n|\le2^{\ell_n}$.
\end{proof}

\subsection{Generic sharpness in the identity carrier}

The formula-footprint theorem is close to tight in the standard truth-table model.  Let the identity carrier have \(M_n=1\), \(G_{n,1}(x)=x\), \(A_1(z)=z\), and \(U_n=\bits^n\).  Then a formula labeling is exactly a Boolean function \(f:\bits^n\to\bits\), and \(\HardF_{n,s,A}(f)\) is the predicate \(\Ckt(f)>s\).

\begin{theorem}[Generic sharpness in the identity carrier]\label{thm:identity-generic-sharpness}
Fix the finite circuit basis.  There is a constant \(C\) such that the following holds for all \(n,s,r\ge1\).  Let
\[
  m\ge C s\log(s+n)+r
  \qquad\text{and}\qquad
  m\le 2^n.
\]
For a uniformly random Boolean function \(f:\bits^n\to\bits\), with probability at least
\[
  1-2^{-r}-2e^{-m/18},
\]
there is a \(1\)-certificate \(\sigma:T\to\bits\) for the predicate \(\Ckt(f)>s\), at \(f\), such that
\[
  |T|=m,
  \qquad
  |\sigma^{-1}(0)|,|\sigma^{-1}(1)|\le \frac{2m}{3}.
\]
Consequently, in the identity carrier, the lower bound in \Cref{thm:formula-footprint} is sharp up to polylogarithmic factors for generic hard truth tables.
\end{theorem}

\begin{proof}
Let \(\mathcal C_{n,s}\) be the set of \(n\)-input Boolean functions computable by size-\(s\) circuits over the fixed basis.  A standard circuit count gives
\[
  |\mathcal C_{n,s}|\le 2^{C_0s\log(s+n)}
\]
for an absolute constant \(C_0\): a circuit of size \(s\) is specified by its gate types and by two input wires per gate, each chosen from at most \(n+s+O(1)\) possibilities.

Choose \(f\) uniformly at random and choose an \(m\)-element set \(T\subseteq\bits^n\) uniformly at random, independently of \(f\).  For any fixed \(h\in\mathcal C_{n,s}\),
\[
  \Pr[h|_T=f|_T]=2^{-m}.
\]
Taking \(C\ge C_0+1\) and applying the union bound gives
\[
  \Pr[\exists h\in\mathcal C_{n,s}:h|_T=f|_T]\le 2^{-r}.
\]
If no size-\(s\) function agrees with \(f\) on \(T\), then \(f|_T\) is a \(1\)-certificate for \(\Ckt(f)>s\): every completion agreeing with \(f\) on \(T\) is inconsistent with every size-\(s\) circuit.

The values \(f(x)\) on \(T\) are independent fair bits.  A Chernoff bound gives
\[
  \Pr\left[
    \max_{b\in\bits}|\{x\in T:f(x)=b\}|>\frac{2m}{3}
  \right]
  \le 2e^{-m/18}.
\]
Thus, over the joint choice of \(f\) and \(T\), the probability that \(T\) fails to be a balanced certificate is at most \(2^{-r}+2e^{-m/18}\).  If a fixed \(f\) has no such \(T\), then every choice of \(T\) fails.  Hence the same upper bound holds for the probability, over \(f\), that no balanced certificate \(T\) exists.

Taking \(m=O(s\log(s+n))\) gives certificates of size \(O(s\log(s+n))\), while \Cref{thm:formula-footprint} gives the lower bound \(\Omega(s/\polylog s)\) when the generator overhead is negligible.  This proves the claimed polylogarithmic sharpness in the identity carrier.
\end{proof}

\subsection{Adaptive deterministic verifiers}

The certificate formulation immediately covers adaptive oracle verifiers.  Let $V$ be deterministic, let $\pi$ be a proposed proof string, and let $V^\lambda(\pi)$ denote its execution with formula-label oracle $\lambda$.

\begin{proposition}[Accepted query transcripts are certificates]\label{prop:adaptive-formula}
Suppose
\[
  V^\lambda(\pi)=1
  \quad\Longrightarrow\quad
  \HardF_{n,s,A}(\lambda)=1
\]
for every formula labeling $\lambda$.  Fix an accepted execution and let $\sigma$ be its query-answer transcript.  Then $\sigma$ is a $1$-certificate for $\HardF_{n,s,A}$, so \Cref{thm:formula-footprint} applies to its $0$- and $1$-answers separately.
\end{proposition}

\begin{proof}
Every labeling extending the transcript answers all adaptive queries exactly as in the accepted execution, follows the same path, and is accepted.  Functional soundness then makes every such completion hard.
\end{proof}

Thus the formula theorem lower-bounds nondeterministic formula-query complexity without any largeness assumption on the accepted set and without requiring truth-table constructivity.

\subsection{Bounded-error accepting-union certificates}

The deterministic transcript argument has a bounded-error analogue for the union of labels seen on accepting executions.  This does not lower-bound the expected number of queries in one randomized run, but it gives a clean randomized robustness statement for the resource most directly comparable to a proof's displayed local information.

\begin{theorem}[Accepting-query union is a certificate]\label{thm:accepting-union}
Let \(P:\bits^V\to\bits\) be any Boolean property and let \(V_R\) be a randomized adaptive oracle verifier with randomness \(r\in R\).  Fix a total assignment \(y\) with \(P(y)=1\).  Let
\[
  A_y:=\{r\in R: V_R^y(r)=1\}
\]
be the accepting random strings on \(y\), and let \(U_y\subseteq V\) be the union of all variables queried during the executions \(V_R^y(r)\) with \(r\in A_y\).  Suppose there is a soundness threshold \(\beta\) such that every \(y'\) with \(P(y')=0\) satisfies
\[
  \Pr_r[V_R^{y'}(r)=1]\le \beta,
\]
while
\[
  \Pr_r[V_R^y(r)=1]>\beta.
\]
Then the partial assignment \(y|_{U_y}\) is a \(1\)-certificate for \(P\) at \(y\).
\end{theorem}

\begin{proof}
Assume toward a contradiction that \(y|_{U_y}\) is not a \(1\)-certificate.  Then some \(y'\) with \(P(y')=0\) agrees with \(y\) on every variable in \(U_y\).  For each \(r\in A_y\), the adaptive execution under \(y'\) follows the same query path as under \(y\): every query made on the accepting execution under \(y\) lies in \(U_y\) and receives the same answer under \(y'\).  Hence \(V_R^{y'}(r)=1\) for every \(r\in A_y\).  Therefore
\[
  \Pr_r[V_R^{y'}(r)=1]
  \ge
  \Pr_r[V_R^y(r)=1]
  >\beta,
\]
contradicting soundness.
\end{proof}

\begin{corollary}[Bounded-error formula union footprint]\label{cor:randomized-formula-union}
Let \(V_R\) be a randomized adaptive formula-label verifier for \(\HardF_{n,s,A}\).  Suppose it accepts a hard labeling \(\lambda\) with probability \(>\beta\), and accepts every labeling \(\lambda'\) with \(\HardF_{n,s,A}(\lambda')=0\) with probability at most \(\beta\).  Let \(U_\lambda\) be the union of generated formula labels queried on accepting executions at \(\lambda\), and write
\[
  U_{\lambda,b}:=\{u\in U_\lambda:\lambda(u)=b\}.
\]
Then, for each \(b\in\bits\),
\[
  g_n+c\Lookup(M_n,|U_{\lambda,b}|,\ell_n)+cM_n>s,
\]
with the same constant \(c\) as in \Cref{thm:formula-footprint}.  In the near-threshold regime of \Cref{cor:formula-threshold}, every bounded-error verifier has accepting union footprint
\[
  |U_{\lambda,0}|,|U_{\lambda,1}|\ge n^{a-o(1)}.
\]
\end{corollary}

\begin{proof}
Apply \Cref{thm:accepting-union} to the property \(\HardF_{n,s,A}\).  The accepting union assignment \(\lambda|_{U_\lambda}\) is a deterministic \(1\)-certificate, so \Cref{thm:formula-footprint} and \Cref{cor:formula-threshold} apply to its two polarity classes.
\end{proof}

The union resource is deliberately strong: a randomized verifier may have small expected query complexity while spreading accepting executions over a large union.  Proving expected-footprint, information, or transcript-support lower bounds requires additional distributional arguments and remains open.

\subsection{Displayed-atom proof interfaces}

The preceding verifier propositions can be stated as an extraction theorem for a restricted proof interface.  This is the formal bridge used by the paper; stronger proof-complexity consequences require normalizing a concrete proof system into this interface.

\begin{definition}[Displayed formula-atom interface]\label{def:displayed-atom-interface}
Fix \(n,s,A\).  A displayed formula-atom interface for \(\HardF_{n,s,A}\) is an acceptance relation together with an extractor that maps every accepted proof \(\pi\) to a partial assignment
\[
  \sigma_\pi:S_\pi\to\bits,
  \qquad S_\pi\subseteq U_n.
\]
For a labeling \(\lambda\), the extracted assignment is \emph{faithful at \(\lambda\)} if \(\sigma_\pi(u)=\lambda(u)\) for every \(u\in S_\pi\).  It is \emph{completion-sound} if every formula labeling extending \(\sigma_\pi\) satisfies \(\HardF_{n,s,A}=1\).  The interface is \emph{literal-listing} if each extracted atom is displayed in the proof by writing its generated key \(u\in\bits^{\ell_n}\) and its bit value.
\end{definition}

\begin{theorem}[Displayed-atom extraction]\label{thm:displayed-atom-extraction}
Let \(\lambda\) satisfy \(\HardF_{n,s,A}(\lambda)=1\), and let \(\pi\) be an accepted proof in a displayed formula-atom interface whose extracted assignment \(\sigma_\pi:S_\pi\to\bits\) is faithful at \(\lambda\) and completion-sound.  Write
\[
  S_{\pi,b}:=\{u\in S_\pi:\sigma_\pi(u)=b\}.
\]
Then, for each \(b\in\bits\),
\[
  g_n+c\Lookup(M_n,|S_{\pi,b}|,\ell_n)+cM_n>s,
\]
with the same constant \(c\) as in \Cref{thm:formula-footprint}.  In the near-threshold regime of \Cref{cor:formula-threshold},
\[
  |S_{\pi,0}|,|S_{\pi,1}|\ge n^{a-o(1)}.
\]
If the interface is literal-listing, every such proof displays at least
\[
  (|S_{\pi,0}|+|S_{\pi,1}|)(\ell_n+1)
\]
atom bits before any additional proof overhead.
\end{theorem}

\begin{proof}
Faithfulness says that \(\sigma_\pi\) is a partial assignment at the actual hard labeling.  Completion soundness says that every extension of \(\sigma_\pi\) remains in the hardness predicate.  Hence \(\sigma_\pi\) is a \(1\)-certificate for \(\HardF_{n,s,A}\) at \(\lambda\).  The footprint inequalities follow from \Cref{thm:formula-footprint}, and the near-threshold lower bound follows from \Cref{cor:formula-threshold}.  In a literal-listing interface, each extracted atom explicitly contains an \(\ell_n\)-bit key and a one-bit value, giving the displayed atom-bit lower bound by counting the listed atoms.
\end{proof}

\begin{remark}[What this extraction theorem does and does not say]\label{rem:displayed-atom-scope}
\Cref{thm:displayed-atom-extraction} is intentionally restricted.  It covers proof interfaces whose local semantic content is an explicit faithful list of formula labels and whose soundness is stable under all completions of unlisted labels.  It does not by itself lower-bound Frege, Extended Frege, IPS, implicit proofs, or algebraic proof systems, which may encode semantic information nonlocally.  Its role is to make the required normalization target precise: to use the formula-footprint theorem in a standard proof system, one must extract a displayed completion-sound atom set with controlled overhead.
\end{remark}

\subsection{A restricted resolution instantiation}

The displayed-atom theorem gives a concrete proof-system consequence for a standard but deliberately restricted resolution format: resolution refutations with explicit formula-label assumptions.  This is not a lower bound for unrestricted resolution encodings of circuit lower bounds.  It is a normalization theorem for the case where the proof derives contradiction from an explicitly listed set of local label facts.

\begin{definition}[Easy-completion CNF]\label{def:easy-completion-cnf}
An \emph{easy-completion CNF} for \(\HardF_{n,s,A}\) is a CNF
\[
  \Phi_{n,s,A}(Y,Z)
\]
over formula-label variables \(Y=\{Y_u:u\in U_n\}\) and auxiliary variables \(Z\) such that, for every total formula labeling \(\mu:U_n\to\bits\),
\[
  \exists z\ \Phi_{n,s,A}(\mu,z)=1
  \quad\Longleftrightarrow\quad
  \EasyF_{n,s,A}(\mu),
\]
where \(\EasyF_{n,s,A}\) is the complement of \(\HardF_{n,s,A}\).
\end{definition}

\begin{theorem}[Resolution with explicit label assumptions]\label{thm:assumption-resolution}
Let \(\lambda\) satisfy \(\HardF_{n,s,A}(\lambda)=1\).  Let \(\Phi_{n,s,A}\) be an easy-completion CNF, and suppose there is a resolution refutation of
\[
  \Phi_{n,s,A}(Y,Z)\wedge \Gamma
\]
where \(\Gamma\) is a set of unit clauses over formula-label variables of the form \(Y_u=\lambda(u)\).  Let
\[
  S_\Gamma:=\{u\in U_n: Y_u=\lambda(u)\text{ appears in }\Gamma\},
  \qquad
  S_{\Gamma,b}:=\{u\in S_\Gamma:\lambda(u)=b\}.
\]
Then, for each \(b\in\bits\),
\[
  g_n+c\Lookup(M_n,|S_{\Gamma,b}|,\ell_n)+cM_n>s,
\]
with the same constant \(c\) as in \Cref{thm:formula-footprint}.  In the near-threshold regime of \Cref{cor:formula-threshold},
\[
  |S_{\Gamma,0}|,|S_{\Gamma,1}|\ge n^{a-o(1)}.
\]
Consequently, any resolution proof format that counts these label assumptions as displayed unit clauses has at least
\[
  |S_{\Gamma,0}|+|S_{\Gamma,1}|
\]
displayed label assumptions, and at least
\[
  (|S_{\Gamma,0}|+|S_{\Gamma,1}|)(\ell_n+1)
\]
displayed label-atom bits, before counting the rest of the resolution derivation.
\end{theorem}

\begin{proof}
Resolution soundness implies that \(\Phi_{n,s,A}\wedge\Gamma\) has no satisfying assignment.  By the defining completeness of the easy-completion CNF, there is therefore no total formula labeling \(\mu\) extending the unit assumptions in \(\Gamma\) with \(\EasyF_{n,s,A}(\mu)\).  Equivalently, every formula labeling extending \(\Gamma\) lies in \(\HardF_{n,s,A}\).  Since the assumptions are faithful to the actual hard labeling \(\lambda\), the partial assignment specified by \(\Gamma\) is a \(1\)-certificate for \(\HardF_{n,s,A}\) at \(\lambda\).  The two polarity inequalities and the near-threshold bounds are exactly \Cref{thm:formula-footprint,cor:formula-threshold}.  If the proof format explicitly lists label assumptions as unit clauses, each distinct assumption displays one generated key and one bit, giving the stated displayed-assumption and displayed-bit lower bounds.
\end{proof}

\begin{remark}[Scope of the resolution instantiation]\label{rem:assumption-resolution-scope}
\Cref{thm:assumption-resolution} is a genuine resolution statement, but only for assumption-based encodings where local label facts are part of the measured proof object.  It does not rule out short resolution, Frege, Extended Frege, IPS, or implicit proofs that encode the needed semantic information without explicitly listing faithful label assumptions.  Its value is as a first controlled normalization theorem: the proof-system theorem becomes strong exactly when an external argument shows that a natural proof format must expose such assumptions.
\end{remark}

\begin{corollary}[Resolution unit-core lower bound]\label{cor:resolution-unit-core}
Let \(\lambda\) satisfy \(\HardF_{n,s,A}(\lambda)=1\), and let \(\Lambda_\lambda\) be the full set of unit clauses \(Y_u=\lambda(u)\) for \(u\in U_n\).  For any resolution refutation \(\Pi\) of
\[
  \Phi_{n,s,A}(Y,Z)\wedge\Lambda_\lambda,
\]
let \(T_\Pi\subseteq U_n\) be the set of formula-label variables whose unit clauses from \(\Lambda_\lambda\) occur as initial clauses used in \(\Pi\), and set
\[
  T_{\Pi,b}:=\{u\in T_\Pi:\lambda(u)=b\}.
\]
Then, for each \(b\in\bits\),
\[
  g_n+c\Lookup(M_n,|T_{\Pi,b}|,\ell_n)+cM_n>s.
\]
In the near-threshold regime of \Cref{cor:formula-threshold},
\[
  |T_{\Pi,0}|,|T_{\Pi,1}|\ge n^{a-o(1)}.
\]
Thus the lower bound applies to the \emph{used unit core} of any full-assignment resolution refutation, not only to a preselected assumption set.
\end{corollary}

\begin{proof}
Discard every initial unit clause from \(\Lambda_\lambda\) that is not used as a leaf of the refutation \(\Pi\).  The remaining proof is still a resolution refutation of \(\Phi_{n,s,A}\wedge\Gamma_\Pi\), where
\[
  \Gamma_\Pi:=\{Y_u=\lambda(u):u\in T_\Pi\}.
\]
Apply \Cref{thm:assumption-resolution} to \(\Gamma_\Pi\).
\end{proof}

\begin{theorem}[Resolution local-clause support]\label{thm:resolution-local-clause-support}
Let \(\lambda\) satisfy \(\HardF_{n,s,A}(\lambda)=1\), and let \(\Phi_{n,s,A}\) be an easy-completion CNF.  Let \(\Gamma\) be any CNF over formula-label variables \(Y=\{Y_u:u\in U_n\}\), with no auxiliary variables, such that \(\lambda\models\Gamma\).  Suppose there is a resolution refutation of
\[
  \Phi_{n,s,A}(Y,Z)\wedge\Gamma.
\]
Let \(S_\Gamma\subseteq U_n\) be the set of formula labels whose variables occur in \(\Gamma\), and set
\[
  S_{\Gamma,b}:=\{u\in S_\Gamma:\lambda(u)=b\}.
\]
Then, for each \(b\in\bits\),
\[
  g_n+c\Lookup(M_n,|S_{\Gamma,b}|,\ell_n)+cM_n>s.
\]
In the near-threshold regime of \Cref{cor:formula-threshold},
\[
  |S_{\Gamma,0}|,|S_{\Gamma,1}|\ge n^{a-o(1)}.
\]
Thus the lower bound applies to the two-polarity label support of any explicitly displayed local label clauses, even when those clauses are not unit clauses.
\end{theorem}

\begin{proof}
Resolution soundness implies that \(\Phi_{n,s,A}\wedge\Gamma\) is unsatisfiable.  Let \(\mu:U_n\to\bits\) be any formula labeling extending \(\lambda|_{S_\Gamma}\).  Since \(\Gamma\) contains only variables from \(S_\Gamma\) and \(\lambda\models\Gamma\), the labeling \(\mu\) also satisfies \(\Gamma\).  If \(\mu\) were easy, then by \Cref{def:easy-completion-cnf} there would be an auxiliary assignment \(z\) satisfying \(\Phi_{n,s,A}(\mu,z)\), and hence \((\mu,z)\) would satisfy \(\Phi_{n,s,A}\wedge\Gamma\), a contradiction.  Therefore every extension of \(\lambda|_{S_\Gamma}\) lies in \(\HardF_{n,s,A}\).  The partial assignment \(\lambda|_{S_\Gamma}\) is a \(1\)-certificate at \(\lambda\), and \Cref{thm:formula-footprint,cor:formula-threshold} give the claimed polarity bounds.
\end{proof}

\begin{remark}[Local clauses versus nonlocal encodings]\label{rem:local-clause-scope}
\Cref{thm:resolution-local-clause-support} counts the explicit label support of clauses over formula-label variables.  It does not apply to encodings that mix label variables with auxiliary variables in a way that lets the proof hide semantic information outside a fixed label support.  Such encodings are precisely where a stronger normalization theorem would be needed.
\end{remark}

\begin{theorem}[Support-local auxiliary resolution]\label{thm:support-local-aux-resolution}
Let \(\lambda\) satisfy \(\HardF_{n,s,A}(\lambda)=1\), and let \(\Phi_{n,s,A}\) be an easy-completion CNF.  Let \(\Gamma(Y,W)\) be a CNF over formula-label variables \(Y=\{Y_u:u\in U_n\}\) and auxiliary variables \(W\).  Suppose there is a set \(S\subseteq U_n\) such that \(\Gamma\) is \emph{support-local at \(\lambda\) over \(S\)}: for every formula labeling \(\mu:U_n\to\bits\) with \(\mu|_S=\lambda|_S\), there exists an assignment \(w\) to \(W\) satisfying \(\Gamma(\mu,w)\).  If there is a resolution refutation of
\[
  \Phi_{n,s,A}(Y,Z)\wedge\Gamma(Y,W),
\]
then, writing
\[
  S_b:=\{u\in S:\lambda(u)=b\},
\]
we have, for each \(b\in\bits\),
\[
  g_n+c\Lookup(M_n,|S_b|,\ell_n)+cM_n>s.
\]
In the near-threshold regime of \Cref{cor:formula-threshold},
\[
  |S_0|,|S_1|\ge n^{a-o(1)}.
\]
\end{theorem}

\begin{proof}
Resolution soundness implies that \(\Phi_{n,s,A}\wedge\Gamma\) is unsatisfiable.  Let \(\mu\) be any formula labeling extending \(\lambda|_S\).  By support-locality, choose \(w\) satisfying \(\Gamma(\mu,w)\).  If \(\mu\) were easy, then by \Cref{def:easy-completion-cnf} there would be an auxiliary assignment \(z\) satisfying \(\Phi_{n,s,A}(\mu,z)\).  Then \((\mu,z,w)\) would satisfy \(\Phi_{n,s,A}\wedge\Gamma\), contradicting the resolution refutation.  Thus every completion of \(\lambda|_S\) lies in \(\HardF_{n,s,A}\), so \(\lambda|_S\) is a \(1\)-certificate at \(\lambda\).  The bounds follow from \Cref{thm:formula-footprint,cor:formula-threshold}.
\end{proof}

\begin{remark}[Where auxiliary encodings still escape]\label{rem:support-local-aux-scope}
\Cref{thm:support-local-aux-resolution} allows auxiliary variables and mixed clauses, but only under a semantic locality invariant: once the labels in \(S\) are fixed, the auxiliary constraints remain satisfiable for every completion of the other labels.  A proof system that uses auxiliary variables to encode global information about the unexposed labels need not satisfy this invariant.  Proving that such a system admits a support-local normalization is exactly the stronger proof-system problem left open here.
\end{remark}

\begin{theorem}[Support-local extension resolution]\label{thm:support-local-extension-resolution}
Let \(\lambda\) satisfy \(\HardF_{n,s,A}(\lambda)=1\), and let \(\Phi_{n,s,A}\) be an easy-completion CNF.  Fix \(S\subseteq U_n\).  Let \(C\) be an acyclic Boolean extension circuit whose input variables are the formula-label variables \(Y_S=\{Y_u:u\in S\}\), and let \(W\) be its extension gates.  Let \(E_C(Y_S,W)\) be the standard CNF encoding of the extension definitions, so every assignment to \(Y_S\) has a satisfying extension \(w=C(Y_S)\).  Let \(\Theta(W)\) be a CNF over extension variables only, and suppose \(\Theta(C(\lambda|_S))=1\).  If there is a resolution refutation of
\[
  \Phi_{n,s,A}(Y,Z)\wedge E_C(Y_S,W)\wedge \Theta(W),
\]
then, writing
\[
  S_b:=\{u\in S:\lambda(u)=b\},
\]
we have, for each \(b\in\bits\),
\[
  g_n+c\Lookup(M_n,|S_b|,\ell_n)+cM_n>s.
\]
In the near-threshold regime of \Cref{cor:formula-threshold},
\[
  |S_0|,|S_1|\ge n^{a-o(1)}.
\]
\end{theorem}

\begin{proof}
Let \(\mu:U_n\to\bits\) be any formula labeling with \(\mu|_S=\lambda|_S\).  Because \(C\) reads only \(Y_S\), the extension value \(w=C(\mu|_S)\) equals \(C(\lambda|_S)\), and hence satisfies both \(E_C(\mu|_S,w)\) and \(\Theta(w)\).  Therefore \(E_C\wedge\Theta\) is support-local at \(\lambda\) over \(S\).  Applying \Cref{thm:support-local-aux-resolution} to \(\Gamma(Y,W)=E_C(Y_S,W)\wedge\Theta(W)\) shows that \(\lambda|_S\) is a completion-sound formula-label certificate.  The claimed inequalities are exactly the formula-footprint and near-threshold bounds.
\end{proof}

\begin{remark}[Scope of the extension-resolution statement]\label{rem:support-local-er-scope}
\Cref{thm:support-local-extension-resolution} is an extended-resolution-style instantiation of the support-local bridge, not a lower bound for unrestricted Extended Frege, unrestricted extended resolution, or algebraic proof systems.  Its extension definitions may compute derived predicates, but only from the fixed formula-label support \(S\), and its final extension clauses must be stable once \(\lambda|_S\) is fixed.  A proof system whose extension variables depend on nonlocal semantic information still requires a separate normalizer.
\end{remark}

\begin{definition}[Dependency-cone extension package]\label{def:dependency-cone-extension}
Let \(C\) be an acyclic Boolean extension circuit with formula-label inputs \(Y=\{Y_u:u\in U_n\}\) and extension variables \(W\).  A \emph{dependency-cone annotation} assigns to every label variable \(Y_u\) the cone \(\Delta(Y_u)=\{u\}\), and to every extension variable \(w\in W\) a set \(\Delta(w)\subseteq U_n\), such that if \(w\) is defined from earlier variables \(v_1,\ldots,v_r\), then
\[
  \Delta(w)\supseteq\Delta(v_1)\cup\cdots\cup\Delta(v_r).
\]
For a CNF \(\Theta\) over label and extension variables, its \emph{declared final support} is
\[
  S_\Theta:=\bigcup_{X\text{ occurs in }\Theta}\Delta(X).
\]
\end{definition}

\begin{theorem}[Dependency-cone extension normalizer]\label{thm:dependency-cone-extension-normalizer}
Let \(\lambda\) satisfy \(\HardF_{n,s,A}(\lambda)=1\), and let \(\Phi_{n,s,A}\) be an easy-completion CNF.  Let \(C\) be an acyclic Boolean extension circuit with a dependency-cone annotation \(\Delta\), let \(E_C(Y,W)\) be the standard CNF encoding of its extension definitions, and let \(\Theta(Y,W)\) be a CNF over label and extension variables.  Suppose the extension assignment obtained by evaluating \(C\) on \(\lambda\) satisfies \(\Theta\).  If there is a resolution refutation of
\[
  \Phi_{n,s,A}(Y,Z)\wedge E_C(Y,W)\wedge\Theta(Y,W),
\]
then, writing
\[
  S_{\Theta,b}:=\{u\in S_\Theta:\lambda(u)=b\},
\]
we have, for each \(b\in\bits\),
\[
  g_n+c\Lookup(M_n,|S_{\Theta,b}|,\ell_n)+cM_n>s.
\]
In the near-threshold regime of \Cref{cor:formula-threshold},
\[
  |S_{\Theta,0}|,|S_{\Theta,1}|\ge n^{a-o(1)}.
\]
In particular, any proof system whose accepted extension proof \(\pi\) supplies such a cone annotation with \(|S_\Theta|\le B(|\pi|)\) has a support-local normalization with budget \(B\).
\end{theorem}

\begin{proof}
Let \(\mu:U_n\to\bits\) be any formula labeling with \(\mu|_{S_\Theta}=\lambda|_{S_\Theta}\).  We first prove by induction over the acyclic extension order that every variable \(X\) with \(\Delta(X)\subseteq S_\Theta\) has the same value under evaluation on \(\mu\) as under evaluation on \(\lambda\).  This is immediate for label variables because their cones are singletons.  For an extension variable \(w\), every input variable to the definition of \(w\) has its cone contained in \(\Delta(w)\).  If \(\Delta(w)\subseteq S_\Theta\), the induction hypothesis gives equality of all input values, so the value of \(w\) is also the same under \(\mu\) and \(\lambda\).

Thus the evaluation assignment \(w=C(\mu)\) satisfies \(\Theta(\mu,w)\), because \(C(\lambda)\) satisfies \(\Theta(\lambda,C(\lambda))\).  The extension clauses \(E_C(\mu,w)\) are satisfied by construction.  Hence \(E_C\wedge\Theta\) is support-local at \(\lambda\) over \(S_\Theta\).  Applying \Cref{thm:support-local-aux-resolution} to \(\Gamma(Y,W)=E_C(Y,W)\wedge\Theta(Y,W)\) yields that \(\lambda|_{S_\Theta}\) is a completion-sound formula-label certificate, and \Cref{thm:formula-footprint,cor:formula-threshold} give the polarity lower bounds.  The final sentence is exactly \Cref{def:support-local-normalizer} with the extracted support \(S_\Theta\).
\end{proof}

\begin{remark}[What cone annotations buy]\label{rem:dependency-cone-scope}
\Cref{thm:dependency-cone-extension-normalizer} is still not an unrestricted extended-resolution lower bound.  It says that a proof artifact becomes usable by the formula-footprint theorem once every extension variable exposes a syntactic upper bound on the formula labels it can semantically depend on.  This is a normalizer one step closer to ordinary proof objects: the support is not chosen in advance, but extracted from the declared dependency cones of the final contradiction.
\end{remark}

\begin{corollary}[Canonical syntactic-cone normalizer]\label{cor:canonical-syntactic-cone}
Let \(\lambda\), \(\Phi_{n,s,A}\), \(C\), \(E_C\), and \(\Theta\) be as in \Cref{thm:dependency-cone-extension-normalizer}, but assume no external cone annotation is supplied.  Define the canonical syntactic cone \(\operatorname{Syn}_C(X)\subseteq U_n\) by
\[
  \operatorname{Syn}_C(Y_u)=\{u\},\qquad
  \operatorname{Syn}_C(w)=\operatorname{Syn}_C(v_1)\cup\cdots\cup\operatorname{Syn}_C(v_r)
\]
whenever the extension variable \(w\) is defined from earlier variables \(v_1,\ldots,v_r\).  Let
\[
  S^{\mathrm{syn}}_\Theta:=\bigcup_{X\text{ occurs in }\Theta}\operatorname{Syn}_C(X).
\]
If the extension assignment obtained by evaluating \(C\) on \(\lambda\) satisfies \(\Theta\), and there is a resolution refutation of
\[
  \Phi_{n,s,A}(Y,Z)\wedge E_C(Y,W)\wedge\Theta(Y,W),
\]
then, writing
\[
  S^{\mathrm{syn}}_{\Theta,b}:=\{u\in S^{\mathrm{syn}}_\Theta:\lambda(u)=b\},
\]
we have, for each \(b\in\bits\),
\[
  g_n+c\Lookup(M_n,|S^{\mathrm{syn}}_{\Theta,b}|,\ell_n)+cM_n>s.
\]
In the near-threshold regime of \Cref{cor:formula-threshold},
\[
  |S^{\mathrm{syn}}_{\Theta,0}|,|S^{\mathrm{syn}}_{\Theta,1}|\ge n^{a-o(1)}.
\]
In particular, any explicit acyclic extension proof whose final syntactic-cone support satisfies \(|S^{\mathrm{syn}}_\Theta|\le B(|\pi|)\) has a support-local normalization with budget \(B\).
\end{corollary}

\begin{proof}
The recursively defined sets \(\operatorname{Syn}_C(X)\) form a valid dependency-cone annotation: label variables have singleton cones, and every extension variable's cone contains the cones of its defining inputs with equality.  Applying \Cref{thm:dependency-cone-extension-normalizer} to this canonical annotation gives the stated support \(S^{\mathrm{syn}}_\Theta\), the two polarity inequalities, and the near-threshold bounds.  The final sentence is the support-local normalization clause of \Cref{thm:dependency-cone-extension-normalizer} with the canonical support.
\end{proof}

\begin{remark}[Scope of the canonical cone normalizer]\label{rem:canonical-syntactic-cone-scope}
\Cref{cor:canonical-syntactic-cone} removes the need for separate cone annotations: the cone is read from the extension-definition graph itself.  It still does not prove a lower bound for unrestricted extended resolution, Extended Frege, Frege, IPS, or implicit proofs.  If the final contradiction mentions variables with huge syntactic cones, or if a stronger proof system compresses nonlocal semantic information in a way not represented by acyclic label-dependency syntax, the bound may be vacuous.  The theorem's payoff is therefore precise: ordinary acyclic extension syntax is enough for normalization, but a selective proof-system consequence still needs a nontrivial syntactic-cone support bound.
\end{remark}

\begin{corollary}[Final backward-slice lower bound]\label{cor:final-backward-slice}
In the setting of \Cref{cor:canonical-syntactic-cone}, let \(\operatorname{Slice}_C(\Theta)\) be the backward slice of the acyclic extension-definition graph containing every label or extension variable that reaches a variable occurring in \(\Theta\).  Let
\[
  \partial_Y\operatorname{Slice}_C(\Theta):=\{u\in U_n:Y_u\text{ occurs as a label input in }\operatorname{Slice}_C(\Theta)\}.
\]
If the extension assignment induced by \(\lambda\) satisfies \(\Theta\), and
\[
  \Phi_{n,s,A}(Y,Z)\wedge E_C(Y,W)\wedge\Theta(Y,W)
\]
has a resolution refutation, then, for each \(b\in\bits\),
\[
  g_n+c\Lookup(M_n,|\partial_Y\operatorname{Slice}_C(\Theta)\cap\lambda^{-1}(b)|,\ell_n)+cM_n>s.
\]
In the near-threshold regime of \Cref{cor:formula-threshold},
\[
  |\partial_Y\operatorname{Slice}_C(\Theta)\cap\lambda^{-1}(0)|,\,
  |\partial_Y\operatorname{Slice}_C(\Theta)\cap\lambda^{-1}(1)|
  \ge n^{a-o(1)}.
\]
Consequently, any explicit acyclic extension proof representation whose length counts the distinct formula-label inputs appearing in the final backward slice has length at least \(2n^{a-o(1)}\) in this regime.
\end{corollary}

\begin{proof}
By construction, \(\partial_Y\operatorname{Slice}_C(\Theta)=S^{\mathrm{syn}}_\Theta\): a label \(u\) lies in the final backward slice exactly when there is a directed path in the extension-definition graph from \(Y_u\) to some variable occurring in \(\Theta\), which is exactly the recursive definition of the canonical syntactic cone union.  The two polarity inequalities and near-threshold lower bounds are therefore \Cref{cor:canonical-syntactic-cone}.  The final statement follows because the proof representation is assumed to count each distinct formula-label input in the backward slice at least once, and the two polarity sets are disjoint.
\end{proof}

\begin{remark}[What the slice lower bound does and does not say]\label{rem:final-backward-slice-scope}
\Cref{cor:final-backward-slice} is the proof-object budget form of the canonical syntactic-cone normalizer.  It applies to explicit acyclic extension syntax and says that the final contradiction cannot be supported by a small local backward slice.  It is still not an unrestricted extended-resolution, Extended Frege, Frege, or IPS lower bound: a stronger proof system can escape the statement if the relevant semantic information is compressed outside such an explicit acyclic label-dependency slice, or if the final backward slice is itself large.
\end{remark}

\begin{definition}[Slice-explicit extension refutation]\label{def:slice-explicit-er}
A \emph{slice-explicit extension refutation} of the easy-completion CNF consists of:
\begin{enumerate}
  \item an explicit acyclic extension circuit \(C\) over formula-label inputs \(Y=\{Y_u:u\in U_n\}\);
  \item its CNF definition \(E_C(Y,W)\);
  \item a final CNF \(\Theta(Y,W)\);
  \item a resolution refutation of \(\Phi_{n,s,A}(Y,Z)\wedge E_C(Y,W)\wedge\Theta(Y,W)\).
\end{enumerate}
Its slice size \(\operatorname{size}_{\mathrm{slice}}(\pi)\) is any proof-size measure that is at least the number of distinct formula-label inputs in \(\partial_Y\operatorname{Slice}_C(\Theta)\).
\end{definition}

\begin{corollary}[Slice-explicit extension-resolution lower bound]\label{cor:slice-explicit-er}
Let \(\lambda\) satisfy \(\HardF_{n,s,A}(\lambda)=1\), and let \(\pi\) be a slice-explicit extension refutation whose final CNF \(\Theta\) is satisfied by evaluating its extension circuit \(C\) on \(\lambda\).  Then, for each \(b\in\bits\),
\[
  g_n+c\Lookup(M_n,|\partial_Y\operatorname{Slice}_C(\Theta)\cap\lambda^{-1}(b)|,\ell_n)+cM_n>s.
\]
In the near-threshold regime of \Cref{cor:formula-threshold},
\[
  \operatorname{size}_{\mathrm{slice}}(\pi)\ge 2n^{a-o(1)}.
\]
\end{corollary}

\begin{proof}
\Cref{cor:final-backward-slice} gives the two polarity lower bounds for \(\partial_Y\operatorname{Slice}_C(\Theta)\).  In the near-threshold regime, both polarity subsets have size at least \(n^{a-o(1)}\), and they are disjoint.  By the definition of slice-explicit size, \(\operatorname{size}_{\mathrm{slice}}(\pi)\) is at least the number of distinct formula-label inputs in their union, so it is at least \(2n^{a-o(1)}\).
\end{proof}

\begin{remark}[Restricted ER-style payoff]\label{rem:slice-explicit-er-scope}
\Cref{cor:slice-explicit-er} is a restricted extended-resolution-style lower bound for proof representations that explicitly charge the formula-label boundary of the final extension slice.  It is more recognizable than an abstract support-local normalizer, because the measured object is read from an ordinary acyclic extension-definition graph.  It is still not an unrestricted extended-resolution, Extended Frege, Frege, or IPS lower bound: the theorem assumes the final slice is explicitly represented and charged.
\end{remark}

\begin{corollary}[Displayed acyclic ER length lower bound]\label{cor:displayed-er-length}
Let \(\lambda\) satisfy \(\HardF_{n,s,A}(\lambda)=1\), and let \(\pi\) be a slice-explicit extension refutation as in \Cref{def:slice-explicit-er}.  Assume that every extension definition is displayed by CNF clauses in \(E_C\) with its operand variables written explicitly, and define the displayed extension-object length \(\operatorname{len}_{\mathrm{disp}}(\pi)\) to be the number of literal occurrences in \(E_C\) and \(\Theta\), plus the literal occurrences in the resolution derivation.  If the final CNF \(\Theta\) is satisfied by evaluating \(C\) on \(\lambda\), then, in the near-threshold regime of \Cref{cor:formula-threshold},
\[
  \operatorname{len}_{\mathrm{disp}}(\pi)\ge 2n^{a-o(1)}.
\]
\end{corollary}

\begin{proof}
\Cref{cor:final-backward-slice} gives
\[
  |\partial_Y\operatorname{Slice}_C(\Theta)\cap\lambda^{-1}(0)|,\,
  |\partial_Y\operatorname{Slice}_C(\Theta)\cap\lambda^{-1}(1)|
  \ge n^{a-o(1)}.
\]
It remains to compare this boundary with ordinary displayed length.  If \(u\in\partial_Y\operatorname{Slice}_C(\Theta)\), then \(Y_u\) lies on a directed path to a variable occurring in \(\Theta\).  If the path has length \(0\), then \(Y_u\) itself occurs in \(\Theta\).  Otherwise, let \(X\) be the first extension variable on such a path.  The definition of \(X\) has \(Y_u\) as a displayed operand, so the CNF clauses defining \(X\) in \(E_C\) contain a literal over \(Y_u\).  Thus every distinct formula-label input in the final backward slice contributes at least one displayed literal occurrence in \(E_C\cup\Theta\).  The two polarity subsets are disjoint, so the displayed extension-object length is at least their sum, which is at least \(2n^{a-o(1)}\).
\end{proof}

\begin{remark}[Displayed length versus unrestricted ER]\label{rem:displayed-er-length-scope}
\Cref{cor:displayed-er-length} removes the separate slice-size charging convention for explicitly written acyclic extension objects: ordinary displayed literal length already pays for every formula-label input in the final backward slice.  The boundary is now the standard one for compressed proof systems.  If a proof format hides extension definitions behind macros, encodes label identities implicitly, or supplies nonlocal semantic information without a displayed acyclic dependency graph, this corollary does not apply.  In particular, it is not a lower bound for Frege, Extended Frege, IPS, implicit proofs, or unrestricted ER encodings whose effective dependency support is not visible in the displayed clauses.
\end{remark}

\begin{corollary}[Expansion-bounded compressed ER transfer]\label{cor:compressed-er-expansion}
Let \(\lambda\) satisfy \(\HardF_{n,s,A}(\lambda)=1\).  Let \(\rho\) be a compressed acyclic extension refutation that expands, by a sound decompression map \(\operatorname{Exp}\), to a displayed acyclic extension refutation \(\pi=\operatorname{Exp}(\rho)\) satisfying the hypotheses of \Cref{cor:displayed-er-length}.  Suppose the displayed expansion length is bounded by a monotone function \(B\) of compressed length:
\[
  \operatorname{len}_{\mathrm{disp}}(\operatorname{Exp}(\rho))
  \le B(|\rho|).
\]
Then, in the near-threshold regime of \Cref{cor:formula-threshold},
\[
  B(|\rho|)\ge 2n^{a-o(1)}.
\]
In particular, if \(B(L)\le C L^k\) for constants \(C,k\), then every such compressed refutation has
\[
  |\rho|\ge C^{-1/k}2^{1/k}n^{a/k-o(1)}.
\]
\end{corollary}

\begin{proof}
The expanded object \(\pi=\operatorname{Exp}(\rho)\) is a displayed acyclic extension refutation satisfying the hypotheses of \Cref{cor:displayed-er-length}.  Therefore
\[
  \operatorname{len}_{\mathrm{disp}}(\operatorname{Exp}(\rho))
  \ge 2n^{a-o(1)}.
\]
Combining this with the assumed expansion upper bound gives \(B(|\rho|)\ge 2n^{a-o(1)}\).  If \(B(L)\le C L^k\), rearranging \(C|\rho|^k\ge 2n^{a-o(1)}\) yields the stated compressed-length lower bound.
\end{proof}

\begin{remark}[What compression must prove]\label{rem:compressed-er-expansion-scope}
\Cref{cor:compressed-er-expansion} is not a lower bound for arbitrary compressed proofs.  It says that compression helps only by making the displayed expansion large, by hiding the relevant dependency support, or by leaving the displayed acyclic ER model.  Thus a future proof-system theorem must either prove a controlled expansion/support budget for the compressed representation, or explain why its proof object is genuinely outside the local completion models used here.
\end{remark}

\begin{corollary}[Macro-compressed ER reuse tradeoff]\label{cor:macro-compressed-er}
Let \(\lambda\) satisfy \(\HardF_{n,s,A}(\lambda)=1\).  Suppose a macro-compressed acyclic extension refutation \(\rho\) consists of:
\begin{enumerate}
  \item a displayed skeleton of literal length \(S\);
  \item a displayed macro library of total literal length \(D\);
  \item \(r\) macro invocations, each expanding to a library body;
  \item a sound expansion \(\operatorname{Exp}(\rho)\) to a displayed acyclic extension refutation satisfying the hypotheses of \Cref{cor:displayed-er-length}.
\end{enumerate}
Then, in the near-threshold regime of \Cref{cor:formula-threshold},
\[
  S+rD\ge 2n^{a-o(1)}.
\]
Consequently, if the encoded length \(L\) of \(\rho\) satisfies \(S,D,r\le L\), then
\[
  L^2+L\ge 2n^{a-o(1)},
  \qquad\text{and hence}\qquad
  L\ge n^{a/2-o(1)}.
\]
\end{corollary}

\begin{proof}
The displayed expansion consists of the skeleton plus, for each of the \(r\) invocations, one displayed macro body from a library whose total body length is \(D\).  Hence
\[
  \operatorname{len}_{\mathrm{disp}}(\operatorname{Exp}(\rho))\le S+rD.
\]
By \Cref{cor:displayed-er-length}, the same displayed expansion has length at least \(2n^{a-o(1)}\).  Thus \(S+rD\ge 2n^{a-o(1)}\).  If \(S,D,r\le L\), then \(S+rD\le L+L^2\), giving \(L^2+L\ge 2n^{a-o(1)}\).  The final lower bound follows by taking square roots and absorbing constants and the lower-order exponent loss.
\end{proof}

\begin{remark}[Macro reuse is the remaining escape]\label{rem:macro-compressed-er-scope}
\Cref{cor:macro-compressed-er} is the concrete macro-library version of \Cref{cor:compressed-er-expansion}.  It does not rule out proof systems whose macros have nonlocal semantics rather than displayed acyclic expansions, nor systems where a small proof denotes an exponentially large object without a controlled expansion accounting.  It says that ordinary macro sharing can compress the displayed proof only by paying a visible reuse product \(rD\).
\end{remark}

\begin{corollary}[Support-declared macro ER normalizer]\label{cor:support-declared-macro-er}
Let \(\lambda\) satisfy \(\HardF_{n,s,A}(\lambda)=1\), and let \(\Phi_{n,s,A}\) be an easy-completion CNF.  Suppose a macro-compressed extension refutation \(\rho\) supplies:
\begin{enumerate}
  \item a compressed macro or extension constraint package \(\Gamma_\rho(Y,W)\), representing its skeleton, macro library, macro invocations, and final clauses;
  \item a resolution refutation of \(\Phi_{n,s,A}(Y,Z)\wedge\Gamma_\rho(Y,W)\);
  \item a declared formula-label support \(S_\rho\subseteq U_n\);
  \item the support-soundness invariant that, for every formula labeling \(\mu\) with \(\mu|_{S_\rho}=\lambda|_{S_\rho}\), there is an auxiliary assignment \(w_\mu\) satisfying \(\Gamma_\rho(\mu,w_\mu)\).
\end{enumerate}
Write \(S_{\rho,b}:=\{u\in S_\rho:\lambda(u)=b\}\).  Then, for each \(b\in\bits\),
\[
  g_n+c\Lookup(M_n,|S_{\rho,b}|,\ell_n)+cM_n>s.
\]
In the near-threshold regime of \Cref{cor:formula-threshold},
\[
  |S_{\rho,0}|,|S_{\rho,1}|\ge n^{a-o(1)}.
\]
Consequently, if the support declaration has budget \(|S_\rho|\le B(|\rho|)\), then
\[
  B(|\rho|)\ge 2n^{a-o(1)}.
\]
\end{corollary}

\begin{proof}
The support-soundness invariant says exactly that \(\Gamma_\rho\) is support-local at \(\lambda\) over \(S_\rho\).  Since \(\Phi_{n,s,A}\wedge\Gamma_\rho\) has a resolution refutation, \Cref{thm:support-local-aux-resolution} applies with \(S=S_\rho\).  Hence \(\lambda|_{S_\rho}\) is a completion-sound formula-label \(1\)-certificate, and \Cref{thm:formula-footprint,cor:formula-threshold} give the two polarity inequalities and near-threshold lower bounds.  The budget statement follows from
\[
  B(|\rho|)\ge |S_\rho|\ge |S_{\rho,0}|+|S_{\rho,1}|\ge 2n^{a-o(1)}.
\]
\end{proof}

\begin{remark}[Support declarations are an obligation]\label{rem:support-declared-macro-er-scope}
\Cref{cor:support-declared-macro-er} is not a lower bound for arbitrary compressed proof systems.  It gives the support-budget counterpart to \Cref{cor:macro-compressed-er}: a compressed macro proof may avoid displaying its whole expansion, but if it claims a local support summary, that summary must itself be large near threshold.  The theorem does not apply to macros with unsound declarations, hidden nonlocal semantics, or proof systems whose compressed witnesses affect unexposed labels without a separate support-local normalizer.
\end{remark}

\begin{corollary}[Support-audited LRAT trace normalizer]\label{cor:support-audited-lrat}
Let \(\lambda\) satisfy \(\HardF_{n,s,A}(\lambda)=1\), and let \(\Phi_{n,s,A}\) be an easy-completion CNF.  Suppose an LRAT-style checked clausal trace \(\tau\) certifies unsatisfiability of
\[
  \Phi_{n,s,A}(Y,Z)\wedge\Gamma_\tau(Y,W),
\]
where \(\Gamma_\tau\) is the collection of checker-visible auxiliary, macro, or support clauses supplied with the trace.  Assume that the trace also carries a support audit: a declared set \(S_\tau\subseteq U_n\) and a verified invariant saying that, for every formula labeling \(\mu\) with \(\mu|_{S_\tau}=\lambda|_{S_\tau}\), there is an auxiliary assignment \(w_\mu\) satisfying \(\Gamma_\tau(\mu,w_\mu)\).  Write \(S_{\tau,b}:=\{u\in S_\tau:\lambda(u)=b\}\).  Then, for each \(b\in\bits\),
\[
  g_n+c\Lookup(M_n,|S_{\tau,b}|,\ell_n)+cM_n>s.
\]
In the near-threshold regime of \Cref{cor:formula-threshold},
\[
  |S_{\tau,0}|,|S_{\tau,1}|\ge n^{a-o(1)}.
\]
Consequently, if the trace or its checker-visible annotation has support budget \(|S_\tau|\le B(|\tau|)\), then
\[
  B(|\tau|)\ge 2n^{a-o(1)}.
\]
\end{corollary}

\begin{proof}
LRAT soundness gives an ordinary clausal refutation of \(\Phi_{n,s,A}\wedge\Gamma_\tau\).  The verified support audit says exactly that \(\Gamma_\tau\) is support-local at \(\lambda\) over \(S_\tau\).  Applying \Cref{thm:support-local-aux-resolution} with \(S=S_\tau\) shows that \(\lambda|_{S_\tau}\) is a completion-sound formula-label certificate.  The polarity inequalities and near-threshold bounds follow from \Cref{thm:formula-footprint,cor:formula-threshold}; the final budget inequality follows from \(|S_\tau|\ge |S_{\tau,0}|+|S_{\tau,1}|\).
\end{proof}

\begin{remark}[LRAT without a support audit is outside the theorem]\label{rem:support-audited-lrat-scope}
\Cref{cor:support-audited-lrat} is a checker-facing version of the support-declared macro normalizer.  It does not prove a lower bound for arbitrary LRAT or DRAT traces.  The new object is the support audit: without a sound, independently checkable reason why all completions preserving \(S_\tau\) satisfy the auxiliary trace constraints, the formula-footprint theorem has no local object to count.
\end{remark}

\begin{corollary}[Provenance-typed LRAT support audit]\label{cor:provenance-typed-lrat}
Let \(\lambda\) satisfy \(\HardF_{n,s,A}(\lambda)=1\), and let \(\Phi_{n,s,A}\) be an easy-completion CNF.  Suppose an LRAT-style checked trace, or a DRAT trace accompanied by an LRAT certificate, refutes
\[
  \Phi_{n,s,A}(Y,Z)\wedge\bigwedge_{j\in J}\Gamma_j(Y,W_j).
\]
Assume that the trace carries a trusted provenance package with the following data:
\begin{enumerate}
  \item a variable-role map identifying the formula-label variables \(Y_u\) and the auxiliary variables;
  \item pairwise disjoint auxiliary blocks \(W_j\);
  \item a declared formula-label support \(P_j\subseteq U_n\) for each auxiliary block \(\Gamma_j\);
  \item a checked local witness invariant saying that, for every formula labeling \(\mu\) with \(\mu|_{P_j}=\lambda|_{P_j}\), there is an assignment \(w_j(\mu)\) satisfying \(\Gamma_j(\mu,w_j(\mu))\);
  \item a provenance rule that every derived clause is assigned the union of the provenance sets of its antecedents, with final declared support \(S_\tau:=\bigcup_{j\in J}P_j\).
\end{enumerate}
Write \(S_{\tau,b}:=\{u\in S_\tau:\lambda(u)=b\}\).  Then, for each \(b\in\bits\),
\[
  g_n+c\Lookup(M_n,|S_{\tau,b}|,\ell_n)+cM_n>s.
\]
In the near-threshold regime of \Cref{cor:formula-threshold},
\[
  |S_{\tau,0}|,|S_{\tau,1}|\ge n^{a-o(1)}.
\]
Consequently, if the trace or checker-visible provenance annotation has support budget \(|S_\tau|\le B(|\tau|)\), then
\[
  B(|\tau|)\ge 2n^{a-o(1)}.
\]
\end{corollary}

\begin{proof}
It is enough to show that the provenance package supplies the support audit required by \Cref{cor:support-audited-lrat}.  Let \(\mu\) be any formula labeling with \(\mu|_{S_\tau}=\lambda|_{S_\tau}\).  Since \(P_j\subseteq S_\tau\) for every \(j\), the same \(\mu\) agrees with \(\lambda\) on every \(P_j\).  The local witness invariant gives an assignment \(w_j(\mu)\) satisfying \(\Gamma_j(\mu,w_j(\mu))\).  Because the auxiliary blocks \(W_j\) are pairwise disjoint, these assignments combine into one assignment \(w_\mu\) satisfying all clauses in \(\bigwedge_j\Gamma_j(Y,W_j)\).  Thus the whole auxiliary package is support-local at \(\lambda\) over \(S_\tau\).

LRAT soundness, or the supplied LRAT certificate for the DRAT trace, gives the required clausal refutation.  Applying \Cref{cor:support-audited-lrat} yields the polarity inequalities, the near-threshold support lower bounds, and the final budget inequality.
\end{proof}

\begin{remark}[Typed provenance versus raw traces]\label{rem:provenance-typed-lrat-scope}
\Cref{cor:provenance-typed-lrat} closes the trusted-provenance route, not the raw-trace route.  It does not extract formula-label support from untyped LRAT or DRAT syntax.  If auxiliary blocks share variables, a separate compatibility proof or global witness constructor must replace the disjoint-block assumption.  Without such typed provenance or compatibility data, \Cref{prop:support-erasure-indistinguishability} still applies.
\end{remark}

\begin{corollary}[Namespace-typed LRAT normalizer]\label{cor:namespace-typed-lrat}
Let \(\lambda\) satisfy \(\HardF_{n,s,A}(\lambda)=1\), and let \(\Phi_{n,s,A}\) be an easy-completion CNF.  Suppose an LRAT-style checked trace refutes
\[
  \Phi_{n,s,A}(Y,Z)\wedge\Gamma(Y,W),
\]
where the input uses a trusted syntactic namespace distinguishing formula-label variables \(Y_u\) from auxiliary variables \(W\).  Let \(\mathcal A_\tau\) be the LRAT antecedent closure of the final empty clause, and let \(\Gamma_\tau\) be the original auxiliary clauses from \(\Gamma\) that occur in \(\mathcal A_\tau\).  Define
\[
  S_\tau:=\{u\in U_n: Y_u \text{ occurs in some clause of }\Gamma_\tau\}.
\]
Assume there is an auxiliary assignment \(w_\lambda\) satisfying \(\Gamma_\tau(\lambda,w_\lambda)\).  Write \(S_{\tau,b}:=\{u\in S_\tau:\lambda(u)=b\}\).  Then, for each \(b\in\bits\),
\[
  g_n+c\Lookup(M_n,|S_{\tau,b}|,\ell_n)+cM_n>s.
\]
In the near-threshold regime of \Cref{cor:formula-threshold},
\[
  |S_{\tau,0}|,|S_{\tau,1}|\ge n^{a-o(1)}.
\]
Consequently, if the trace length or checker-visible namespace accounting bounds \(|S_\tau|\le B(|\tau|)\), then
\[
  B(|\tau|)\ge 2n^{a-o(1)}.
\]
\end{corollary}

\begin{proof}
Clauses outside the LRAT antecedent closure of the empty clause are unused by the checked derivation, so the same trace restricted to \(\mathcal A_\tau\) gives an LRAT refutation of \(\Phi_{n,s,A}\wedge\Gamma_\tau\).  It remains to show that \(S_\tau\) is a sound support audit for \(\Gamma_\tau\).

Let \(\mu\) be any formula labeling with \(\mu|_{S_\tau}=\lambda|_{S_\tau}\).  Every formula-label variable that occurs in \(\Gamma_\tau\) is indexed by an element of \(S_\tau\).  Therefore every label literal in \(\Gamma_\tau(\mu,w_\lambda)\) has the same truth value as in \(\Gamma_\tau(\lambda,w_\lambda)\), while the auxiliary assignment is unchanged.  Since \(w_\lambda\) satisfies \(\Gamma_\tau(\lambda,w_\lambda)\), the same \(w_\lambda\) satisfies \(\Gamma_\tau(\mu,w_\lambda)\).  Thus \(\Gamma_\tau\) is support-local at \(\lambda\) over \(S_\tau\).

The support-audited LRAT theorem \(\Cref{cor:support-audited-lrat}\) now applies to the restricted trace and support \(S_\tau\), giving the polarity inequalities, near-threshold support bounds, and final budget lower bound.
\end{proof}

\begin{remark}[Namespace typing is weaker than full provenance but stronger than raw syntax]\label{rem:namespace-typed-lrat-scope}
\Cref{cor:namespace-typed-lrat} removes the need for a separate derived-clause provenance annotation in the special case where the proof log has a trusted formula-label namespace and the used auxiliary clauses have a hard witness.  It still does not make plain untyped LRAT or DRAT a support normalizer: erasing the namespace returns to the ambiguity in \Cref{prop:support-erasure-indistinguishability}, and clauses whose auxiliary part has no hard witness still require a separate support-locality proof.
\end{remark}

\begin{proposition}[Support-audit necessity separator]\label{prop:support-audit-necessity}
Let \(T\subseteq U_n\) be finite and let \(\alpha:T\to\bits\).  Introduce formula-label variables \(\{Y_u:u\in T\}\).  Let \(\Gamma_T\) be the unit CNF fixing each \(Y_u\) to \(\alpha(u)\), and let \(\Phi_T\) be the single blocking clause falsified exactly by that joint assignment:
\[
  \Phi_T
  :=
  \bigvee_{u\in T:\alpha(u)=1}\neg Y_u
  \ \vee\
  \bigvee_{u\in T:\alpha(u)=0}Y_u .
\]
Then \(\Phi_T\wedge\Gamma_T\) has an LRAT-style clausal refutation of length \(O(|T|)\).  However, if a declared support \(S\subseteq T\) is sound in the support-local sense that every labeling \(\mu\) with \(\mu|_S=\alpha|_S\) satisfies \(\Gamma_T(\mu)\), then \(T\subseteq S\).
\end{proposition}

\begin{proof}
Resolve the blocking clause successively against the unit clause for each \(u\in T\).  Each resolution step removes the corresponding literal, so after \(|T|\) steps one derives the empty clause.  This derivation is a standard linear clausal refutation and can be written with LRAT-style antecedent hints.

For the support claim, suppose \(u\in T\setminus S\).  Define \(\mu\) to agree with \(\alpha\) on \(S\), to flip \(\alpha(u)\) at \(u\), and arbitrarily on \(T\setminus(S\cup\{u\})\).  Then \(\mu|_S=\alpha|_S\), but the unit clause in \(\Gamma_T\) fixing \(Y_u\) to \(\alpha(u)\) is false under \(\mu\).  Since \(\Gamma_T\) has no auxiliary variables, no witness can repair this failure.  Hence a sound support declaration must contain \(u\).  As \(u\) was arbitrary, \(T\subseteq S\).
\end{proof}

\begin{remark}[What the separator says]\label{rem:support-audit-necessity-scope}
\Cref{prop:support-audit-necessity} does not prove a lower bound for unaudited LRAT or DRAT.  It proves the narrower boundary needed by \Cref{cor:support-audited-lrat}: clausal refutation checking and completion-local support checking are different obligations.  A short trace may certify unsatisfiability while every sound support audit has to expose the whole set of labels fixed by the auxiliary constraints.
\end{remark}

\begin{proposition}[Support-erasure indistinguishability separator]\label{prop:support-erasure-indistinguishability}
Let \(T\) be finite and let \(\alpha:T\to\bits\).  Let \(C_T\) be the clausal skeleton over variables \(\{P_u:u\in T\}\) consisting of the unit clauses fixing each \(P_u\) to \(\alpha(u)\) and the single blocking clause falsified exactly by that joint assignment.  The same LRAT-style linear resolution trace refutes \(C_T\) under both of the following semantic lifts.

In the auxiliary lift, the variables \(P_u\) are auxiliary variables unrelated to formula-label coordinates.  The constraint package has empty formula-label support: for every formula labeling \(\mu\), setting \(P_u=\alpha(u)\) satisfies the unit constraints.

In the label lift, the variables \(P_u\) are the formula-label variables \(Y_u\) for \(u\in T\).  Then every sound support declaration for the unit constraints must contain \(T\).  Consequently the untyped clausal skeleton and its LRAT antecedent hints do not determine the completion-local formula-label support.
\end{proposition}

\begin{proof}
The syntactic refutation is the same in both lifts.  Resolve the blocking clause successively against the unit clauses fixing \(P_u\) to \(\alpha(u)\); after all variables in \(T\) have been resolved away, the empty clause is derived.  This gives one LRAT-style trace with the same clauses and antecedent hints.

For the auxiliary lift, the variables \(P_u\) are witnesses.  No formula-label coordinate is constrained: for every labeling of the formula variables, choose the auxiliary assignment \(P_u=\alpha(u)\).  Thus the unit part of the trace is support-local over the empty formula-label support.

For the label lift, identify \(P_u\) with \(Y_u\).  If a declared support \(S\subseteq T\) omits some \(u\), choose a completion \(\mu\) agreeing with \(\alpha\) on \(S\) and flipping \(u\).  The unit clause fixing \(Y_u\) to \(\alpha(u)\) is then false, with no auxiliary witness able to repair it.  Hence every sound support declaration contains all of \(T\).

The raw clausal refutation is identical in the two cases; only the erased role map saying which variables are formula labels and which are auxiliaries differs.  Therefore support locality is not certified by untyped LRAT syntax alone.
\end{proof}

\begin{remark}[Support provenance]\label{rem:support-erasure-scope}
\Cref{prop:support-erasure-indistinguishability} is not an impossibility theorem for every unannotated clausal proof.  It says that a trace checker which sees only untyped clauses and antecedent hints has not checked the semantic support object needed by \Cref{cor:support-audited-lrat}.  A trusted variable-role map, clause provenance, or support-locality checker supplies extra information; that extra information is precisely the support audit being charged.
\end{remark}

\begin{proposition}[Plain LRAT extraction trichotomy]\label{prop:plain-lrat-extraction-trichotomy}
Call an extractor \emph{plain} if its input is only the untyped clausal trace: clause strings together with LRAT antecedent hints, but no variable-role map, clause provenance, support declaration, or support-locality checker.  Suppose such an extractor returns a set \(S(\tau)\) of variables and claims that, whenever \(\tau\) is accepted, \(S(\tau)\) is a sound completion-local formula-label support for the semantic instance represented by the trace.

For the unit-plus-blocking skeleton \(C_T\) from \Cref{prop:support-erasure-indistinguishability}, exactly one of the following must occur.
\begin{enumerate}
  \item The extractor uses extra semantic information, such as a variable-role map, clause provenance, or support-locality checker.  Then it is not plain; it has supplied the additional audit object charged in \Cref{cor:support-audited-lrat,cor:provenance-typed-lrat}.
  \item The extractor is sound for every semantic lift of the raw trace.  Then its output on \(C_T\) contains all variables in \(T\).
  \item The extractor outputs a strict subset of \(T\) from the raw trace alone.  Then the label lift in \Cref{prop:support-erasure-indistinguishability} falsifies its claimed support soundness.
\end{enumerate}
Consequently, plain unaudited LRAT or DRAT syntax alone does not give a nontrivial support-extraction normalizer for the present formula-footprint theorem.  A positive proof-system application must either charge provenance or prove a separate system-specific invariant that rules out the bad semantic lift.
\end{proposition}

\begin{proof}
The auxiliary and label lifts in \Cref{prop:support-erasure-indistinguishability} have the same untyped clauses and the same LRAT antecedent hints.  A plain extractor therefore receives the same input in both lifts and must return the same set \(S(\tau)\).

In the label lift, the variables of the skeleton are the formula-label variables \(Y_u\).  By the support claim in \Cref{prop:support-audit-necessity}, every sound completion-local support for the unit constraints fixing those labels must contain \(T\).  Hence, if \(S(\tau)\) is sound for every semantic lift, \(T\subseteq S(\tau)\).

If instead \(S(\tau)\) omits some \(u\in T\), choose the label lift and flip \(Y_u\) while preserving all coordinates in \(S(\tau)\), as in the proof of \Cref{prop:support-audit-necessity}.  The corresponding unit clause is falsified and no auxiliary witness repairs it.  Thus the claimed support is unsound for that semantic lift.

The only way for the extractor to return different supports on the auxiliary and label lifts is to use information not present in the raw trace, namely a role map, provenance, support declaration, or an equivalent semantic invariant.  That is the first case and is precisely the additional object charged by the audited/provenance-typed LRAT theorems.
\end{proof}

\begin{remark}[What remains after the trichotomy]\label{rem:plain-lrat-trichotomy-scope}
\Cref{prop:plain-lrat-extraction-trichotomy} is a no-free-normalizer statement, not a lower bound for every LRAT or DRAT refutation.  It closes only the route that tries to read nontrivial formula-label support from raw untyped trace syntax.  A proof format may still enter the theorem package by supplying trusted provenance, by making support-locality checker-visible, or by proving a system-specific invariant that prevents the ambiguous semantic lift.
\end{remark}

\begin{corollary}[Cone-budget extension lower bound]\label{cor:cone-budget-extension}
In the setting of \Cref{thm:dependency-cone-extension-normalizer}, suppose the final CNF \(\Theta\) mentions at most \(q\) distinct label or extension variables, and every mentioned variable \(X\) has declared cone size \(|\Delta(X)|\le d\).  Then
\[
  qd\ge |S_\Theta|\ge |S_{\Theta,0}|+|S_{\Theta,1}|.
\]
In the near-threshold regime of \Cref{cor:formula-threshold},
\[
  qd\ge 2n^{a-o(1)}.
\]
Consequently a cone-annotated extension proof with \(q\le L\) final variables and cone bound \(d\le D(L)\) must satisfy
\[
  L\,D(L)\ge 2n^{a-o(1)}.
\]
\end{corollary}

\begin{proof}
The support \(S_\Theta\) is the union of at most \(q\) cones, each of size at most \(d\), so \(|S_\Theta|\le qd\).  The sets \(S_{\Theta,0}\) and \(S_{\Theta,1}\) are disjoint subsets of \(S_\Theta\), and \Cref{thm:dependency-cone-extension-normalizer} gives the near-threshold lower bound \(|S_{\Theta,0}|,|S_{\Theta,1}|\ge n^{a-o(1)}\).  The final inequality follows by substituting \(q\le L\) and \(d\le D(L)\).
\end{proof}

\begin{corollary}[Bounded-depth extension tradeoff]\label{cor:bounded-depth-extension}
In the setting of \Cref{cor:cone-budget-extension}, suppose every extension definition has fan-in at most \(r\), and every label or extension variable mentioned in the final CNF \(\Theta\) has extension depth at most \(h\), where label variables have depth \(0\).  If \(\Theta\) mentions at most \(q\) distinct label or extension variables, then in the near-threshold regime of \Cref{cor:formula-threshold},
\[
  q\,r^h\ge 2n^{a-o(1)}.
\]
Consequently, if \(q\le L\), then
\[
  h\ge \log_r\!\left(\frac{2n^{a-o(1)}}{L}\right)
\]
whenever the right-hand side is positive.
\end{corollary}

\begin{proof}
A label variable has cone size \(1=r^0\).  If an extension variable at depth \(t\) is computed from at most \(r\) earlier variables, and every input variable has cone size at most \(r^{t-1}\), then its cone has size at most \(r\cdot r^{t-1}=r^t\).  Induction on the extension depth therefore gives \(|\Delta(X)|\le r^h\) for every variable \(X\) mentioned in \(\Theta\).  Apply \Cref{cor:cone-budget-extension} with \(d=r^h\).  The logarithmic form is just the rearranged inequality when the quotient is larger than \(1\).
\end{proof}

\begin{definition}[Support-local normalization with support budget]\label{def:support-local-normalizer}
Fix an easy-completion CNF \(\Phi_{n,s,A}\).  A proof or certificate system \(\mathcal P\) has a \emph{support-local normalization with support budget \(B\)} for a hard labeling \(\lambda\) if every accepted \(\mathcal P\)-proof \(\pi\) of the corresponding lower-bound statement yields:
\begin{enumerate}
  \item a support set \(S_\pi\subseteq U_n\);
  \item a CNF \(\Gamma_\pi(Y,W)\) that is support-local at \(\lambda\) over \(S_\pi\);
  \item a sound refutation, in the underlying refutation calculus, of
  \[
    \Phi_{n,s,A}(Y,Z)\wedge\Gamma_\pi(Y,W);
  \]
  \item the size control \(|S_\pi|\le B(|\pi|)\), where \(|\pi|\) is the measured proof length.
\end{enumerate}
\end{definition}

\begin{theorem}[Support-local normalization transfer]\label{thm:support-local-normalization-transfer}
Let \(\lambda\) satisfy \(\HardF_{n,s,A}(\lambda)=1\).  Suppose a proof or certificate system \(\mathcal P\) has a support-local normalization with support budget \(B\) for \(\lambda\).  For every accepted proof \(\pi\), write
\[
  S_{\pi,b}:=\{u\in S_\pi:\lambda(u)=b\}.
\]
Then, for each \(b\in\bits\),
\[
  g_n+c\Lookup(M_n,|S_{\pi,b}|,\ell_n)+cM_n>s.
\]
In the near-threshold regime of \Cref{cor:formula-threshold},
\[
  |S_{\pi,0}|,|S_{\pi,1}|\ge n^{a-o(1)},
  \qquad
  B(|\pi|)\ge 2n^{a-o(1)}.
\]
In particular, if \(B(L)=O(L)\), then every such accepted proof has length \(n^{a-o(1)}\).
\end{theorem}

\begin{proof}
Let \(\Gamma_\pi\) and \(S_\pi\) be supplied by the normalizer.  The sound refutation of \(\Phi_{n,s,A}\wedge\Gamma_\pi\) implies that this CNF is unsatisfiable.  The proof of \Cref{thm:support-local-aux-resolution} used only this unsatisfiability, the easy-completion semantics of \(\Phi_{n,s,A}\), and support-locality of \(\Gamma_\pi\).  Therefore \(\lambda|_{S_\pi}\) is a \(1\)-certificate for \(\HardF_{n,s,A}\), and \Cref{thm:formula-footprint,cor:formula-threshold} give the two polarity inequalities and the two near-threshold lower bounds.  Since \(S_{\pi,0}\) and \(S_{\pi,1}\) are disjoint subsets of \(S_\pi\), the support budget satisfies
\[
  B(|\pi|)\ge |S_\pi|\ge |S_{\pi,0}|+|S_{\pi,1}|\ge 2n^{a-o(1)}.
\]
The final statement is immediate when \(B(L)=O(L)\).
\end{proof}

\begin{corollary}[Carrier-independent audit upgrade criterion]\label{cor:carrier-independent-audit-upgrade}
Let \(\mathcal P\) be a recognizable proof or certificate system for the
easy-completion lower-bound statements \(\Phi_{n,s,A}\).  Suppose there is a
polynomial-time, proof-system-intrinsic audit extractor
\[
  \mathsf{Aud}_{\mathcal P}(\pi,n)
  =
  (S_\pi,\Gamma_\pi,\mathcal W_\pi)
\]
with the following properties for every accepted \(\mathcal P\)-proof \(\pi\).
\begin{enumerate}[label=(\roman*)]
  \item \(S_\pi\subseteq U_n\) is a set of formula-label coordinates computed
  from the proof object and the public proof-system syntax, not from a
  carrier-specific namespace manifest, hidden role advice, or semantic oracle.
  \item \(\Gamma_\pi(Y,W)\) is support-local at the hard labeling \(\lambda\)
  over \(S_\pi\), with \(\mathcal W_\pi\) certifying this support-locality by
  exposed controller completeness, local equivalence, dependency cones,
  provenance, or another system-intrinsic invariant.
  \item The underlying checker verifies a sound refutation of
  \[
    \Phi_{n,s,A}(Y,Z)\wedge\Gamma_\pi(Y,W).
  \]
  \item A budget \(B(|\pi|)\) upper-bounds the extracted support and audit
  data:
  \[
    |S_\pi|+|\mathcal W_\pi|\le B(|\pi|).
  \]
\end{enumerate}
Then \(\mathcal P\) satisfies the support-local normalization hypothesis of
\(\Cref{thm:support-local-normalization-transfer}\).  In particular, in the
near-threshold regime of \(\Cref{cor:formula-threshold}\),
\[
  B(|\pi|)\ge 2n^{a-o(1)}.
\]
If \(B(L)=O(L^k)\), then every accepted proof has length
\[
  |\pi|\ge n^{a/k-o(1)}.
\]
\end{corollary}

\begin{proof}
The data output by \(\mathsf{Aud}_{\mathcal P}\) are exactly the four pieces in
the definition of support-local normalization: a support set, a support-local
CNF, a sound refutation after adjoining that CNF to the easy-completion
formula, and a size budget.  The additional requirement that the audit is
proof-system-intrinsic records what is needed for a top-paper proof-system
upgrade: the support object must be forced by the proof format itself, rather
than imported through a carrier-specific manifest or hidden role annotation.
Thus \(\Cref{thm:support-local-normalization-transfer}\) applies with the same
budget \(B\), giving \(B(|\pi|)\ge 2n^{a-o(1)}\).  If \(B(L)=O(L^k)\), then
\(O(|\pi|^k)\ge 2n^{a-o(1)}\), which rearranges to the stated proof-length
lower bound.

This is an upgrade criterion, not an instantiation for Frege, Extended Frege,
IPS, arbitrary compressed proofs, or raw proof traces.  To use it for such a
system one must still prove the system-intrinsic audit extractor and its
support-locality certificate.
\end{proof}

\begin{corollary}[Self-auditing RUP-CDCL instantiation]\label{cor:self-auditing-rup-cdcl}
Consider an enriched RUP-CDCL proof format for the easy-completion CNF
\(\Phi_{n,s,A}\) in which the proof object contains:
\begin{enumerate}[label=(\roman*)]
  \item a checked, proof-internal label decoder mapping formula-label
  coordinates \(u\in U_n\) to the Boolean variables representing them in the
  clausal encoding;
  \item a live-database reconstruction log for clause additions and deletions;
  \item for every learned or deleted-clause consequence used in the final
  replay, a checker-visible RUP propagation certificate over the reconstructed
  live database;
  \item retained support annotations obtained by intersecting every touched
  clause in those RUP certificates with the decoded formula-label variables.
\end{enumerate}
Assume the checker verifies that all auxiliary clauses touched by the replay
are satisfiable under every completion preserving the retained support, and
that the final reconstructed database refutes
\(\Phi_{n,s,A}\wedge\Gamma_\pi\).  Then this self-auditing RUP-CDCL format
instantiates the carrier-independent audit upgrade criterion with linear audit
budget.  Hence, in the near-threshold regime of
\(\Cref{cor:formula-threshold}\), every accepted proof in this format has
length \(n^{a-o(1)}\).
\end{corollary}

\begin{proof}
The extractor is the generic checker for the enriched proof syntax.  From the
proof-internal decoder and live-database log it computes the formula-label
variables touched by the RUP certificates used in the final replay.  Their
coordinate preimages form \(S_\pi\).  The clauses whose touched support is
contained in this set form \(\Gamma_\pi\), while the RUP propagation
certificates, live-database reconstruction, deleted-clause support records, and
support-local auxiliary witnesses form \(\mathcal W_\pi\).  All of this data is
read from the proof artifact itself, with no external carrier-specific
namespace manifest, hidden role advice, or semantic oracle.

The soundness and support-locality assumptions are precisely the hypotheses
needed by the RUP-complete and canonical-namespace corollaries,
\Cref{cor:rup-complete-cdcl,cor:canonical-namespace-rup-cdcl}, except that the
namespace parser is internalized into the proof format rather than supplied as a
separate manifest.  Therefore the extractor returns the data required by
\Cref{cor:carrier-independent-audit-upgrade}.  Because every
support element, propagation step, live-database event, and retained
deleted-clause support record is explicitly present in the proof object, the
audit budget is \(B(L)=O(L)\).  Applying
\(\Cref{cor:carrier-independent-audit-upgrade}\) gives the near-threshold
support lower bound and hence proof length \(n^{a-o(1)}\).

This is not a lower bound for raw DIMACS, raw RUP/DRUP, arbitrary CDCL, Frege,
Extended Frege, IPS, or an unannotated compressed proof system.  Its role is to
show that the carrier-independent criterion is an honest proof-format theorem
once the proof system itself carries the label decoder and support audit that
the separators prove cannot be erased for free.
\end{proof}

\begin{corollary}[Append-only self-auditing RUP-CDCL]\label{cor:append-only-self-auditing-rup-cdcl}
Consider the self-auditing RUP-CDCL format of
\Cref{cor:self-auditing-rup-cdcl} under the additional syntactic restriction
that the reconstructed clause database is append-only: once a clause is added,
it remains available until the final empty-clause replay.  Then the proof object
need not carry separate retained deleted-clause support records.  It is enough
that it contains the checked proof-internal label decoder, the append-only clause-addition log, and checker-visible RUP propagation certificates for the
used learned consequences.  If the checker verifies support-locality for the
touched auxiliary clauses relative to the decoded formula-label support, then
this deletion-free format instantiates the carrier-independent audit upgrade
criterion with linear audit budget.  Hence, in the near-threshold regime of
\Cref{cor:formula-threshold}, every accepted proof in this append-only
self-auditing RUP-CDCL format has length \(n^{a-o(1)}\).
\end{corollary}

\begin{proof}
Run the extractor from \Cref{cor:self-auditing-rup-cdcl}, but compute support by
a single scan over the monotone clause database and the RUP propagation
certificates used in the final replay.  Since no clause is deleted, every clause
that can participate in a later propagation remains present in the append-only
log.  The support of a used clause is therefore just the set of decoded
formula-label variables occurring in that clause, and no retained
deleted-clause support record is needed to remember information that vanished
from the live database.

The support set \(S_\pi\), support-local CNF \(\Gamma_\pi\), RUP witnesses, and
support-local auxiliary witnesses are all read from the ordinary append-only
proof artifact plus its internal decoder.  The size of this audit data is linear
in the proof length.  Thus the hypotheses of
\Cref{cor:carrier-independent-audit-upgrade} hold with \(B(L)=O(L)\), giving the
near-threshold support lower bound and hence proof length \(n^{a-o(1)}\).

This removes one annotation layer from
\Cref{cor:self-auditing-rup-cdcl}; it does not remove the need for a checked
label decoder, RUP propagation witnesses, or support-locality checks.  In
particular, it is still not a lower bound for raw DIMACS, arbitrary deletion
traces, provenance-free CDCL, Frege, Extended Frege, IPS, or unannotated
compressed proofs.
\end{proof}

\begin{corollary}[Prefix-typed append-only RUP-CDCL]\label{cor:prefix-typed-append-only-rup-cdcl}
Fix a typed append-only RUP-CDCL proof system for the easy-completion CNF
\(\Phi_{n,s,A}\) in which the formula-label variables form a public syntactic
class: from a variable name alone the checker can decide whether it is a
formula-label variable and, if so, recover its coordinate \(u\in U_n\).  Suppose
the proof object contains an append-only clause-addition log, checker-visible
RUP propagation certificates for every used learned consequence, and
support-local auxiliary witnesses for every touched auxiliary clause.  Then no
proof-internal label decoder and no retained deleted-clause support records are
needed.  This typed append-only RUP-CDCL system instantiates the
carrier-independent audit upgrade criterion with linear audit budget.  Hence,
in the near-threshold regime of \Cref{cor:formula-threshold}, every accepted
proof in this format has length \(n^{a-o(1)}\).
\end{corollary}

\begin{proof}
The public typed namespace is itself the decoder.  The checker scans the
append-only clause-addition log and the RUP propagation certificates used in
the final replay, collects exactly the variables belonging to the public
formula-label syntactic class, and maps them to their coordinates in \(U_n\).
This gives the support set \(S_\pi\) without any proof-specific decoder
annotation.

Because the clause database is append-only, no later replay step can depend on
a deleted clause whose formula-label support must be remembered separately.
Thus the monotone scan also recovers the support-local CNF \(\Gamma_\pi\), the
RUP witnesses, and the touched auxiliary witnesses with total size linear in
the proof length.  The assumed support-local auxiliary witnesses certify that
the touched auxiliary clauses are satisfiable under every completion preserving
\(S_\pi\), and the RUP replay certifies the final clausal contradiction for
\(\Phi_{n,s,A}\wedge\Gamma_\pi\).  Therefore the hypotheses of
\Cref{cor:carrier-independent-audit-upgrade} hold with \(B(L)=O(L)\), giving
the claimed lower bound.

The theorem removes the proof-internal label decoder only because the typed
proof system makes the formula-label namespace public and checkable.  It does
not prove a lower bound for untyped DIMACS, raw deletion-bearing RUP/DRUP,
arbitrary CDCL, provenance-free learned clauses, Frege, Extended Frege, IPS, or
arbitrary compressed proofs, and it does not remove the support-locality
witness obligation.
\end{proof}

\begin{proposition}[Support-local witnesses are necessary for prefix-typed append-only RUP-CDCL]\label{prop:prefix-typed-rup-support-witness-necessity}
In \Cref{cor:prefix-typed-append-only-rup-cdcl}, the support-local auxiliary
witness hypothesis cannot be replaced by append-only RUP checking plus public
formula-label typing alone.  More precisely, there is no sound extractor that,
for every typed append-only RUP-CDCL trace, scans only the public formula-label
variables occurring in the append-only log and RUP certificates and then
certifies completion-local support for touched auxiliary clauses without an
additional support-locality witness or an equivalent proof-system invariant.
\end{proposition}

\begin{proof}
This is the same-skeleton obstruction from
\Cref{prop:support-audit-necessity,prop:support-erasure-indistinguishability}
with the variable-role ambiguity removed.  Let \(W\) be the formula-label
support obtained by the public typed scan of the append-only log and the RUP
certificates.  Add a touched auxiliary clause, or a bounded auxiliary gadget,
whose displayed clausal use and RUP replay are unchanged but whose semantic
soundness depends on a hidden formula-label coordinate \(h\notin W\).  There
are two completions that agree on every coordinate in \(W\) and differ only at
\(h\).  The typed append-only trace, the public support scan, and the RUP
certificates are identical under the two completions, but the auxiliary clause
is completion-sound in one completion and not in the other.

Thus \(W\) is not a completion cylinder support for the touched auxiliary
object.  Public typing identifies formula-label variables that occur in the
trace, append-only logging prevents deleted-clause support from being lost, and
RUP certificates verify displayed clausal implication over the replayed
database.  None of these syntactic checks proves that every touched auxiliary
constraint is satisfiable under all completions preserving \(W\).  That last
fact is exactly what the support-local auxiliary witness, or an equivalent
intrinsic invariant, supplies.  The positive conclusion of
\Cref{cor:prefix-typed-append-only-rup-cdcl} is therefore unchanged, but its
remaining witness assumption is a real hypothesis rather than removable
bookkeeping.
\end{proof}

\begin{corollary}[Intrinsic-locality prefix-typed append-only RUP-CDCL]\label{cor:intrinsic-locality-prefix-typed-rup-cdcl}
Fix a prefix-typed append-only RUP-CDCL proof system as in
\Cref{cor:prefix-typed-append-only-rup-cdcl}, except that the proof object need
not carry separate support-local auxiliary witnesses.  Instead assume that
touched auxiliary clauses are introduced only by a fixed finite family of
public auxiliary rules.  For every rule instance producing an auxiliary clause
\(C\), the checker can compute from the displayed rule name and displayed
premises a finite formula-label support \(T(C)\), and the rule specification
has a system-level soundness theorem saying that \(C\) is satisfiable under
every completion preserving \(T(C)\).  Suppose rule applicability, the sets
\(T(C)\), the append-only clause log, and the RUP certificates are all
checker-visible, and the total displayed size of the touched rule instances is
linear in the proof length.  Then this intrinsic-locality typed append-only
RUP-CDCL system instantiates the carrier-independent audit upgrade criterion
with linear audit budget.  Hence, in the near-threshold regime of
\Cref{cor:formula-threshold}, every accepted proof in this format has length
\(n^{a-o(1)}\).
\end{corollary}

\begin{proof}
The public formula-label namespace gives the decoder exactly as in
\Cref{cor:prefix-typed-append-only-rup-cdcl}.  The append-only log and RUP
certificates give the used clausal replay and formula-label support without
deleted-clause support records.  For each touched auxiliary clause \(C\), the
checker reads the public rule name and premises, computes \(T(C)\), and invokes
the fixed rule-soundness theorem.  Thus the support-local auxiliary witness
required in \Cref{cor:prefix-typed-append-only-rup-cdcl} is not an extra
per-proof annotation: it is generated canonically by the proof system's own
rule discipline.

Let \(S_\pi\) be the union of the formula-label variables scanned from the
append-only clausal replay, the RUP certificates, and all sets \(T(C)\) for
touched auxiliary clauses.  If two completions agree on \(S_\pi\), then every
touched auxiliary rule instance remains satisfiable by its fixed
completion-local soundness theorem, and the RUP replay remains the same
clausal contradiction over the support-local CNF.  The total audit data is
linear because the rule instances, computed supports, append-only log, and RUP
certificates are all checker-visible with linear total displayed size.
Therefore \Cref{cor:carrier-independent-audit-upgrade} applies with
\(B(L)=O(L)\), and the near-threshold lower bound follows.

This corollary is the positive counterpart to
\Cref{prop:prefix-typed-rup-support-witness-necessity}.  Public label typing
and RUP checking alone still do not certify auxiliary support-locality; the
new ingredient is the fixed, public, completion-stable auxiliary rule
discipline.  The result does not cover raw untyped RUP/DRUP, deletion-bearing
traces without retained support, arbitrary CDCL, or proof systems whose
auxiliary clauses are justified only by hidden semantic choices.
\end{proof}

\begin{proposition}[Strict-prefix Qrep replay normalizer]\label{prop:strict-prefix-qrep}
Let \(\pi\) be a finite proof transcript with an ordinary displayed support set \(K_0(\pi)\).  Suppose the only compressed support steps in \(\pi\) are quotient-representative dictionary entries
\[
  e=(a_e,r_e,F_e,t_e,u_e,\operatorname{rank}(e)),
\]
where \(a_e\) is the row or local object replaced by representative \(r_e\), \(F_e\) is a finite same-theory certificate forest proving that \(a_e\) and \(r_e\) have the same label or semantic value in the declared local theory, \(t_e\) is the line where \(F_e\) is displayed, and \(u_e\) is the first line using \(r_e\) in place of \(a_e\).  Assume every used entry is strict-prefix and well-founded:
\begin{enumerate}
  \item \(t_e<u_e\);
  \item every dictionary dependency in \(F_e\) has strictly smaller rank and has already been accepted in the prefix before \(t_e\);
  \item dictionary entries are append-only, and no representative is rebound after first use;
  \item the entries, ranks, dependency pointers, certificate forests, and first-use lines are determined by \(N\) and the transcript \(\pi\), with no SAT-truth or row-selection advice.
\end{enumerate}
Then the quotient uses in \(\pi\) have a canonical replay expansion: each representative use can be replaced, in line order and increasing rank, by the representative row together with the already charged forest \(F_e\).  The expanded transcript has local support contained in
\[
  K_{\mathrm{Qrep}}(\pi)
    :=K_0(\pi)\cup
      \bigcup_{e\in\operatorname{Used}(\pi)}
        \bigl(\operatorname{supp}(F_e)\cup\{a_e,r_e\}\bigr),
\]
and its displayed support budget is enlarged by at most a constant factor times
\[
  D_Q(\pi)
    :=\sum_{e\in\operatorname{Used}(\pi)}
       \bigl(|F_e|+|\operatorname{supp}(F_e)|+1\bigr).
\]
Consequently, if a proof-system normalizer is support-local after expanding all non-Qrep steps and its only remaining compression is strict-prefix Qrep, then \Cref{thm:support-local-normalization-transfer} applies with support budget enlarged by \(O(D_Q(\pi))\).

If any first representative use is not strict-prefix, rank-decreasing, dependency-closed, append-only, or advice-free, no transparent-support conclusion follows from this proposition.  In the abstract label model there is a private-row falsifier at the first unsafe use: two completions can agree on the visible prefix while differing on the hidden row, hidden certificate bit, or hidden representative choice that controls the quotient merge.
\end{proposition}

\begin{proof}
Process dictionary entries in proof-prefix order, breaking ties by increasing rank.  The strict inequality \(t_e<u_e\) says that when a representative is first used, its certificate forest is already displayed.  The rank condition says that every quotient dependency needed by that forest has already been replayed or belongs to the ordinary displayed support.  Thus the recursive replay terminates and is independent of later proof lines.  Append-only dictionaries and the no-rebinding condition prevent a later line from changing the meaning of an earlier representative.

For each use of \(r_e\) in place of \(a_e\), splice in the certificate forest \(F_e\) and the displayed representative row.  Since \(F_e\) is a same-theory certificate, the splice preserves the local semantic value used by the original line.  Since all leaves of \(F_e\), together with \(a_e\) and \(r_e\), are included in \(K_{\mathrm{Qrep}}(\pi)\), the replay introduces no uncharged local object.  Summing the certificate forest sizes and newly exposed supports over used entries gives the stated \(O(D_Q(\pi))\) budget increase.

After replay, the transcript is an ordinary displayed-support transcript for purposes of any support-local normalizer.  Applying \Cref{thm:support-local-normalization-transfer} to the replayed proof gives the advertised inherited footprint bound.

For the boundary statement, choose the earliest unsafe representative use.  All earlier lines are fixed.  If the missing condition is late certification, cyclic or nondecreasing dependency, rebinding, or advice-chosen representative selection, introduce a private row or certificate bit not determined by the visible prefix and set two completions to agree on every displayed object before that use while disagreeing on the hidden object controlling the merge.  The two completions make the same visible prefix acceptable but disagree on whether the representative use is sound.  Hence the unexpanded quotient step is not transparent support.
\end{proof}

\begin{example}[Two-line prefix scan for Qrep]\label{ex:qrep-prefix-scan}
The finite checker used in the accompanying artifact has two one-entry transcripts.  In the valid transcript, \(a=\texttt{row\_B}\) is replaced by \(r=\texttt{row\_A}\), the same-theory certificate forest \(F=\{\texttt{cert\_A\_equiv\_B}\}\) is displayed at line \(3\), and the representative is first used at line \(4\).  The prefix scan accepts.  Its support consists of the final lower-bound line, the FormulaEval parent, and the two rows \(\texttt{row\_A},\texttt{row\_B}\), so
\[
  |K_{\mathrm{Qrep}}|=4,
  \qquad
  D_Q=|F|+2+1=4.
\]
In the delayed transcript the same representative is used at line \(4\), but the certificate is displayed only at line \(5\).  The checker rejects with the reason \(\texttt{certificate\_not\_strict\_prefix}\).  This is exactly the temporal private-row falsifier in \Cref{prop:strict-prefix-qrep}: before the first use, the visible prefix has not paid for the same-theory evidence.
\end{example}

\section{Occurrence-label certificates and full rows}\label{sec:rows}

Fix $z$ with $\HardO_{n,s,A}(z)=1$, and let $\sigma:Q\to\bits$ be a $1$-certificate at $z$, where $Q\subseteq\Omega_n$.  For a row $x\in\bits^n$, define
\[
  Q_x:=\{j\in[M_n]:(x,j)\in Q\}
\]
and let $\sigma_x$ be the induced partial assignment on that row.

\begin{definition}[Forced rows]\label{def:forced-rows}
For $b\in\bits$, let
\[
  R_b(\sigma)
  :=\{x\in\bits^n:\sigma_x\text{ forces }A_{M_n}=b\}.
\]
\end{definition}

The sets $R_0(\sigma)$ and $R_1(\sigma)$ are disjoint.

\begin{theorem}[Row-forcing theorem]\label{thm:row-forcing}
There is an absolute constant $c\ge1$ such that every $1$-certificate $\sigma$ for $\HardO_{n,s,A}$ satisfies, for each $b\in\bits$,
\begin{equation}\label{eq:forced-rows}
  |R_b(\sigma)|
  >\frac{s-c}{c(n+1)}.
\end{equation}
Consequently,
\begin{equation}\label{eq:occ-query}
  |Q|
  \ge
  \Force_0(A_{M_n})|R_0(\sigma)|
  +\Force_1(A_{M_n})|R_1(\sigma)|,
\end{equation}
and therefore, for $s=\Omega(n)$,
\begin{equation}\label{eq:occ-query-asymptotic}
  |Q|
  =\Omega\!\left(
    (\Force_0(A_{M_n})+\Force_1(A_{M_n}))\frac{s}{n}
  \right).
\end{equation}
\end{theorem}

\begin{proof}
Fix $b\in\bits$.  For every row $x\notin R_b(\sigma)$, the partial assignment $\sigma_x$ does not force output $b$.  Hence it has a completion whose aggregator value is $1-b$.  Choose such a completion independently for every row outside $R_b(\sigma)$, and complete the rows in $R_b(\sigma)$ arbitrarily.  The resulting total occurrence labeling $z^{(b)}$ extends the certificate and satisfies
\[
  F_{n,A}^{z^{(b)}}(x)=1-b
  \qquad\text{for every }x\notin R_b(\sigma).
\]
On a row in $R_b(\sigma)$, every completion has output $b$.  Thus $F_{n,A}^{z^{(b)}}$ differs from the constant function $1-b$ exactly on $R_b(\sigma)$.

By \Cref{lem:patching},
\[
  \Ckt(F_{n,A}^{z^{(b)}})
  \le c+c(n+1)|R_b(\sigma)|
\]
for an absolute constant $c$.  Since $z^{(b)}$ extends a $1$-certificate, its induced function has circuit complexity greater than $s$.  Rearranging proves \eqref{eq:forced-rows}.

Every row in $R_b(\sigma)$ contains at least $\Force_b(A_{M_n})$ queried coordinates.  The forced-row sets are disjoint, so summing their rowwise query counts proves \eqref{eq:occ-query}.
\end{proof}

\subsection{Examples of aggregator costs}

For the threshold aggregator $T_{M,t}(u)=1$ iff $\sum_i u_i\ge t$,
\[
  \Force_1(T_{M,t})=t,
  \qquad
  \Force_0(T_{M,t})=M-t+1.
\]
This includes OR and AND.  The resulting bounds are shown in \Cref{tab:aggregators}.

\begin{table}[H]
\centering
\small
\begin{tabular}{@{}lccc@{}}
\toprule
Aggregator $A_M$ & $\Force_0(A_M)$ & $\Force_1(A_M)$ & Query lower bound \\
\midrule
$\mathrm{OR}_M$ & $M$ & $1$ & $\Omega((M+1)s/n)$ \\
$\mathrm{AND}_M$ & $1$ & $M$ & $\Omega((M+1)s/n)$ \\
Majority, odd $M$ & $(M+1)/2$ & $(M+1)/2$ & $\Omega(Ms/n)$ \\
Parity $\oplus_M$ & $M$ & $M$ & $\Omega(Ms/n)$ \\
\bottomrule
\end{tabular}
\caption{Minimum forcing numbers and the consequence of \Cref{thm:row-forcing}.  Parity is maximally expensive on both output values.}
\label{tab:aggregators}
\end{table}

The total lower bound for OR is still $\Omega(Ms/n)$ because many rows must be forced to zero, but its positive rows can be certified by one occurrence each.  Parity has no such one-sided shortcut.

\begin{corollary}[Full-row barrier for parity]\label{cor:full-row}
For parity, every $1$-certificate for $\HardO_{n,s,\oplus}$ satisfies
\[
  |Q|=\Omega(M_ns/n).
\]
Moreover, every row in $R_0(\sigma)\cup R_1(\sigma)$ is complete: all $M_n$ coordinates of that row belong to $Q$.
If
\[
  M_n=n^{q+o(1)},
  \qquad
  s_n=n^{1+\eps-o(1)},
\]
then
\[
  |Q|\ge n^{q+\eps-o(1)}.
\]
\end{corollary}

\begin{proof}
Any proper partial assignment to parity leaves a free coordinate whose flip changes the output.  Hence both forcing numbers equal $M_n$.  Apply \Cref{thm:row-forcing}.
\end{proof}

\subsection{Adaptive occurrence-query verifiers}

\begin{proposition}[Occurrence-query lower bound]\label{prop:adaptive-occ}
Let $V$ be deterministic and suppose
\[
  V^z(\pi)=1
  \quad\Longrightarrow\quad
  \HardO_{n,s,A}(z)=1
\]
for every occurrence labeling $z$.  Then the query-answer transcript of every accepted execution satisfies \eqref{eq:occ-query}.  For parity it makes $\Omega(M_ns/n)$ oracle queries.
\end{proposition}

The conclusion is a lower bound on semantic occurrence access.  Turning it into a proof-length lower bound requires a separate statement relating proof length to the number of concrete occurrence labels exposed by the proof system.

\section{Equality quotients and residual affine distance}\label{sec:quotient}

The occurrence model treats all positions as independent.  Equality links can invalidate that relaxation.

\begin{example}[Parity cancellation under an equality link]\label{ex:cancellation}
Suppose a parity row contains two formula strings $u$ and $u'$ and the admissible labelings satisfy $\lambda(u)=\lambda(u')$.  Their contribution is always
\[
  \lambda(u)\oplus\lambda(u')=0,
\]
although neither label is fixed.  Independent occurrence completion would incorrectly allow the pair to contribute $1$.
\end{example}

We therefore quotient labels before taking row parity.

\subsection{The quotient incidence map}

Let $E_n$ be an equivalence relation on $U_n$, and write
\[
  \mathcal C_n:=U_n/E_n
\]
for its set of classes.  For $C\in\mathcal C_n$, define the parity incidence
\[
  B_{E_n}(x,C)
  :=\bigoplus_{j:G_{n,j}(x)\in C}1.
\]
For a quotient labeling $z\in\bits^{\mathcal C_n}$, define
\begin{equation}\label{eq:quotient-carrier}
  H_{n,E_n}^{z}(x)
  :=\bigoplus_{C\in\mathcal C_n}B_{E_n}(x,C)z_C.
\end{equation}
Equivalently, $H_{n,E_n}^{z}=L_{E_n}z$, where
\[
  L_{E_n}:\F^{\mathcal C_n}\longrightarrow\F^{\bits^n}
\]
is the binary incidence map.

Define
\[
  \HardQ_{n,s,E_n}(z)
  :=\mathbf 1[\Ckt(H_{n,E_n}^{z})>s].
\]
Let $\sigma:W\to\bits$ be a $1$-certificate at some hard quotient labeling.  Its set of source completions is the affine set
\begin{equation}\label{eq:affine-completions}
  \mathcal A(\sigma)
  :=\{L_{E_n}z:z|_W=\sigma\}
  \subseteq\F^{\bits^n}.
\end{equation}

\begin{definition}[Residual constant distance]\label{def:residual-distance}
Define
\begin{equation}\label{eq:delta-pm}
  \delta_{\pm}(\sigma)
  :=\min_{h\in\mathcal A(\sigma)}
      \min_{b\in\bits}\dist(h,b\mathbf 1),
\end{equation}
where $\mathbf 1$ denotes the all-one source truth table.
\end{definition}

\begin{theorem}[Quotient residual-distance barrier]\label{thm:quotient-distance}
There is an absolute constant $c\ge1$ such that every $1$-certificate $\sigma$ for $\HardQ_{n,s,E_n}$ satisfies
\begin{equation}\label{eq:quotient-distance}
  \delta_{\pm}(\sigma)
  >\frac{s-c}{c(n+1)}.
\end{equation}
In particular, if $s_n=n^{1+\eps-o(1)}$, then
\[
  \delta_{\pm}(\sigma)\ge n^{\eps-o(1)}.
\]
\end{theorem}

\begin{proof}
Choose $h=L_{E_n}z\in\mathcal A(\sigma)$ and $b\in\bits$ attaining \eqref{eq:delta-pm}.  The quotient labeling $z$ extends the certificate, so $\Ckt(h)>s$.  On the other hand, $h$ differs from the constant function $b$ on exactly $\dist(h,b\mathbf 1)$ source rows.  Equality patching gives
\[
  \Ckt(h)
  \le c+c(n+1)\dist(h,b\mathbf 1).
\]
Rearranging proves the claim.
\end{proof}

\begin{remark}[Coding interpretation]\label{rem:coding}
Before labels are fixed, $\operatorname{im}(L_{E_n})$ is a binary linear code of source truth tables.  A certificate restricts the label vector to an affine subcube, whose image is an affine subset of this code.  \Cref{thm:quotient-distance} says that completion soundness forces the entire residual affine set outside Hamming balls of radius $\Theta(s/n)$ around both constant functions.
\end{remark}

\begin{remark}[Scope of quotienting]\label{rem:quotient-scope}
The theorem handles constraints saying that certain labels are equal.  It does not by itself handle nonlinear semantic laws.  For example, SAT self-reduction imposes
\[
  Y(u)=Y(u|_{x_i=0})\vee Y(u|_{x_i=1}),
\]
which is not captured by an equality quotient.  The next proposition gives the local support-accounting analogue for one relation instance.  A global residual-completion theorem for a whole SAT self-reduction closure remains an open direction discussed in \Cref{sec:related}.
\end{remark}

\begin{proposition}[Semantic-relation support certificate]\label{prop:semantic-relation-support}
Let \(J\) be a finite index set, let \(R\subseteq\bits^J\) be a Boolean relation, and let \(a\in R\).  A set \(S\subseteq J\) is support-local for \(R\) at \(a\) if every \(b\in\bits^J\) with \(b|_S=a|_S\) also lies in \(R\).  Then \(S\) is exactly a standard \(1\)-certificate for membership of \(a\) in \(R\).

For the SAT self-reduction relation
\[
  R_{\vee}:=\{(y,y_0,y_1): y=y_0\vee y_1\},
\]
the tuple \((0,0,0)\) requires all three coordinates in any support-local certificate.  The satisfying tuples \((1,1,0)\), \((1,0,1)\), and \((1,1,1)\) have support-local certificates of size \(2\), namely \(y=1\) together with one child equal to \(1\).
\end{proposition}

\begin{proof}
The first statement is a restatement of the two definitions: \(S\) is support-local exactly when fixing \(a|_S\) forces every completion to remain in \(R\), which is precisely the usual certificate condition for the relation membership predicate.

For \(R_{\vee}\), if \(a=(0,0,0)\) and any one of the three coordinates is omitted, flipping that coordinate while preserving the others violates \(y=y_0\vee y_1\).  Thus all three coordinates are necessary and sufficient.  If \(a=(1,1,0)\) or \((1,1,1)\), fixing \(y=1\) and \(y_0=1\) forces the relation regardless of \(y_1\).  If \(a=(1,0,1)\), fixing \(y=1\) and \(y_1=1\) does the same.
\end{proof}

\begin{proposition}[Acyclic OR-tree residual certificate]\label{prop:acyclic-or-tree-residual}
Let \(T\) be a finite rooted tree whose internal nodes compute OR of their children.  Let \(L(T)\) be the leaf set and let \(r\) be the root.  After eliminating internal nodes, the frontier relation on \(\{r\}\cup L(T)\) is
\[
  R_T
  :=\left\{(y_r,(y_\ell)_{\ell\in L(T)}):
      y_r=\bigvee_{\ell\in L(T)}y_\ell\right\}.
\]
Let \(a\in R_T\), and let \(S\subseteq\{r\}\cup L(T)\) be support-local for \(R_T\) at \(a\).  If \(a_r=0\), then \(r\in S\) and \(L(T)\subseteq S\).  If \(a_r=1\), then \(r\in S\) and \(S\) contains at least one leaf \(\ell\) with \(a_\ell=1\).  Conversely, these conditions are sufficient.
\end{proposition}

\begin{proof}
The frontier relation is the ordinary OR relation on the leaves.  If \(a_r=0\), then every leaf has value \(0\).  If \(r\notin S\), changing only \(y_r\) to \(1\) violates \(R_T\).  If some leaf \(\ell\notin S\), changing only \(y_\ell\) to \(1\) violates \(R_T\).  Thus \(r\) and all leaves are necessary, and fixing them to \(0\) is plainly sufficient.

If \(a_r=1\), then at least one true leaf exists.  If \(r\notin S\), changing \(y_r\) to \(0\) while leaving the true leaf unchanged violates \(R_T\).  If \(S\) contains no true leaf, complete every leaf not in \(S\) to \(0\); the root remains fixed to \(1\), so the relation is violated.  Hence \(S\) must contain \(r\) and a true leaf.  Conversely, fixing \(r=1\) and any leaf \(\ell\) with \(a_\ell=1\) forces the OR relation under all completions.
\end{proof}

\begin{proposition}[OR-DAG frontier residual certificate]\label{prop:or-dag-frontier-residual}
Let \(D\) be a finite acyclic directed graph with root \(r\).  Each non-frontier node computes the OR of its children.  Let \(Q_r\) be the set of distinct quotient classes of frontier coordinates reachable from \(r\), where the quotient identifies only soundly equal frontier formulas.  After eliminating internal OR nodes, the residual frontier relation is
\[
  R_D
  :=\left\{(y_r,(z_q)_{q\in Q_r}):
      y_r=\bigvee_{q\in Q_r} z_q\right\}.
\]
Let \(a\in R_D\), and let \(S\subseteq\{r\}\cup Q_r\) be support-local for \(R_D\) at \(a\).  If \(a_r=0\), then \(r\in S\) and \(Q_r\subseteq S\).  If \(a_r=1\), then \(r\in S\) and \(S\) contains at least one reachable frontier class \(q\in Q_r\) with \(a_q=1\).  Conversely, these conditions are sufficient.
\end{proposition}

\begin{proof}
Eliminate the internal nodes in reverse topological order.  Each internal node computes the OR of the frontier coordinates reachable below it.  Shared subgraphs only repeat the same OR subexpression, and repeated frontier occurrences collapse under idempotence of OR.  Sound frontier quotienting merely replaces repeated equal formulas by one coordinate.  Hence the root relation is exactly the displayed OR over the distinct reachable frontier classes \(Q_r\).

The support-local certificates for this relation are the ordinary OR certificates.  If \(a_r=0\), every \(q\in Q_r\) has \(a_q=0\).  Omitting \(r\) lets us flip \(y_r\) to \(1\), and omitting any reachable frontier class \(q\) lets us flip \(z_q\) to \(1\); either completion violates \(R_D\).  Thus \(r\) and all classes in \(Q_r\) are necessary and plainly sufficient.  If \(a_r=1\), omitting \(r\) lets us flip the root to \(0\) while keeping a true reachable class fixed.  If no true reachable class is included in \(S\), complete all unexposed reachable classes to \(0\), leaving the fixed root equal to \(1\); this violates \(R_D\).  Fixing \(r=1\) and one true reachable class forces the relation under all completions.
\end{proof}

\begin{remark}[Boundary of the acyclic tree theorem]\label{rem:acyclic-or-boundary}
\Cref{prop:acyclic-or-tree-residual,prop:or-dag-frontier-residual} are residual-completion theorems only for eliminated pure-OR frontier relations.  The DAG theorem removes the tree-only and repeated-frontier caveats for this pure OR fragment after sound quotienting.  It does not say that a proof may ignore local relation audits on displayed internal nodes.  Exposed acyclic non-OR frontiers are handled next by \Cref{prop:acyclic-boolean-frontier-residual}, checker-visible adaptive decision trees by \Cref{prop:visible-adaptive-frontier-residual}, and checker-visible adaptive branching programs by \Cref{prop:visible-adaptive-branching-frontier-residual}; learned clauses, cyclic dependencies, and arbitrary hidden real-SAT semantic influence still require a separate normalizer or realizer.
\end{remark}

\begin{proposition}[Acyclic Boolean frontier residual certificate]\label{prop:acyclic-boolean-frontier-residual}
Let \(C\) be a finite acyclic Boolean circuit with frontier coordinates \(Q\), root coordinate \(r\), and induced function \(F_C:\bits^Q\to\bits\).  After eliminating internal gates, consider the residual frontier relation
\[
  R_C
  :=
  \{(y_r,z)\in\bits^{\{r\}\cup Q}: y_r=F_C(z)\}.
\]
Let \(a\in R_C\), and let \(S\subseteq\{r\}\cup Q\).  Then \(S\) is support-local for \(R_C\) at \(a\) iff \(r\in S\) and \(S\cap Q\) is a standard pointwise certificate forcing \(F_C(z)=a_r\) at the frontier assignment \(a|_Q\).  In particular, the minimum support size is
\[
  1+C_{a_r}(F_C,a|_Q),
\]
where \(C_b(F_C,a|_Q)\) denotes ordinary Boolean certificate complexity at \(a|_Q\) for forcing value \(b\).
\end{proposition}

\begin{proof}
If \(r\notin S\), keep the frontier assignment equal to \(a|_Q\) and flip only \(y_r\).  The completion preserves \(S\) but violates \(y_r=F_C(z)\).  Hence \(r\in S\).

Assume \(r\in S\).  If \(S\cap Q\) is not a certificate forcing \(F_C(z)=a_r\), then there is a frontier assignment \(z'\) agreeing with \(a\) on \(S\cap Q\) but satisfying \(F_C(z')\ne a_r\).  Keeping \(y_r=a_r\) and using \(z'\) gives a completion preserving \(S\) that violates \(R_C\).  Thus support-locality implies that \(S\cap Q\) is a certificate.

Conversely, if \(r\in S\) and \(S\cap Q\) certifies \(F_C(z)=a_r\), then every completion preserving \(S\) keeps \(y_r=a_r\) and keeps \(F_C(z)=a_r\).  Therefore \(y_r=F_C(z)\) holds for every such completion, which is exactly support-locality.  The minimum-size identity follows by minimizing the frontier certificate size and adding the necessary root coordinate.
\end{proof}

\begin{proposition}[Visible adaptive frontier residual certificate]\label{prop:visible-adaptive-frontier-residual}
Let \(T\) be a finite binary decision tree whose internal nodes query frontier coordinates in \(Q\) and whose leaves are labelled by bits.  Let \(r\) be a root/output coordinate and let \(F_T:\bits^Q\to\bits\) be the Boolean function computed by following the queried branch outcomes to a leaf.  Consider the residual frontier relation
\[
  R_T
  :=
  \{(y_r,z)\in\bits^{\{r\}\cup Q}: y_r=F_T(z)\}.
\]
Let \(a\in R_T\), and let \(S\subseteq\{r\}\cup Q\).  Then \(S\) is support-local for \(R_T\) at \(a\) iff \(r\in S\) and \(S\cap Q\) is a standard pointwise certificate forcing \(F_T(z)=a_r\) at the frontier assignment \(a|_Q\).  In particular, the minimum support size is
\[
  1+C_{a_r}(F_T,a|_Q),
\]
and the frontier coordinates queried along the actual \(a|_Q\)-transcript path form one sufficient, though not always minimum, frontier certificate.
\end{proposition}

\begin{proof}
The first three steps are the same root-plus-certificate argument as in \Cref{prop:acyclic-boolean-frontier-residual}.  If \(r\notin S\), keep the frontier assignment equal to \(a|_Q\) and flip \(y_r\), violating \(y_r=F_T(z)\).  If \(r\in S\) but \(S\cap Q\) is not a certificate for \(F_T(z)=a_r\), choose a frontier assignment \(z'\) that agrees with \(a\) on \(S\cap Q\) but has \(F_T(z')\ne a_r\); with \(y_r=a_r\), this preserves \(S\) and violates \(R_T\).  Conversely, if \(r\in S\) and \(S\cap Q\) certifies \(F_T(z)=a_r\), every completion preserving \(S\) keeps both sides of \(y_r=F_T(z)\) equal to \(a_r\).

For the transcript-path claim, fix the values of all coordinates queried along the path followed by \(a|_Q\).  Every completion preserving those values follows the same path to the same leaf and hence has the same value \(F_T(z)=a_r\).  Thus the path variables certify the decision-tree value.  A smaller certificate may exist if different subtrees have the same leaf value, so the path certificate is sufficient rather than necessarily minimum.
\end{proof}

\begin{proposition}[Visible adaptive branching-program frontier residual certificate]\label{prop:visible-adaptive-branching-frontier-residual}
Let \(B\) be a finite acyclic deterministic branching program with start node \(s\), internal nodes labelled by queries to frontier coordinates in \(Q\), two outgoing edges labelled \(0\) and \(1\) from each internal node, and sink nodes labelled by bits.  Let \(r\) be a root/output coordinate and let \(F_B:\bits^Q\to\bits\) be the Boolean function computed by following the queried edge labels from \(s\) to a sink.  Consider the residual frontier relation
\[
  R_B
  :=
  \{(y_r,z)\in\bits^{\{r\}\cup Q}: y_r=F_B(z)\}.
\]
Let \(a\in R_B\), and let \(S\subseteq\{r\}\cup Q\).  Then \(S\) is support-local for \(R_B\) at \(a\) iff \(r\in S\) and \(S\cap Q\) is a standard pointwise certificate forcing \(F_B(z)=a_r\) at the frontier assignment \(a|_Q\).  In particular, the minimum support size is
\[
  1+C_{a_r}(F_B,a|_Q).
\]
The frontier coordinates queried along the actual \(a|_Q\)-computation path form one sufficient, though not always minimum, frontier certificate.
\end{proposition}

\begin{proof}
The root-coordinate and certificate-complexity directions are identical to \Cref{prop:visible-adaptive-frontier-residual}, with \(F_B\) in place of \(F_T\).  If \(r\notin S\), keep the frontier assignment equal to \(a|_Q\) and flip \(y_r\).  This preserves \(S\) but violates \(y_r=F_B(z)\).  If \(r\in S\) and \(S\cap Q\) is not a certificate for \(F_B(z)=a_r\), choose a frontier completion \(z'\) agreeing with \(a\) on \(S\cap Q\) but satisfying \(F_B(z')\ne a_r\).  With \(y_r=a_r\), this completion preserves \(S\) and violates \(R_B\).  Conversely, root plus a frontier certificate for \(F_B(z)=a_r\) forces both sides of the residual equation for every completion preserving \(S\).

For the path-support claim, fix all coordinate values queried along the computation path followed by \(a|_Q\).  Any completion preserving those queried values traverses the same visible branching-program edges to the same sink, hence has the same output \(F_B(z)=a_r\).  Thus the path variables are a sufficient support certificate.  Sharing in the branching program may make the represented function smaller than its unfolded decision tree, and other certificates may be smaller than the actual path, so the path bound is not asserted to be minimum.
\end{proof}

\begin{proposition}[Hidden adaptive transcript insufficiency]\label{prop:hidden-adaptive-transcript-insufficiency}
Fix \(m\ge 2\), frontier coordinates \(Q=\{q_1,\ldots,q_m\}\), root coordinate \(r\), and hard tuple \(a_r=0\), \(a_{q_i}=0\) for all \(i\).  Consider the reported adaptive transcript summary
\[
  \tau=(q_1=0,\ \mathrm{output}=0),
\]
with no certified claim that the transcript is a complete computation path or that the controller has no further hidden dependencies.  Then \(\tau\) is compatible at \(a\) with both residual relations
\[
  R_1=\{(y_r,z):y_r=z_{q_1}\}
  \qquad\text{and}\qquad
  R_{\vee}=\{(y_r,z):y_r=\bigvee_{i=1}^m z_{q_i}\}.
\]
For \(R_1\) at \(a\), the support \(\{r,q_1\}\) is support-local.  For \(R_{\vee}\) at \(a\), every support-local set contains \(r\) and all of \(Q\).  Consequently any support extractor that sees only \(\tau\) and outputs a proper subset of \(\{r\}\cup Q\) is unsound for some hidden adaptive controller compatible with the same reported transcript.  A sound transcript-based normalizer must therefore either certify controller/path completeness, expose enough dependency support, or pay worst-case frontier support.
\end{proposition}

\begin{proof}
Both displayed relations agree with the reported summary at the all-zero tuple: the first queried value is \(q_1=0\) and the output/root value is \(0\).  For \(R_1\), fixing \(r=0\) and \(q_1=0\) forces \(y_r=z_{q_1}\) under every completion preserving those two coordinates, so \(\{r,q_1\}\) is support-local.

For \(R_{\vee}\), omitting \(r\) lets us flip the root while keeping the frontier fixed.  Omitting any frontier coordinate \(q_j\) lets us preserve all coordinates in the alleged support, set \(z_{q_j}=1\), keep all other unforced frontier coordinates arbitrary, and keep \(y_r=0\).  The resulting completion violates \(y_r=\bigvee_i z_{q_i}\).  Thus \(r\) and every coordinate in \(Q\) are necessary, and they are plainly sufficient.

The same reported transcript summary \(\tau\) cannot distinguish these two cases.  Hence a transcript-only extractor returning less than \(\{r\}\cup Q\) is correct for the first relation but unsound for the second.  The conclusion is exactly the hidden-controller boundary: the positive visible-adaptive propositions rely on an audited controller and a certified complete path, not merely on a small reported transcript summary.
\end{proof}

\begin{proposition}[Controller-complete adaptive audit normalizer]\label{prop:controller-complete-adaptive-audit-normalizer}
Let \(H:\bits^Q\to\bits\) be a hidden adaptive controller and let
\[
  R_H=\{(y_r,z)\in\bits^{\{r\}\cup Q}:y_r=H(z)\}.
\]
Fix \(a\in R_H\).  A controller-complete adaptive audit for \(a\) and a support \(S\subseteq\{r\}\cup Q\) consists of a finite acyclic deterministic branching program \(B\) over \(Q\), a replay path \(\pi\) of \(B\) on \(a|_Q\), and a certificate of local controller equivalence such that:
\begin{enumerate}
\item \(r\in S\), and every frontier coordinate queried along \(\pi\) belongs to \(S\cap Q\);
\item the sink reached by \(\pi\) has label \(a_r\);
\item for every \(z'\in\bits^Q\) agreeing with \(a\) on \(S\cap Q\), the hidden controller agrees with the exposed one: \(H(z')=F_B(z')\).
\end{enumerate}
Then \(S\) is support-local for \(R_H\) at \(a\).  If the audit certifies the global identity \(H=F_B\) on all of \(\bits^Q\), then the exact minimum support is again
\[
  1+C_{a_r}(F_B,a|_Q),
\]
as in \Cref{prop:visible-adaptive-branching-frontier-residual}.
\end{proposition}

\begin{proof}
Let \((y_r,z')\) be any completion preserving \(S\).  Since \(r\in S\), the root coordinate remains \(y_r=a_r\).  Since every query along \(\pi\) lies in \(S\cap Q\), the branching program \(B\) sees the same answers along that path on \(z'\) as it sees on \(a|_Q\).  Therefore the exposed controller follows \(\pi\) and outputs \(F_B(z')=a_r\).

The local controller-equivalence certificate transfers this value from the exposed controller to the hidden controller on the entire completion cylinder preserving \(S\cap Q\): \(H(z')=F_B(z')=a_r\).  Hence \(y_r=H(z')\), so the residual relation holds for every completion preserving \(S\).  This is exactly support-locality.

If the audit certifies \(H=F_B\) globally, then \(R_H=R_B\).  The exact minimum support statement is therefore precisely the visible branching-program certificate theorem \(\Cref{prop:visible-adaptive-branching-frontier-residual}\).  The point of the audit is not that a transcript alone determines \(B\); it is that a proof object paying for controller completeness and local equivalence restores the visible-adaptive normalizer that \(\Cref{prop:hidden-adaptive-transcript-insufficiency}\) says is unavailable from the transcript summary alone.
\end{proof}

\begin{corollary}[Trace-complete DPLL log normalizer]\label{cor:trace-complete-dpll}
Consider a DPLL-style decision-log object for a residual relation
\[
  R_H=\{(y_r,z)\in\bits^{\{r\}\cup Q}:y_r=H(z)\},
\]
where the hidden computation may make adaptive queries to frontier coordinates in \(Q\).  Suppose the log is trace-complete at a hard tuple \(a\in R_H\): it contains a finite acyclic DPLL decision controller \(B_\Pi\) over \(Q\), the replay path \(\pi\) followed by \(a|_Q\), terminal output \(a_r\), and clause, propagation, or pruning witnesses certifying that \(H(z')=F_{B_\Pi}(z')\) for every completion \(z'\) agreeing with \(a\) on the root and all frontier coordinates queried along \(\pi\).  Let
\[
  S_\Pi:=\{r\}\cup\{q\in Q:q\text{ is queried along }\pi\}\cup W_\Pi,
\]
where \(W_\Pi\) is the set of formula-label frontier coordinates appearing in those local witness certificates.  Then \(S_\Pi\) is support-local for \(R_H\) at \(a\).  Consequently any proof-size measure or checker budget \(B(|\Pi|)\) that upper-bounds \(|S_\Pi|\) is subject to the same near-threshold formula-footprint lower bound when the DPLL log is used as the proof object for \(\Phi_{n,s,A}\):
\[
  B(|\Pi|)\ge 2n^{a-o(1)}
\]
for \(s=n^{a-o(1)}\) in the polarity-balanced regime of \(\Cref{cor:formula-threshold}\).
\end{corollary}

\begin{proof}
The trace-complete DPLL log is exactly a controller-complete adaptive audit after adding the witness coordinates \(W_\Pi\) to the support: \(B_\Pi\) is the exposed acyclic controller, \(\pi\) is the replay path, and the logged clause, propagation, or pruning witnesses are the local equivalence certificate on the completion cylinder determined by \(S_\Pi\).  \Cref{prop:controller-complete-adaptive-audit-normalizer} therefore makes \(S_\Pi\) support-local for \(R_H\) at \(a\).  When the residual relation is the easy-completion proof object for \(\Phi_{n,s,A}\), the support-local normalization transfer and \(\Cref{cor:formula-threshold}\) give the displayed support-budget lower bound.
\end{proof}

\begin{corollary}[Antecedent-complete CDCL normalizer]\label{cor:antecedent-complete-cdcl}
Consider a CDCL-style learned-clause log for the same residual relation \(R_H\).  Suppose the log is antecedent-complete at the hard tuple \(a\): it contains a finite acyclic decision-and-restart controller \(B_\Pi\), the replay path \(\pi\), terminal output \(a_r\), and, for every learned clause, propagation, or pruning step used along \(\pi\), an acyclic antecedent DAG whose leaves are original support-local clauses, explicit decisions, unit-propagation witnesses, or earlier learned clauses with smaller antecedent rank.  Let \(W_\Pi\) be the set of formula-label frontier coordinates appearing in the antecedent closure of all learned clauses and propagation witnesses used along \(\pi\), and set
\[
  S_\Pi:=\{r\}\cup\{q\in Q:q\text{ is queried along }\pi\}\cup W_\Pi .
\]
Assume the antecedent DAGs certify local equivalence \(H(z')=F_{B_\Pi}(z')\) for every completion \(z'\) agreeing with \(a\) on \(S_\Pi\cap Q\).  Then \(S_\Pi\) is support-local for \(R_H\) at \(a\).  Hence any proof-size measure or checker budget \(B(|\Pi|)\) upper-bounding this replay path plus antecedent-closure support obeys
\[
  B(|\Pi|)\ge 2n^{a-o(1)}
\]
near the polarity-balanced threshold of \(\Cref{cor:formula-threshold}\) when the CDCL log is used as the proof object for \(\Phi_{n,s,A}\).
\end{corollary}

\begin{proof}
Expand the learned-clause uses along their antecedent DAGs.  Acyclicity and antecedent ranks ensure that the expansion terminates in original local clauses, explicit decisions, unit-propagation witnesses, or already expanded lower-rank learned clauses.  By definition, all formula-label frontier coordinates appearing in this expansion are included in \(W_\Pi\), hence in \(S_\Pi\).  The antecedent-completeness certificate therefore supplies the local controller-equivalence field required by \(\Cref{prop:controller-complete-adaptive-audit-normalizer}\): on every completion preserving \(S_\Pi\), the exposed controller and all learned-clause justifications replay with the same values and yield the same terminal output \(a_r\).  Thus \(S_\Pi\) is support-local.  Applying the support-local normalization transfer and the threshold form of the formula-footprint theorem gives the budget lower bound.
\end{proof}

\begin{proposition}[CDCL antecedent-erasure trichotomy]\label{prop:cdcl-antecedent-erasure-trichotomy}
Let a CDCL-erased extractor be a procedure that sees only a finite decision/restart skeleton, the replay path, the sequence of learned-clause names and uses on that path, and the untyped displayed clauses, but does not see learned-clause antecedent DAGs, variable-role provenance, clause-deletion provenance, or an independent support-locality checker.  Suppose the extractor outputs a formula-label support intended to certify completion-local soundness for \(R_H\).  For the unit-plus-blocking skeleton \(C_T\) of \(\Cref{prop:support-erasure-indistinguishability}\), packaged as a single learned-clause use on the replay path, exactly one of the following holds.
\begin{enumerate}[label=(\roman*)]
  \item The extractor obtains additional antecedent provenance, a trusted variable-role map, a namespace certificate, or a system-specific semantic invariant.  Then it is no longer an antecedent-erased CDCL extractor; it has supplied the support audit or learned-clause provenance charged by the positive theorems above.
  \item The extractor is sound for every same-skeleton semantic lift of the erased log.  Then, on \(C_T\), its output must contain all variables in \(T\), hence it pays the role-agnostic support budget for the possible label lift.
  \item The extractor outputs a strict subset of \(T\) while using only the erased CDCL skeleton.  Then the label lift from \(\Cref{prop:support-erasure-indistinguishability}\), with the same erased learned-clause skeleton, falsifies its claimed support soundness.
\end{enumerate}
Consequently, arbitrary CDCL logs with erased learned-clause antecedents do not instantiate \(\Cref{thm:support-local-normalization-transfer}\) by syntax alone.  A positive CDCL theorem must either expose antecedent-complete learned-clause DAGs, add equivalent provenance or namespace data, charge a conservative support budget, or prove a genuinely non-plain semantic invariant that rules out the bad lift.
\end{proposition}

\begin{proof}
Use the auxiliary and label lifts from \(\Cref{prop:support-erasure-indistinguishability}\).  In the auxiliary lift, the displayed learned clause can be realized through auxiliary variables, so formula-label support is empty after roles are known.  In the label lift, the same untyped clause skeleton is a blocking constraint over the variables in \(T\), and \(\Cref{prop:support-audit-necessity}\) forces every sound completion-local support to contain all of \(T\).  Erasing learned-clause antecedents, variable roles, and deletion provenance makes these two lifts identical to the CDCL-erased extractor: it sees the same controller skeleton, replay path, learned-clause name, and untyped displayed clause.  Therefore it must output the same set in both lifts.  If it outputs all of \(T\), it is conservative for this skeleton and pays the role-agnostic budget.  If it outputs a strict subset of \(T\), the label lift witnesses unsoundness.  If it avoids the ambiguity by distinguishing the two lifts, then it used antecedent provenance, role information, namespace information, a support-locality checker, or another semantic invariant beyond the erased CDCL syntax.  These are precisely the three cases.
\end{proof}

\begin{corollary}[RUP-complete CDCL support normalizer]\label{cor:rup-complete-cdcl}
Consider a CDCL/DRUP-style learned-clause log for the residual relation \(R_H\).  Suppose the log is RUP-complete at the hard tuple \(a\): it contains a finite acyclic decision-and-restart controller \(B_\Pi\), the replay path \(\pi\), terminal output \(a_r\), and, for every learned clause, propagation, pruning step, or clause-deletion-dependent use along \(\pi\), a checker-visible RUP propagation certificate over the known live clause database at that moment.  Each RUP certificate records the unit-propagation trace, the clauses touched, the propagated literals, the conflict clause, and the formula-label frontier coordinates touched by those clauses.  The proof object also carries either a trusted formula-label namespace, a trusted variable-role map, or an equivalent local-support parser separating label variables from auxiliary variables; if a deleted clause is used only through a retained learned clause, the retained RUP certificate and its support are still charged.

Let \(W_\Pi\) be the union of all formula-label frontier coordinates appearing in the replayed RUP traces, in the local-support parser output for clauses used by those traces, and in retained certificates for deleted clauses whose consequences are still used.  Set
\[
  S_\Pi:=\{r\}\cup\{q\in Q:q\text{ is queried along }\pi\}\cup W_\Pi .
\]
Assume the RUP traces certify local equivalence \(H(z')=F_{B_\Pi}(z')\) for every completion \(z'\) agreeing with \(a\) on \(S_\Pi\cap Q\).  Then \(S_\Pi\) is support-local for \(R_H\) at \(a\).  Consequently any proof-size measure or checker budget \(B(|\Pi|)\) upper-bounding the replay path together with the parsed RUP-trace support obeys
\[
  B(|\Pi|)\ge 2n^{a-o(1)}
\]
near the polarity-balanced threshold of \(\Cref{cor:formula-threshold}\) when the RUP-complete CDCL log is used as the proof object for \(\Phi_{n,s,A}\).
\end{corollary}

\begin{proof}
A RUP certificate is an executable local witness: starting from the known live clause database, the recorded unit-propagation trace derives the displayed conflict after adding the negation of the learned clause.  Because the certificate lists the clauses touched, propagated literals, conflict clause, and parsed formula-label coordinates, preserving \(S_\Pi\) preserves the truth values of every label coordinate on which that local check depends.  The trusted namespace, role map, or equivalent parser is what makes this a formula-label support statement rather than an untyped syntactic statement.  Retained certificates for deleted clauses ensure that deletion does not erase already used support.

Induct over the accepted learned clauses and propagation steps in the order in which the replay uses them.  The base clauses are original local clauses, explicit decisions, or previously retained audited certificates whose parsed formula-label support lies in \(W_\Pi\).  The induction step is the RUP trace itself: every clause and literal read by the unit-propagation check has support inside \(W_\Pi\), so the same local implication is valid on every completion preserving \(S_\Pi\).  Thus the replayed controller, all learned clauses used by it, and the terminal output remain locally equivalent to the hidden residual computation on the completion cylinder determined by \(S_\Pi\).

This is exactly the local controller-equivalence field required by \(\Cref{prop:controller-complete-adaptive-audit-normalizer}\).  Therefore \(S_\Pi\) is support-local for \(R_H\) at \(a\).  Applying the support-local normalization transfer and the formula-footprint threshold lower bound gives the displayed budget bound.  The statement is deliberately non-plain: without the live-database record, parsed formula-label namespace or role map, and retained support for deleted-clause consequences, the erasure ambiguity of \(\Cref{prop:cdcl-antecedent-erasure-trichotomy}\) remains.
\end{proof}

\begin{corollary}[Canonical-namespace RUP-CDCL standardizer]\label{cor:canonical-namespace-rup-cdcl}
In the setting of \(\Cref{cor:rup-complete-cdcl}\), suppose the CDCL/DRUP-style trace is accompanied by a deterministic encoding manifest \(\mathcal M\) with the following checker-visible fields.
\begin{enumerate}[label=(\roman*)]
  \item A reserved formula-label variable interval \(I_{\mathrm{lab}}\) and an injective map \(\Lambda:I_{\mathrm{lab}}\to Q\) from those variables to frontier coordinates.
  \item A disjoint auxiliary-variable region \(I_{\mathrm{aux}}\), with every displayed variable belonging to exactly one of \(I_{\mathrm{lab}}\) or \(I_{\mathrm{aux}}\).
  \item A rule saying that the formula-label support of any displayed clause \(C\) is exactly \(\Lambda(C\cap I_{\mathrm{lab}})\), ignoring literal signs and auxiliary variables.
  \item A live-database log sufficient for the RUP checker to reconstruct, at each learned-clause check, the currently live original clauses, accepted learned clauses, and deletions.
  \item A retention rule saying that when a deleted clause contributes through a retained learned clause, the already checked RUP certificate for that retained learned clause remains available and its parsed label support remains charged.
\end{enumerate}
Assume every RUP propagation certificate in the trace is checker-visible over the live database reconstructed from \(\mathcal M\).  Then the manifest induces an equivalent local-support parser, so the trace satisfies the parsing hypothesis of \(\Cref{cor:rup-complete-cdcl}\) without an additional trusted role map.  Hence the canonical support
\[
  S_\Pi^{\mathcal M}
  :=
  \{r\}\cup\{q\in Q:q\text{ is queried along }\pi\}
  \cup
  \bigcup_{C\text{ touched by replayed RUP traces}}\Lambda(C\cap I_{\mathrm{lab}})
\]
together with retained deleted-clause certificate supports is support-local for \(R_H\) at \(a\).  Any checker budget upper-bounding this canonical parsed support obeys the same near-threshold formula-footprint lower bound.
\end{corollary}

\begin{proof}
The manifest makes variable roles syntactic rather than trusted semantic advice.  The checker verifies disjointness of \(I_{\mathrm{lab}}\) and \(I_{\mathrm{aux}}\), verifies that every displayed variable lies in exactly one region, and computes the label support of each touched clause by the declared interval intersection \(C\cap I_{\mathrm{lab}}\) followed by \(\Lambda\).  This computation is deterministic and local to the displayed clause and manifest; no hidden role judgment is used.

The live-database log reconstructs the clause database used by each RUP check, while the retention rule prevents deletion from erasing support already charged to retained consequences.  Therefore every replayed RUP certificate has exactly the data required by \(\Cref{cor:rup-complete-cdcl}\): touched clauses, propagated literals, conflict clause, live database, and a parsed formula-label support for all touched clauses.  Applying \(\Cref{cor:rup-complete-cdcl}\) gives support-locality for \(S_\Pi^{\mathcal M}\) and then the formula-footprint budget bound.

This corollary standardizes the namespace route; it does not make unmanifested DIMACS, raw RUP, or raw DRUP syntax into a support normalizer.  If the interval map, live database, or retention rule is erased, the extractor is back in the ambiguity regime of \(\Cref{prop:cdcl-antecedent-erasure-trichotomy}\).
\end{proof}

\begin{remark}[Boundary of exposed Boolean frontiers]\label{rem:acyclic-boolean-frontier-boundary}
\Cref{prop:acyclic-or-tree-residual,prop:or-dag-frontier-residual} are the pure-OR special cases of \Cref{prop:acyclic-boolean-frontier-residual}.  The Boolean-frontier theorem allows exposed finite acyclic AND/OR/NOT frontier circuits, \Cref{prop:visible-adaptive-frontier-residual} records the checker-visible adaptive decision-tree case, and \Cref{prop:visible-adaptive-branching-frontier-residual} records the shared-state acyclic branching-program case.  \Cref{prop:hidden-adaptive-transcript-insufficiency} shows why merely reporting a small adaptive transcript is not a substitute for that visibility, while \Cref{prop:controller-complete-adaptive-audit-normalizer} records the positive audit that is sufficient once controller completeness and local equivalence are paid for.  \Cref{cor:trace-complete-dpll} packages the first recognizable proof-object instantiation of that audit: a DPLL decision log is admissible only when it is trace-complete and carries the local witnesses that make the replayed controller sound on the claimed completion cylinder.  \Cref{cor:antecedent-complete-cdcl} adds the learned-clause version: CDCL-style compression is admissible when the used learned clauses carry acyclic antecedent DAGs whose leaves expose the local support they depend on.  \Cref{prop:cdcl-antecedent-erasure-trichotomy} records the matching boundary for erased learned-clause logs: add provenance, pay conservative support, or become unsound on a same-skeleton lift.  \Cref{cor:rup-complete-cdcl} gives the positive non-plain response: a CDCL/DRUP-style log can supply the same audit through checker-visible RUP propagation certificates, a known live clause database, trusted label namespace or role parsing, and retained deleted-clause support.  \Cref{cor:canonical-namespace-rup-cdcl} standardizes one such parsing route: a deterministic encoding manifest with a reserved label interval lets the checker compute support directly from touched clauses.  The positive results still assume that the semantic influence has been exposed as an auditable finite object.  They do not handle raw RUP/DRUP traces without a manifest or support parsing, unaudited learned clauses, hidden clause minimization, cyclic restart/search graphs, hidden proof-system influence, or an arbitrary SAT self-reduction proof whose frontier circuit, adaptive control, and quotienting have not been made checker-visible.  Those cases still require a separate normalizer or realizer.
\end{remark}

\section{Application to polylogarithmic SAT parity carriers}\label{sec:sat}

We now specialize the labels to satisfiability.  Fix a standard total encoding of Boolean formulas: malformed strings are assigned SAT value $0$.  Let
\[
  \ell_n=\lceil(\log n)^d\rceil,
  \qquad d>1,
\]
and let the batched generator have
\[
  M_n=n^{q+o(1)},
  \qquad
  g_n=n^{\rho+o(1)}.
\]
Define
\begin{equation}\label{eq:sat-carrier}
  H_n^{\SAT}(x)
  :=\bigoplus_{j=1}^{M_n}\SAT_{\ell_n}(G_{n,j}(x)).
\end{equation}
The condition $d>1$ is not needed for substitution.  It merely prevents the complete $\ell_n$-bit SAT table, of size $2^{\ell_n}$, from being polynomial in $n$.

\subsection{Exponent debt and the all-exponent regime}

The following accounting proposition separates the general issue from this carrier.

\begin{proposition}[Exponent-window lemma]\label{prop:window}
Suppose that for every fixed $c\ge1$,
\[
  \SAT\in\SIZE(m^c)
  \quad\Longrightarrow\quad
  \Ckt(K_n)\le n^{q+\beta c+o(1)}.
\]
If $\Ckt(K_n)>n^{a-o(1)}$ infinitely often, then
\[
  \SAT\notin\SIZE(m^c)
\]
for every fixed $c$ satisfying $q+\beta c<a$.  Hence one fixed source exponent excludes all polynomial target exponents only when $\beta=0$ and $a>q$.
\end{proposition}

\begin{proof}
For a fixed $c$ with $q+\beta c<a$, choose a positive constant gap between the two exponents.  The conditional upper bound then contradicts the infinitely-often lower bound on all sufficiently large hard lengths.
\end{proof}

Polylogarithmic query length has $\beta=0$.

\begin{lemma}[Batched substitution]\label{lem:substitution}
If a label function $y:\bits^{\ell_n}\to\bits$ has circuit size $t_n$, then for every size-$O(M_n)$ aggregator $A$,
\[
  \Ckt(H_{n,A}^{y})
  \le g_n+M_nt_n+O(M_n).
\]
\end{lemma}

\begin{proof}
Compute the whole query batch once, feed every query into a separate copy of the size-$t_n$ label circuit, and apply the aggregator.
\end{proof}

\begin{theorem}[Polylogarithmic-query transfer]\label{thm:beta-zero}
For every fixed $c\ge1$,
\begin{equation}\label{eq:beta-zero}
  \SAT\in\SIZE(m^c)
  \quad\Longrightarrow\quad
  \Ckt(H_n^{\SAT})
  \le n^{\max\{q,\rho\}+o(1)}.
\end{equation}
Consequently, if for some $a>\max\{q,\rho\}$,
\[
  \Ckt(H_n^{\SAT})>n^{a-o(1)}
\]
holds infinitely often, then $\SAT\notin\Ppoly$.
\end{theorem}

\begin{proof}
Under the target assumption, $\SAT_{\ell_n}$ has circuits of size $\ell_n^c$.  Since $\ell_n^c=(\log n)^{dc+o(1)}=n^{o(1)}$, \Cref{lem:substitution} gives
\[
  \Ckt(H_n^{\SAT})
  \le n^{\rho+o(1)}+n^{q+o(1)}
  =n^{\max\{q,\rho\}+o(1)}.
\]
The final claim follows from \Cref{prop:window}, or directly by fixing the polynomial exponent of a hypothetical SAT circuit family.
\end{proof}

\subsection{Reverse calibration}

The same substitution inequality can be read backward.

\begin{proposition}[Reverse substitution burden]\label{prop:reverse}
If $\Ckt(H_n^{\SAT})>s_n$, then
\begin{equation}\label{eq:reverse}
  \Ckt(\SAT_{\ell_n})
  >\frac{s_n-g_n-O(M_n)}{M_n}.
\end{equation}
The conclusion is vacuous when the numerator is nonpositive.
\end{proposition}

\begin{proof}
Apply \Cref{lem:substitution} with an optimal circuit for $\SAT_{\ell_n}$ and rearrange.
\end{proof}

At $s_n=n^{1+\eps-o(1)}$ with $q,\rho<1+\eps$, \eqref{eq:reverse} yields
\[
  \Ckt(\SAT_{\ell_n})
  >n^{1+\eps-q-o(1)}
  =2^{\Omega(\ell_n^{1/d})}
\]
on the corresponding hard lengths.  Thus the transfer removes dependence on the unknown polynomial exponent, but a direct source lower bound already contains a stretched-exponential lower bound for the matched short SAT slices.

\subsection{Certificate consequences}

Let
\[
  \lambda_n^{\SAT}:=\SAT_{\ell_n}|_{U_n}
\]
be the actual formula labeling, and let
\[
  z_n^{\SAT}(x,j):=\SAT_{\ell_n}(G_{n,j}(x))
\]
be its induced occurrence labeling.

\begin{corollary}[Local certificate costs for a hard SAT carrier]\label{cor:sat-costs}
Assume
\[
  \Ckt(H_n^{\SAT})>s_n.
\]
Then the following hold.

\begin{enumerate}
  \item Every $1$-certificate for $\HardF_{n,s_n,\oplus}$ at $\lambda_n^{\SAT}$ satisfies the two inequalities in \Cref{thm:formula-footprint}; in the near-threshold regime it exposes $n^{1+\eps-o(1)}$ distinct satisfiable generated formulas and the same number of unsatisfiable generated formulas.

  \item Every $1$-certificate for $\HardO_{n,s_n,\oplus}$ at $z_n^{\SAT}$ queries $\Omega(M_ns_n/n)$ occurrences and contains $\Omega(s_n/n)$ complete rows of each parity.

  \item Let $E_n$ be any equivalence relation on $U_n$ on whose classes SAT is constant.  Every $1$-certificate for the resulting quotient hardness predicate has residual constant distance $\Omega(s_n/n)$.
\end{enumerate}
\end{corollary}

The three statements have different hypotheses.  A formula-local proof need not be sound under independent occurrence completions, and a proof using SAT self-reduction need not be sound over all equality-quotient completions.  The corollary should therefore be read as a menu of barriers for explicitly specified certificate interfaces, not as an exhaustive theorem about arbitrary SAT-semantic proofs.

\subsection{A brief global naturalization observation}

The local certificate results do not require largeness.  At the opposite extreme, a dense truth-table property that is functionally sound for the carrier enters the natural-proofs regime.

\begin{proposition}[Functional naturalization]\label{prop:naturalization}
Let $\mathcal P_n$ be a property of full truth tables $y:\bits^{\ell_n}\to\bits$ such that
\[
  y\in\mathcal P_n
  \quad\Longrightarrow\quad
  \Ckt(H_{n,\oplus}^{y|_{U_n}})>n^{a-o(1)}
\]
for some $a>\max\{q,\rho\}$.  Then, for every fixed $c$ and all sufficiently large relevant $n$,
\[
  y\in\mathcal P_n
  \quad\Longrightarrow\quad
  \Ckt(y)>\ell_n^c.
\]
If $\mathcal P_n$ is decidable in time $2^{O(\ell_n)}$ from the truth table and has density at least $2^{-O(\ell_n)}$, it is a constructive large property useful against polynomial-size circuits on the matched lengths.
\end{proposition}

\begin{proof}
If $y$ had a circuit of size $\ell_n^c$, \Cref{lem:substitution} would give
\[
  \Ckt(H_{n,\oplus}^{y|_{U_n}})
  \le n^{\max\{q,\rho\}+o(1)},
\]
contradicting functional soundness.
\end{proof}

Under the usual pseudorandom-function assumptions, the final combination of constructivity, largeness, and usefulness is subject to the Razborov--Rudich barrier~\cite{RR97}.  The point is not that every certificate naturalizes: sparse, SAT-specific, or merely local certificates can fall outside these hypotheses.

\subsection{What the application does not establish}

No explicit generator is presently shown here to make $H_n^{\SAT}$ require $n^{1+\eps}$ unrestricted gates.  Such a theorem would already imply $\SAT\notin\Ppoly$ by \Cref{thm:beta-zero} and would contain the short-length lower bound in \Cref{prop:reverse}.  Likewise, the certificate theorems do not automatically yield lower bounds for Frege, Extended Frege, IPS, implicit proofs, or another standard proof system.  \Cref{thm:displayed-atom-extraction} supplies this bridge only for restricted displayed-atom interfaces.  A standard proof-system application would still require a normalization theorem showing that short proofs induce one of the completion-sound certificate models analyzed above with controlled overhead.

The randomized theorem in \Cref{thm:accepting-union} gives a bounded-error lower bound for the union of labels queried on accepting executions.  It should not be confused with an expected-query lower bound.  A verifier could, in principle, query few labels on each random string while spreading its accepting paths over a large union.  Controlling expected footprint, transcript support, or information revealed about the labeling would require a minimax, distributional, or information-theoretic argument beyond the deterministic certificate reduction.

\section{Proof dependency and audit map}\label{sec:audit}

\Cref{tab:proof-dependencies} records the logical dependency structure of the theorem package.  Its purpose is to make clear which claims use only elementary completion arguments, which claims import circuit-counting or sorting-network facts, and which advertised motivations remain outside the proved dependency graph.

{\scriptsize
\begin{longtable}{@{}p{0.22\textwidth}p{0.32\textwidth}p{0.40\textwidth}@{}}
\caption{Proof dependency and audit map.  Every theorem-level claim is traced to its local dependencies and its explicit nonclaim boundary.}
\label{tab:proof-dependencies}\\
\toprule
Claim & Dependencies & Audit boundary \\
\midrule
\endfirsthead
\toprule
Claim & Dependencies & Audit boundary \\
\midrule
\endhead
Equality patching, \Cref{lem:patching} & Equality tests for explicit rows and XOR patching. & Nonuniform upper bound; no lower-bound assumption. \\
Residual-error threshold, \Cref{cor:residual} & \Cref{lem:patching}. & Converts source hardness into distance from any subthreshold circuit. \\
Batched static dictionary, \Cref{lem:dictionary} & Sorting networks with tagged query and dictionary records. & Nonuniform dictionary; dictionary records are masked out before aggregation. \\
Sparse formula-defect upper bound, \Cref{prop:defect} & Batched generator plus \Cref{lem:dictionary}. & Applies to formula-label completions, with repeated generated strings sharing one key. \\
Formula-footprint barrier, \Cref{thm:formula-footprint} & Certificate completion to a constant baseline plus \Cref{prop:defect}. & Standard nondeterministic certificate lower bound; not a proof-system lower bound by itself. \\
Decision-tree lower bound and near-threshold corollaries, \Cref{cor:formula-ndt,cor:formula-threshold,cor:diversity} & \Cref{thm:formula-footprint} and exponent accounting. & Gives standard certificate-complexity payoff; no SAT lower bound. \\
Identity-carrier generic sharpness, \Cref{thm:identity-generic-sharpness} & Circuit counting, union bound, and Chernoff bound. & Generic truth-table sharpness; not an explicit SAT carrier. \\
Deterministic and bounded-error formula verifiers, \Cref{prop:adaptive-formula,thm:accepting-union,cor:randomized-formula-union} & Transcript/certificate equivalence plus \Cref{thm:formula-footprint}. & Bounded-error result is for accepting-query union, not expected query complexity. \\
Displayed-atom extraction, \Cref{thm:displayed-atom-extraction} & Definition of faithful completion-sound displayed atoms plus \Cref{thm:formula-footprint}. & Restricted proof interface only; no Frege, Extended Frege, IPS, or implicit-proof lower bound. \\
Resolution support-local bridge, \Cref{thm:assumption-resolution,cor:resolution-unit-core,thm:resolution-local-clause-support,thm:support-local-aux-resolution,thm:support-local-extension-resolution} & Resolution soundness, easy-completion CNF semantics, unused-leaf deletion, local-clause support, support-local auxiliary witnesses, support-local extension definitions, and \Cref{thm:formula-footprint}. & Named resolution statements only when faithful label information is explicit, auxiliary constraints are support-local, or extension variables are computed from a fixed label support. \\
Dependency-cone and canonical syntactic-cone extension normalizers, final backward-slice, slice-explicit ER, displayed acyclic ER, expansion-bounded compressed ER, macro-compressed ER, support-declared macro ER, support-audited LRAT, provenance-typed LRAT, namespace-typed LRAT, strict-prefix Qrep replay, support-audit necessity, support-erasure indistinguishability, plain LRAT extraction trichotomy, cone-budget, and bounded-depth lower bounds, \Cref{def:dependency-cone-extension,thm:dependency-cone-extension-normalizer,cor:canonical-syntactic-cone,cor:final-backward-slice,def:slice-explicit-er,cor:slice-explicit-er,cor:displayed-er-length,cor:compressed-er-expansion,cor:macro-compressed-er,cor:support-declared-macro-er,cor:support-audited-lrat,cor:provenance-typed-lrat,cor:namespace-typed-lrat,prop:strict-prefix-qrep,prop:support-audit-necessity,prop:support-erasure-indistinguishability,prop:plain-lrat-extraction-trichotomy,cor:cone-budget-extension,cor:bounded-depth-extension} & Acyclic extension evaluation, syntactic cone monotonicity, canonical graph reachability, backward-slice identity, displayed-literal charging, expansion accounting, macro-reuse accounting, support-declaration soundness, LRAT soundness, typed variable-role provenance, namespace parsing, LRAT antecedent closure, block-local witness composition, strict-prefix quotient replay, support-locality, the unit-clause separator for unaudited support, provenance-erasure ambiguity, the raw-trace extractor trichotomy, cone counting, depth induction, and \Cref{thm:formula-footprint}. & Extension and clausal traces only when final cone or slice support is visible in the displayed acyclic extension object, when the compressed representation expands to such an object with controlled overhead or visible macro reuse, when a sound declared/audited support budget is supplied, when typed block provenance or a trusted label-variable namespace plus antecedent closure composes into such an audit, or when quotient representatives replay from already charged same-theory forests; unrestricted extension variables, delayed quotient certificates, hidden semantic pruning, and raw unaudited or untyped LRAT/DRAT traces may still have huge, compressed, nonlocal, or semantically ambiguous dependency support, and checked clausal unsatisfiability plus raw trace syntax alone does not certify completion-local support. \\
Support-local normalization transfer, \Cref{def:support-local-normalizer,thm:support-local-normalization-transfer} & Sound refutation, support-locality, easy-completion semantics, and \Cref{thm:formula-footprint}. & Conditional transfer theorem; a concrete proof system still needs a normalizer. \\
Row-forcing theorem and occurrence consequences, \Cref{thm:row-forcing,cor:full-row,prop:adaptive-occ} & Row completion to a constant and \Cref{lem:patching}; aggregator forcing numbers. & Private occurrence model; equality or nonlinear SAT constraints require separate treatment. \\
Quotient residual-distance barrier, \Cref{thm:quotient-distance} & Affine quotient completion plus \Cref{lem:patching}. & Equality constraints only; nonlinear semantic constraints need separate support accounting. \\
Semantic-relation support certificate, \Cref{prop:semantic-relation-support} & Definition of support-locality for a finite relation; enumeration of \(y=y_0\vee y_1\). & Local relation-instance audit only; the OR-frontier theorems below are bounded eliminated-frontier cases, not full SAT self-reduction closure theorems. \\
Acyclic OR-tree residual certificate, \Cref{prop:acyclic-or-tree-residual} & Elimination of internal OR nodes to the frontier relation \(y_r=\bigvee_{\ell\in L(T)}y_\ell\), plus ordinary OR certificate complexity. & Bounded eliminated-frontier theorem only; shared DAGs and frontier quotienting are handled separately by \Cref{prop:or-dag-frontier-residual}. \\
OR-DAG frontier residual certificate, \Cref{prop:or-dag-frontier-residual} & Topological elimination of finite acyclic OR DAGs, idempotence of OR, and sound quotienting of repeated reachable frontier formulas. & Pure-OR shared-frontier theorem; the exposed acyclic Boolean generalization is handled by \Cref{prop:acyclic-boolean-frontier-residual}. \\
Acyclic Boolean frontier residual certificate, \Cref{prop:acyclic-boolean-frontier-residual} & Exposed finite acyclic Boolean frontier circuit and ordinary pointwise certificate complexity of the induced frontier function. & Exposed-frontier theorem only; visible adaptive control is handled by \Cref{prop:visible-adaptive-frontier-residual}, while learned clauses, cyclic semantic systems, and hidden semantic proof influence remain outside scope. \\
Visible adaptive frontier residual certificate, \Cref{prop:visible-adaptive-frontier-residual} & Exposed finite adaptive decision tree over frontier coordinates and ordinary pointwise certificate complexity of the induced decision-tree function. & Visible-adaptive theorem only; the whole control tree must be checker-visible, and hidden adaptive pruning or learned semantic influence remains outside scope. \\
Visible adaptive branching-program frontier residual certificate, \Cref{prop:visible-adaptive-branching-frontier-residual} & Exposed finite acyclic deterministic branching program over frontier coordinates and ordinary pointwise certificate complexity of the induced branching-program function. & Visible shared-adaptive theorem only; the whole branching program must be checker-visible, and hidden adaptive controllers or learned semantic influence remain outside scope. \\
Hidden adaptive transcript insufficiency, \Cref{prop:hidden-adaptive-transcript-insufficiency} & Same reported transcript summary is compatible with a one-variable function and an all-frontier OR function at the all-zero tuple. & Boundary theorem only; it proves transcript-only extraction is unsound without controller/path completeness, exposed dependencies, or worst-case support. \\
Controller-complete adaptive audit normalizer, \Cref{prop:controller-complete-adaptive-audit-normalizer} & Exposed acyclic controller, replay path, and local equivalence certificate on the claimed completion cylinder, plus \Cref{prop:visible-adaptive-branching-frontier-residual} for the global case. & Positive audited-hidden-adaptive theorem only; it assumes the controller-completeness audit and does not show that arbitrary hidden SAT self-reduction semantics or raw proof traces supply one. \\
Trace-complete DPLL decision-log normalizer, \Cref{cor:trace-complete-dpll} & Complete finite DPLL decision controller, replay path, terminal output, and clause or propagation witnesses certifying local equivalence on the claimed completion cylinder, plus \Cref{prop:controller-complete-adaptive-audit-normalizer}. & Recognizable DPLL-style bridge only for trace-complete logs; partial solver transcripts, unaudited learned clauses, restarts, cyclic search graphs, and arbitrary CDCL traces still need separate provenance or controller-completeness certificates. \\
Antecedent-complete CDCL learned-clause normalizer, \Cref{cor:antecedent-complete-cdcl} & Complete decision/restart controller, replay path, learned-clause antecedent DAGs, rank-decreasing antecedent closure, and local equivalence certificates, plus \Cref{cor:trace-complete-dpll,prop:controller-complete-adaptive-audit-normalizer}. & Recognizable CDCL-style bridge only when every used learned clause has acyclic antecedent provenance down to original local witnesses; arbitrary CDCL logs, hidden clause minimization, deleted-clause provenance gaps, and cyclic restart/search graphs remain outside scope. \\
CDCL antecedent-erasure trichotomy, \Cref{prop:cdcl-antecedent-erasure-trichotomy} & Support-erasure indistinguishability, support-audit necessity, and the same-skeleton learned-clause lift of the unit-plus-blocking separator. & Boundary result for arbitrary provenance-free CDCL logs: add learned-clause antecedent provenance or an equivalent invariant, pay role-agnostic support, or lose support soundness on an indistinguishable semantic lift. \\
RUP-complete CDCL support normalizer, \Cref{cor:rup-complete-cdcl} & Checker-visible RUP propagation certificates over a known live clause database, trusted label namespace or role parsing, retained deleted-clause support, and \Cref{prop:controller-complete-adaptive-audit-normalizer}. & Positive non-plain CDCL bridge only for logs whose RUP traces expose parsed formula-label support; raw RUP/DRUP syntax, erased role information, and uncharged deleted-clause influence remain outside scope. \\
Canonical-namespace RUP-CDCL standardizer, \Cref{cor:canonical-namespace-rup-cdcl} & Deterministic encoding manifest, reserved label-variable interval, injective coordinate map, live-database reconstruction, retained deleted-clause certificate support, and \Cref{cor:rup-complete-cdcl}. & Checkable namespace route only; unmanifested DIMACS/RUP/DRUP syntax, erased interval maps, and hidden variable roles remain outside scope. \\
SAT transfer and reverse calibration, \Cref{prop:window,lem:substitution,thm:beta-zero,prop:reverse,cor:sat-costs,prop:naturalization} & Batched substitution and the certificate theorems above. & Conditional on a source lower bound; does not prove \(\SAT\notin\Ppoly\). \\
\bottomrule
\end{longtable}
}

\section{Related work and open directions}\label{sec:related}

\subsection{Hardness magnification and locality}

Hardness magnification was developed systematically by Oliveira and Santhanam~\cite{OS18}, with further near-frontier results by Oliveira, Pich, and Santhanam~\cite{OPS21} and compression-based magnification by McKay, Murray, and Williams~\cite{MMW19}.  Chen, Hirahara, Oliveira, Pich, Rajgopal, and Santhanam studied the interaction between magnification, natural proofs, and locality~\cite{CHOPRS22}.  Recent work of Atserias and M\"uller develops general magnification through sparse distinguisher matrices and obtains sharp thresholds in several settings~\cite{AM25}.

The present paper does not offer a new general magnification theorem.  Its SAT substitution calculation is an application of the elementary fact that polylogarithmic target length removes the target exponent from the source exponent.  The new focus is the local certificate complexity of the resulting source-hardness predicate.  Formula footprints concern the number of distinct generated labels that must be fixed; row forcing concerns the number and type of aggregate rows that must be made immutable; quotient distance concerns the geometry left after sound equality identifications.

\subsection{Certificate and decision-tree complexity}

Certificate complexity is the nondeterministic analogue of deterministic query complexity and is a standard measure of Boolean functions~\cite{Nisan91,BuhrmanDeWolf02}.  Our hardness predicates have extremely large and highly structured variable sets: distinct generator outputs, generated occurrences, or quotient classes.  The results are inputwise lower bounds for $1$-certificates of these properties.

The row-forcing theorem is not merely the observation that parity has certificate complexity $M$.  Sparse-row patching first forces $\Omega(s/n)$ rows of each output value, after which the minimum forcing number of the aggregator converts rows into label probes.  The formula theorem is different: repeated occurrences share variables, and a batched dictionary upper bound converts a small two-sided formula footprint into a small source circuit.

\subsection{Circuit test suites, anti-checkers, and exact specification}

There is a close conceptual relation between our certificates and finite sets of examples that refute small circuits.  Chapman and Williams study the circuit-input game and the data complexity of testing circuits with labeled examples, showing a duality between small test suites and circuit lower bounds~\cite{ChapmanWilliams15}.  That line of work asks for labeled source inputs that distinguish a fixed target function from all small circuits.  Our formula-label theorem instead fixes labels on generated objects and asks whether every completion of the ambient label table still induces a hard source function.  Thus the roles of examples are different: test suites constrain candidate circuits for a fixed function, whereas formula footprints constrain completions of a generator-induced labeling.

\Cref{thm:identity-generic-sharpness} deliberately uses the standard counting side of this picture as a sharpness check.  In the identity carrier, generated labels are source-input labels, so the formula-footprint certificate becomes an ordinary finite test set against all size-\(s\) circuits.  The theorem does not replace the main lower bound; it shows that the lower bound has the right scale in the baseline model.

The same distinction separates the present work from teaching dimension and exact specification.  Exact specification by examples studies how many labeled examples uniquely specify a target concept inside a hypothesis class~\cite{AnthonyBCS92}; teaching-dimension work studies the information complexity of selected examples for learning or teaching concept classes~\cite{DoliwaFSZ14}.  Our certificates need not identify the whole labeling or teach a unique concept.  They only force membership in the hardness predicate: every completion must induce a function outside the size-\(s\) circuit ball.  This is why the lower bounds count polarity footprints, forced rows, and quotient residual distance rather than uniqueness of the target labeling.

\subsection{Partial MCSP and partial functions}

The formula-label predicate is also near partial-function variants of MCSP.  MCSP asks whether a complete truth table has a circuit of size at most \(s\), while partial MCSP variants ask whether a partial function has a small circuit completion.  Hirahara proves NP-hardness for a partial-function variant of MCSP under randomized reductions and places it in the meta-complexity landscape~\cite{Hirahara22PartialMCSP}.  The present predicate is the complementary, certificate-complexity viewpoint for a generator-induced partial labeling: a certificate for hardness must rule out every small-circuit completion after the unexposed generated labels are filled in.  The batched dictionary lemma is the point at which the generator structure matters; it turns a small set of formula-label exceptions into a small source circuit, yielding the near-threshold two-polarity footprint lower bound.

\begin{table}[htbp]
\centering
\scriptsize
\begin{tabularx}{\textwidth}{@{}p{0.25\textwidth}p{0.29\textwidth}X@{}}
\toprule
Nearby line & Shared feature & Distinction in this paper \\
\midrule
Partial MCSP and partial functions & Completion of a partial object must avoid or admit small circuits. & The partial object is induced by a generator; the lower bound is an inputwise certificate lower bound for forcing hardness under all completions. \\
Circuit test suites and anti-checkers & Labeled examples can refute all small circuits for a target. & Formula footprints fix generated labels, not source inputs, and the theorem exploits shared batched dictionary lookup. \\
Teaching and exact specification & Examples force a concept inside a hypothesis class. & The certificate need not identify a unique labeling; it only forces membership in the hardness predicate. \\
Decision-tree certificate complexity & A partial assignment forces a Boolean property. & The variables have formula, occurrence, or quotient geometry, yielding polarity, row, and affine-distance refinements. \\
Hardness magnification & A source lower bound could imply a stronger target lower bound. & The SAT section is conditional calibration; the paper does not prove the source lower bound. \\
Proof complexity & Proofs may expose local semantic information. & Displayed acyclic ER objects inherit a length lower bound here; expansion-bounded compressed ER objects inherit the corresponding decompression tradeoff; macro-compressed ER objects inherit a reuse-product tradeoff; support-declared macro objects, support-audited LRAT traces, provenance-typed LRAT/DRAT-with-LRAT-certificate traces, and namespace-typed LRAT antecedent-closure traces inherit support-budget lower bounds; strict-prefix Qrep objects inherit a replay condition for quotient representatives; trace-complete DPLL decision logs inherit the controller-complete adaptive audit; antecedent-complete CDCL logs inherit the same audit after expanding learned clauses through acyclic antecedent DAGs; self-auditing RUP-CDCL traces instantiate the carrier-independent audit criterion when the decoder and support audit are internal to the proof object, append-only self-auditing traces remove retained deleted-clause support records by forbidding deletions, and prefix-typed append-only traces remove the proof-internal decoder by making formula-label variables public syntax. The support-audit necessity, support-erasure, hidden-transcript, and plain LRAT trichotomy results show why raw unaudited, untyped, partial-transcript, or antecedent-erased support is not enough. Frege, Extended Frege, IPS, arbitrary compressed proofs, arbitrary CDCL without antecedent provenance or typed audit fields, plain unaudited LRAT/DRAT without extra semantic invariants, delayed quotient certificates, and unrestricted nonlocal ER still require a separate extraction theorem. \\
\bottomrule
\end{tabularx}
\caption{Nearest-work comparison and novelty boundary.}
\label{tab:related-boundary}
\end{table}

\subsection{Novelty accounting and narrative cohesion}

The comparison above is also a cohesion test.  Each theorem-level contribution
must identify the nearest established object, the shared mathematical feature,
the novelty delta, and its role in the formula-footprint arc.  This prevents the
paper from becoming a list of locally correct but separately motivated results:
an admitted result must either be the flagship obstruction, calibrate its
sharpness, expose the same obstruction in a new interface, bridge it to proof
objects, or mark a boundary where the bridge is unsound.

\begin{table}[htbp]
\centering
\small
\begin{tabularx}{\textwidth}{@{}p{0.18\textwidth}p{0.24\textwidth}p{0.26\textwidth}X@{}}
\toprule
Object & Closest literature object & Novelty delta & Role in the flagship arc \\
\midrule
Formula-footprint barrier & Certificate complexity, Partial MCSP, circuit test suites & Completion-sound certificates for generator-induced labels must expose many distinct generated formulas of both polarities; the batched dictionary lemma gives the near-threshold count. & Flagship: the paper's main before/after shift is from examples-as-source-inputs to formula labels as a structured completion footprint. \\
Identity-carrier sharpness & Counting lower bounds against small circuits & In the standard truth-table carrier, generic certificates match the lower bound up to polylogarithmic factors. & Calibration: the flagship lower bound is not merely an artifact of the proof. \\
Row and quotient interfaces & Decision-tree certificates and equality-constrained partial functions & Source rows, output polarities, and quotient residual distance are separated as different interfaces of the same completion obstruction. & Interface variant: these results explain what changes when the local object is not a generated formula label. \\
Semantic frontiers and adaptive certificates & Boolean certificate complexity, decision trees, branching programs, SAT self-reduction frontiers & Exposed finite frontiers reduce to ordinary pointwise certificate complexity, while hidden transcript summaries are separated from sound support audits. & Boundary: semantic constraints help only when the support object is checker-visible. \\
Proof-object bridges & Resolution, ER-style syntax, clausal traces, and implicit proofs & Support-local normalizers, provenance or namespace conditions, strict-prefix replay, RUP-complete audits, the carrier-independent audit criterion, the self-auditing RUP-CDCL instantiation, its append-only deletion-free variant, and the prefix-typed append-only refinement say exactly when proof size inherits the footprint lower bound. & Proof-system bridge: proof-complexity applications require an auditable extractor, not a rhetorical analogy. \\
\bottomrule
\end{tabularx}
\caption{Novelty-delta accounting for the cohered result package.}
\label{tab:novelty-delta-accounting}
\end{table}

Thus the changed understanding is a single statement rather than five separate
payoffs: completion-sound evidence that a generated labeling is hard must either
expose a large local footprint, or supply an auditable proof-system mechanism
that computes the same footprint from the proof artifact.  The surrounding
nonclaims are part of this thesis.  The paper is not a new general hardness
magnification theorem, not a Partial MCSP theorem, and not a Frege, Extended
Frege, or IPS lower bound.  Its contribution is the local obstruction and the
exact audit target needed before such proof-system consequences can be claimed.

\subsection{Sorting networks and static lookup}

Sorting networks originate with Batcher~\cite{Batcher68}; Ajtai, Koml\'os, and Szemer\'edi gave asymptotically optimal $O(N\log N)$ comparator networks~\cite{AKS83}.  We use sorting only as a static nonuniform dictionary primitive.  The construction is elementary, but it is quantitatively decisive: replacing pairwise comparison cost $MK\ell$ by nearly linear batch-plus-dictionary cost upgrades the formula lower bound from roughly $s/(M\ell)$ to $s/\polylog(s)$ in the relevant regime.

\subsection{Natural proofs and constructive lower bounds}

Razborov and Rudich formalized the natural-proofs barrier~\cite{RR97}.  Carmosino, Impagliazzo, Kabanets, and Kolokolova showed that natural lower bounds yield learning algorithms for the same circuit class~\cite{CIKK16}; later work developed broader connections between constructive lower bounds, algorithms, and hardness magnification.  Our functional naturalization proposition is a direct substitution observation.  It is included only to distinguish the global dense regime from the sparse certificate regimes that drive the main results.

\subsection{Proof complexity motivation}

A long-term motivation is to understand whether lower-bound proofs in strong systems must expose local semantic information.  Kraj\'i\v{c}ek's implicit proofs~\cite{Krajicek04} and recent interactive-protocol connections between proof and circuit complexity~\cite{AKPS24} illustrate that strong proof systems may use compressed or nonlocal witnesses.  Therefore no conclusion about such systems follows merely from an occurrence-query lower bound.  \Cref{thm:displayed-atom-extraction} gives the restricted positive case: if a proof literally displays faithful formula atoms and is completion-sound, the formula-footprint theorem lower-bounds the displayed local content.  \Cref{thm:assumption-resolution,cor:resolution-unit-core,thm:resolution-local-clause-support,thm:support-local-aux-resolution,thm:support-local-extension-resolution} instantiate this bridge for resolution refutations of an easy-completion CNF when the proof explicitly lists faithful label assumptions, when a full-assignment refutation uses faithful label-unit leaves, when the proof uses faithful clauses over formula-label variables, when auxiliary constraints are support-local under completions, or when extension definitions read only a fixed support of formula-label variables.  \Cref{thm:dependency-cone-extension-normalizer,cor:canonical-syntactic-cone,cor:final-backward-slice,cor:slice-explicit-er,cor:displayed-er-length,cor:compressed-er-expansion,cor:macro-compressed-er,cor:support-declared-macro-er,cor:support-audited-lrat,cor:provenance-typed-lrat,cor:namespace-typed-lrat,prop:strict-prefix-qrep,cor:trace-complete-dpll,cor:antecedent-complete-cdcl,prop:cdcl-antecedent-erasure-trichotomy,cor:rup-complete-cdcl,cor:canonical-namespace-rup-cdcl,cor:self-auditing-rup-cdcl,cor:append-only-self-auditing-rup-cdcl,cor:prefix-typed-append-only-rup-cdcl,cor:intrinsic-locality-prefix-typed-rup-cdcl,prop:support-audit-necessity,prop:support-erasure-indistinguishability,prop:plain-lrat-extraction-trichotomy,cor:cone-budget-extension,cor:bounded-depth-extension,cor:carrier-independent-audit-upgrade} move the extension, compressed-replay, clausal-trace, DPLL-log, learned-clause, and carrier-independent audit routes one step closer to an auditable proof artifact: nested extension variables are allowed, the needed dependency cones can be declared or read canonically from the acyclic extension-definition graph, the final backward slice must expose near-threshold label mass, slice-explicit ER-style representations must pay for that boundary, ordinary displayed acyclic extension length pays for it without a separate slice-size convention, expansion-bounded compressed representations inherit a decompression tradeoff, macro libraries inherit a skeleton-plus-reuse product tradeoff, sound support declarations, support-audited LRAT traces, provenance-typed LRAT/DRAT-with-LRAT-certificate traces, and namespace-typed LRAT antecedent-closure traces inherit support-budget lower bounds, strict-prefix quotient representatives replay from already charged same-theory forests, trace-complete DPLL logs inherit the controller-complete adaptive audit when the complete controller and local witnesses are present, antecedent-complete CDCL logs inherit it when every used learned clause has acyclic antecedent provenance, the CDCL antecedent-erasure trichotomy shows that provenance-free learned-clause syntax must add an equivalent invariant, pay conservative support, or be unsound, RUP-complete CDCL logs give the positive non-plain route when every used learned clause has a checker-visible RUP propagation certificate over the known live database and parsed formula-label support, the canonical-namespace RUP-CDCL standardizer replaces trusted parsing by a checkable label-interval manifest and live-database reconstruction, the self-auditing RUP-CDCL corollary internalizes that parser and the retained support into the proof object itself, the append-only self-auditing corollary removes the retained deletion-support records when the trace forbids deletions, the prefix-typed append-only corollary replaces the proof-internal decoder by a public syntactic namespace for formula-label variables, and the intrinsic-locality prefix-typed corollary replaces separate auxiliary witnesses by fixed public completion-stable auxiliary rules.  The support-audit and support-erasure separators explain why unaudited clausal traces do not automatically supply that support, the plain LRAT trichotomy says a raw-trace extractor must add provenance, pay worst-case support, or be unsound on a same-skeleton semantic lift, the final-interface size, cone width, and bounded-fanin extension depth obey explicit tradeoffs, and a system-intrinsic audit extractor would be enough to remove the carrier-specific manifest from the proof-system bridge.  \Cref{thm:support-local-normalization-transfer} packages the general implication: once a proof system has a support-local normalizer with a controlled support budget, proof-size lower bounds follow.  A stronger application to Frege, Extended Frege, IPS, arbitrary compressed proofs, unrestricted nonlocal extended resolution, raw untyped LRAT/DRAT, raw RUP/DRUP without typed support fields, deletion-bearing traces without retained support, delayed quotient certificates, CDCL traces with hidden learned-clause provenance and no RUP-complete support audit, or another proof system would still need to supply that normalizer.

\subsection{The normalization target for proof-system applications}\label{subsec:normalization-target}

The required proof-system bridge can be stated independently of any particular system.  A \emph{controlled local normalization theorem} for a proof system \(\mathcal P\) would consist of:
\begin{enumerate}
  \item a family of lower-bound statements \(L_n\) asserting that the induced carrier lies in one of the hardness predicates studied here;
  \item an extractor mapping every accepted \(\mathcal P\)-proof of \(L_n\) to one of the following local objects:
  \[
    \text{formula labels},\qquad
    \text{private occurrence rows},\qquad
    \text{equality-quotient labels};
  \]
  \item a faithfulness invariant saying that extracted labels agree with the intended carrier labeling;
  \item a completion-soundness invariant saying that every admissible completion of the unextracted labels still satisfies the relevant hardness predicate;
  \item an overhead bound relating proof size to extracted local size.
\end{enumerate}
If the extractor returns formula labels, \Cref{thm:displayed-atom-extraction,thm:assumption-resolution,cor:resolution-unit-core,thm:resolution-local-clause-support,thm:support-local-aux-resolution,thm:support-local-extension-resolution,thm:dependency-cone-extension-normalizer,cor:canonical-syntactic-cone,cor:final-backward-slice,cor:slice-explicit-er,cor:displayed-er-length,cor:compressed-er-expansion,cor:macro-compressed-er,cor:support-declared-macro-er,cor:support-audited-lrat,cor:provenance-typed-lrat,cor:namespace-typed-lrat,prop:strict-prefix-qrep,cor:trace-complete-dpll,cor:antecedent-complete-cdcl,cor:rup-complete-cdcl,cor:canonical-namespace-rup-cdcl,cor:self-auditing-rup-cdcl,cor:append-only-self-auditing-rup-cdcl,cor:prefix-typed-append-only-rup-cdcl,cor:intrinsic-locality-prefix-typed-rup-cdcl,cor:cone-budget-extension,cor:bounded-depth-extension,thm:support-local-normalization-transfer} and \Cref{thm:formula-footprint} give the footprint lower bound.  If it returns private occurrence rows, \Cref{thm:row-forcing} applies.  If it returns equality-quotient labels, \Cref{thm:quotient-distance} applies.  \Cref{prop:support-audit-necessity} records a boundary case: checked clausal unsatisfiability does not by itself return a support-local formula-label object.  \Cref{prop:support-erasure-indistinguishability} records the sharper provenance boundary: untyped LRAT syntax does not determine whether a variable is auxiliary or label support.  \Cref{prop:plain-lrat-extraction-trichotomy} packages these two separators into the operational trichotomy for plain trace extractors: add provenance, pay role-agnostic worst-case support, or fail on a same-skeleton semantic lift.  \Cref{cor:provenance-typed-lrat} records the positive trusted-provenance route: typed roles, block-local supports, and compatible witnesses do supply the support audit.  \Cref{cor:namespace-typed-lrat} records the lighter namespace route: when formula-label variables are syntactically reserved, the LRAT antecedent closure of the empty clause identifies the used label support of witnessed auxiliary clauses.  \Cref{prop:strict-prefix-qrep} records the corresponding boundary for quotient compression: a representative use is replayable only when its same-theory certificate is already charged before use.  \Cref{cor:trace-complete-dpll} records the recognizable DPLL route: a decision log supplies the adaptive audit only when it logs the complete finite controller, replay path, and local clause or propagation witnesses.  \Cref{cor:antecedent-complete-cdcl} records the learned-clause route: a CDCL log supplies the audit only when every used learned clause has acyclic antecedent provenance down to local support.  \Cref{prop:cdcl-antecedent-erasure-trichotomy} records the provenance-free CDCL boundary: erasing learned-clause antecedents forces added provenance or a non-plain invariant, conservative support, or unsoundness on a same-skeleton lift.  \Cref{cor:rup-complete-cdcl} records the positive non-plain CDCL route: RUP propagation certificates over a known live database supply the audit when their formula-label support is parsed and retained.  \Cref{cor:canonical-namespace-rup-cdcl} records the checkable manifest route: a reserved label interval and live-database log make that parser deterministic.  \Cref{cor:self-auditing-rup-cdcl} records the carrier-independent positive route: the label decoder, live-database reconstruction, RUP certificates, and retained support are all part of the proof artifact.  \Cref{cor:append-only-self-auditing-rup-cdcl} records the deletion-free refinement: if the live database is append-only, the same audit needs no retained deleted-clause support records.  \Cref{cor:prefix-typed-append-only-rup-cdcl} records the typed proof-system refinement: if formula-label variables are public syntax, the same audit needs no proof-internal label decoder.  \Cref{prop:prefix-typed-rup-support-witness-necessity} records the matching boundary: public label syntax and append-only RUP checking still do not certify completion-local auxiliary support without a witness or equivalent intrinsic invariant.  \Cref{cor:intrinsic-locality-prefix-typed-rup-cdcl} records the positive intrinsic route: fixed public auxiliary rules with completion-stable support sets can generate those witnesses as ordinary proof-system syntax.  This is the exact point at which stronger proof-system theorems would enter the paper; beyond the explicit-local, support-local auxiliary, support-local extension-resolution, dependency-cone extension-normalizer, canonical syntactic-cone normalizer, final backward-slice lower bound, slice-explicit ER lower bound, displayed acyclic ER length lower bound, expansion-bounded compressed ER transfer, macro-compressed ER reuse tradeoff, support-declared macro ER normalizer, support-audited LRAT trace normalizer, provenance-typed LRAT support audit, namespace-typed LRAT antecedent-closure normalizer, strict-prefix Qrep replay normalizer, trace-complete DPLL decision-log normalizer, antecedent-complete CDCL learned-clause normalizer, CDCL antecedent-erasure trichotomy, RUP-complete CDCL support normalizer, canonical-namespace RUP-CDCL standardizer, self-auditing RUP-CDCL instantiation, append-only self-auditing RUP-CDCL instantiation, prefix-typed append-only RUP-CDCL instantiation, intrinsic-locality prefix-typed RUP-CDCL instantiation, prefix-typed RUP support-witness necessity, support-audit necessity separator, support-erasure indistinguishability separator, plain LRAT extraction trichotomy, cone-budget, and bounded-depth cases, none is assumed here.

{\tiny
\begin{longtable}{@{}p{0.23\textwidth}p{0.27\textwidth}p{0.44\textwidth}@{}}
\caption{Normalization targets for proof-system applications.  The explicit-local, support-local auxiliary, support-local extension-resolution, dependency-cone extension, canonical syntactic-cone, final backward-slice, slice-explicit ER, displayed acyclic ER, expansion-bounded compressed ER, macro-compressed ER, support-declared macro ER, support-audited LRAT, provenance-typed LRAT, namespace-typed LRAT, strict-prefix Qrep replay, trace-complete DPLL decision-log, antecedent-complete CDCL learned-clause, CDCL antecedent-erasure trichotomy, RUP-complete CDCL support normalizer, canonical-namespace RUP-CDCL standardizer, self-auditing RUP-CDCL instantiation, append-only self-auditing RUP-CDCL instantiation, prefix-typed append-only RUP-CDCL instantiation, intrinsic-locality prefix-typed RUP-CDCL instantiation, prefix-typed RUP support-witness necessity, support-audit necessity, support-erasure indistinguishability, plain LRAT extraction trichotomy, semantic-relation, acyclic OR-frontier, quotiented OR-DAG frontier, acyclic Boolean frontier, visible adaptive decision-tree frontier, visible adaptive branching-program frontier, hidden adaptive transcript boundary, cone-budget, and bounded-depth routes are proved here, and the support-local normalizer row gives the reusable transfer theorem.  Future stronger proof-complexity theorems must instantiate one of these local extraction routes with controlled overhead, add provenance, namespaces, controller-complete logs, learned-clause antecedent DAGs, RUP-complete support traces, canonical label-interval manifests, self-auditing proof syntax, append-only audit structure, public typed label syntax plus intrinsic support-locality, or semantic invariants, or explain why the proof system is genuinely nonlocal.}
\label{tab:normalization-targets}\\
\toprule
Normalization output & Applicable theorem & What must be proved about the proof system \\
\midrule
\endfirsthead
\toprule
Normalization output & Applicable theorem & What must be proved about the proof system \\
\midrule
\endhead
Displayed formula atoms & \Cref{thm:displayed-atom-extraction,thm:formula-footprint} & The proof exposes faithful formula labels and remains sound under all completions of unexposed formula labels. \\
Explicit resolution assumptions & \Cref{thm:assumption-resolution,thm:formula-footprint} & A resolution refutation of an easy-completion CNF uses faithful formula-label unit assumptions, and those assumptions are part of the measured proof object. \\
Full-assignment resolution unit core & \Cref{cor:resolution-unit-core,thm:formula-footprint} & A resolution refutation using the full hard labeling must use many distinct faithful formula-label unit leaves of each polarity. \\
Local resolution clauses & \Cref{thm:resolution-local-clause-support,thm:formula-footprint} & A resolution refutation using faithful clauses over formula-label variables must expose many distinct labels of each polarity in the clause support. \\
Support-local auxiliary resolution & \Cref{thm:support-local-aux-resolution,thm:formula-footprint} & A resolution refutation with auxiliary variables must have auxiliary constraints satisfiable under every completion preserving a fixed label support. \\
Support-local extension resolution & \Cref{thm:support-local-extension-resolution,thm:formula-footprint} & Extension definitions may compute from a fixed formula-label support, and final extension clauses must remain satisfied whenever that support agrees with the hard labeling. \\
Dependency-cone extension normalizer & \Cref{thm:dependency-cone-extension-normalizer,thm:formula-footprint} & A cone-annotated extension proof must expose a syntactic upper bound on the formula labels each final extension variable depends on. \\
Canonical syntactic-cone normalizer & \Cref{cor:canonical-syntactic-cone,thm:formula-footprint} & An acyclic extension proof has its final formula-label support extracted automatically from the extension-definition graph. \\
Final backward-slice lower bound & \Cref{cor:final-backward-slice,thm:formula-footprint} & The final backward extension slice must contain near-threshold many distinct formula-label inputs of each polarity. \\
Slice-explicit ER lower bound & \Cref{def:slice-explicit-er,cor:slice-explicit-er,thm:formula-footprint} & A slice-explicit extension refutation must charge that final formula-label boundary in its proof-size measure. \\
Displayed acyclic ER length lower bound & \Cref{cor:displayed-er-length,cor:final-backward-slice,thm:formula-footprint} & If the acyclic extension definitions and final clauses are displayed, ordinary literal length already charges every formula-label input in the final backward slice. \\
Expansion-bounded compressed ER transfer & \Cref{cor:compressed-er-expansion,cor:displayed-er-length} & If a compressed proof expands to displayed acyclic ER with expansion length \(B(|\rho|)\), then \(B(|\rho|)\) must be near-threshold large; polynomial expansion gives a polynomial-root lower bound on compressed length. \\
Macro-compressed ER reuse tradeoff & \Cref{cor:macro-compressed-er,cor:displayed-er-length} & A displayed skeleton of length \(S\), macro library of length \(D\), and \(r\) invocations must satisfy \(S+rD\ge 2n^{a-o(1)}\); ordinary macro accounting gives \(L\ge n^{a/2-o(1)}\). \\
Support-declared macro ER normalizer & \Cref{cor:support-declared-macro-er,thm:support-local-aux-resolution} & A compressed macro proof with sound declared support \(S_\rho\) must satisfy \(|S_\rho|\ge 2n^{a-o(1)}\) near threshold; any support budget \(B(|\rho|)\) is therefore near-threshold large. \\
Support-audited LRAT trace normalizer & \Cref{cor:support-audited-lrat,thm:support-local-aux-resolution} & An LRAT-style checked clausal trace with a sound support audit must expose near-threshold formula-label support; unaudited LRAT/DRAT remains outside the theorem. \\
Provenance-typed LRAT support audit & \Cref{cor:provenance-typed-lrat,cor:support-audited-lrat} & A typed variable-role map plus block-local clause provenance and compatible auxiliary witnesses supplies the sound support audit; raw untyped LRAT/DRAT remains outside the theorem. \\
Namespace-typed LRAT antecedent closure & \Cref{cor:namespace-typed-lrat,cor:support-audited-lrat} & A trusted formula-label variable namespace plus the LRAT antecedent closure of the empty clause supplies the used label support for witnessed auxiliary clauses; erasing the namespace returns to the raw-trace boundary. \\
Strict-prefix Qrep replay normalizer & \Cref{prop:strict-prefix-qrep} & Quotient representative uses must replay from already charged same-theory certificate forests; delayed or hidden quotient certificates are private-row falsifiers. \\
Trace-complete DPLL decision-log normalizer & \Cref{cor:trace-complete-dpll,prop:controller-complete-adaptive-audit-normalizer} & A DPLL-style log must expose the complete finite decision controller, hard-tuple replay path, terminal value, and local clause or propagation witnesses certifying equivalence on the completion cylinder; partial transcripts and unaudited learned clauses are not normalized. \\
Antecedent-complete CDCL learned-clause normalizer & \Cref{cor:antecedent-complete-cdcl,cor:trace-complete-dpll} & A CDCL-style log must expose every learned clause used on the replay path through an acyclic antecedent DAG whose leaves are original local clauses, decisions, unit propagations, or lower-rank learned clauses; hidden learned-clause provenance is not normalized. \\
CDCL antecedent-erasure trichotomy & \Cref{prop:cdcl-antecedent-erasure-trichotomy,prop:support-erasure-indistinguishability,prop:support-audit-necessity} & A CDCL log that erases learned-clause antecedents and role provenance must add equivalent provenance or a non-plain semantic invariant, pay conservative role-agnostic support, or become unsound on a same-skeleton semantic lift. \\
RUP-complete CDCL support normalizer & \Cref{cor:rup-complete-cdcl,prop:controller-complete-adaptive-audit-normalizer} & A CDCL/DRUP-style log must expose checker-visible RUP propagation certificates over a known live clause database, parsed formula-label namespace or role information, and retained support for deleted-clause consequences; raw RUP/DRUP syntax without support parsing is not normalized. \\
Canonical-namespace RUP-CDCL standardizer & \Cref{cor:canonical-namespace-rup-cdcl,cor:rup-complete-cdcl} & A deterministic encoding manifest must reserve a formula-label variable interval, map it injectively to frontier coordinates, reconstruct the live clause database, and retain deleted-clause certificate supports; unmanifested DIMACS/RUP/DRUP syntax is not normalized. \\
Self-auditing RUP-CDCL trace & \Cref{cor:self-auditing-rup-cdcl,cor:carrier-independent-audit-upgrade} & The proof syntax itself contains the checked label decoder, live-database reconstruction, RUP propagation witnesses, and retained deleted-clause support; the support audit is computed from the proof artifact rather than a side manifest. \\
Append-only self-auditing RUP-CDCL trace & \Cref{cor:append-only-self-auditing-rup-cdcl,cor:self-auditing-rup-cdcl,cor:carrier-independent-audit-upgrade} & If the self-auditing trace is deletion-free, the monotone clause-addition log and RUP certificates determine the touched formula-label support; no retained deleted-clause support records are needed. \\
Prefix-typed append-only RUP-CDCL trace & \Cref{cor:prefix-typed-append-only-rup-cdcl,cor:append-only-self-auditing-rup-cdcl,cor:carrier-independent-audit-upgrade} & If formula-label variables are public syntax in an append-only RUP-CDCL system, the same monotone scan determines support without a proof-internal label decoder or retained deleted-clause support records. \\
Prefix-typed RUP support-witness necessity & \Cref{prop:prefix-typed-rup-support-witness-necessity} & Public label syntax and append-only RUP checking do not certify completion-local auxiliary support by themselves; the system must provide a support-local witness or an equivalent intrinsic invariant. \\
Intrinsic-locality prefix-typed RUP-CDCL trace & \Cref{cor:intrinsic-locality-prefix-typed-rup-cdcl,cor:carrier-independent-audit-upgrade} & If touched auxiliary clauses are introduced only by fixed public completion-stable rules with checker-computable support sets, the separate auxiliary witness field is replaced by ordinary proof-system syntax. \\
Support-audit necessity separator & \Cref{prop:support-audit-necessity} & A linear-size checked clausal refutation can still require all fixed labels in its sound support audit; checked clausal unsatisfiability alone does not certify completion-local support. \\
Support-erasure indistinguishability separator & \Cref{prop:support-erasure-indistinguishability} & The same untyped LRAT skeleton can be lifted with empty formula-label support or with full label support; role/provenance metadata is part of the support audit. \\
Plain LRAT extraction trichotomy & \Cref{prop:plain-lrat-extraction-trichotomy} & A raw-trace extractor must add provenance, be role-agnostically conservative and pay the full support on the separator skeleton, or become unsound on a same-skeleton semantic lift. \\
Cone-budget extension lower bound & \Cref{cor:cone-budget-extension} & The final-interface variable count times the maximum declared cone size must be near-threshold large. \\
Bounded-depth extension tradeoff & \Cref{cor:bounded-depth-extension} & For bounded-fanin extensions, the final-interface variable count times \(r^h\) must be near-threshold large. \\
Support-local normalizer & \Cref{thm:support-local-normalization-transfer,thm:formula-footprint} & A proof system supplies a support-local CNF and a support budget \(B(|\pi|)\) for every accepted proof. \\
Carrier-independent audit extractor & \Cref{cor:carrier-independent-audit-upgrade,thm:support-local-normalization-transfer} & The proof-system artifact itself computes support, support-locality evidence, and budget without a carrier-specific namespace manifest, hidden role advice, or semantic oracle. \\
Private occurrence rows & \Cref{thm:row-forcing,cor:full-row} & The proof exposes occurrence labels, and unexposed positions can be completed independently without invalidating soundness. \\
Equality quotient labels & \Cref{thm:quotient-distance} & The proof's semantic equalities induce a sound quotient, and residual labels remain freely completable in the quotient affine space. \\
Nonlinear relation instance & \Cref{prop:semantic-relation-support} & The proof's declared support contains a standard membership certificate for each finite semantic relation instance it imposes. \\
Acyclic OR-tree frontier & \Cref{prop:acyclic-or-tree-residual} & After internal OR nodes are eliminated, root \(0\) forces the root and all frontier leaves, while root \(1\) requires the root plus one true leaf. \\
Quotiented OR-DAG frontier & \Cref{prop:or-dag-frontier-residual} & After internal OR nodes are eliminated and repeated reachable frontier formulas are soundly quotiented, root \(0\) forces the root and every distinct reachable frontier class, while root \(1\) requires the root plus one true reachable frontier class. \\
Acyclic Boolean frontier & \Cref{prop:acyclic-boolean-frontier-residual} & After an acyclic Boolean semantic frontier is exposed, support-locality is the root coordinate plus an ordinary pointwise certificate for the induced frontier function. \\
Visible adaptive frontier & \Cref{prop:visible-adaptive-frontier-residual} & After a finite adaptive decision tree over frontier coordinates is exposed, support-locality is the root coordinate plus an ordinary pointwise certificate for the induced decision-tree function; the actual transcript path is sufficient support. \\
Visible adaptive branching-program frontier & \Cref{prop:visible-adaptive-branching-frontier-residual} & After a finite acyclic branching program over frontier coordinates is exposed, support-locality is the root coordinate plus an ordinary pointwise certificate for the induced branching-program function; the actual computation path is sufficient support. \\
Hidden adaptive transcript summary & \Cref{prop:hidden-adaptive-transcript-insufficiency} & A small reported adaptive transcript is not a support normalizer unless completeness of the controller/path or enough dependency support is certified; otherwise the same transcript can hide all-frontier OR support. \\
Controller-complete adaptive audit & \Cref{prop:controller-complete-adaptive-audit-normalizer} & A hidden adaptive transcript becomes a support normalizer when the proof object exposes a complete acyclic controller, replay path, and local equivalence certificate for the completion cylinder. \\
Nonlocal or compressed witness & None from this paper & The proof system escapes the present local completion models; a different argument is required. \\
\bottomrule
\end{longtable}
}

\subsection{Open problems}

\begin{openproblem}[Expected randomized and interactive footprints]\label{op:randomized}
\Cref{thm:accepting-union} extends the certificate lower bounds to the union of labels queried on accepting bounded-error executions.  Prove, or refute, analogous lower bounds for expected footprint, transcript support, communication-query product, or information revealed by public-coin or interactive verifiers with bounded soundness error.  Such a result would require a distributional or information-theoretic completion argument, not merely the deterministic certificate reduction used here.
\end{openproblem}

\begin{openproblem}[A concrete proof-system instantiation]\label{op:proof-system}
Identify a natural proof or certificate system for which every short proof of the carrier lower-bound statement can be normalized into one of the completion-sound models in \Cref{tab:normalization-targets}.  \Cref{thm:displayed-atom-extraction} handles the toy case where formula atoms are already displayed faithfully, while \Cref{thm:assumption-resolution,cor:resolution-unit-core,thm:resolution-local-clause-support,thm:support-local-aux-resolution,thm:support-local-extension-resolution,thm:dependency-cone-extension-normalizer,cor:cone-budget-extension,cor:bounded-depth-extension} handle restricted resolution cases where faithful label information is explicit, auxiliary constraints are support-local, extension definitions read only a fixed label support, or extension variables expose checked dependency cones with bounded final cone budget or bounded-fanin depth.  \Cref{prop:strict-prefix-qrep} handles the quotient-compression case only when every representative use has a prior charged same-theory certificate forest.  \Cref{prop:support-audit-necessity} shows why checked clausal unsatisfiability is not enough by itself: the normalizer must certify completion-local support, not merely a refutation.  \Cref{prop:support-erasure-indistinguishability} adds that untyped trace syntax does not determine the formula-label support without variable-role provenance.  \Cref{prop:plain-lrat-extraction-trichotomy} turns this into the plain-trace alternative: add provenance, pay role-agnostic worst-case support, or lose soundness on an indistinguishable semantic lift.  \Cref{prop:hidden-adaptive-transcript-insufficiency} adds the adaptive analogue: a reported transcript summary cannot replace a certified complete controller/path.  \Cref{prop:controller-complete-adaptive-audit-normalizer} gives the matching positive audited-hidden-adaptive condition: expose the acyclic controller, replay path, and local equivalence certificate for the completion cylinder.  \Cref{cor:trace-complete-dpll} instantiates that condition for DPLL-style decision logs when the whole finite controller and the local clause or propagation witnesses are part of the proof object.  \Cref{cor:antecedent-complete-cdcl} extends the instantiation to learned clauses only when every used learned clause carries an acyclic antecedent DAG down to audited local support.  \Cref{cor:provenance-typed-lrat} proves the positive trusted-provenance case when role maps, block supports, and witness compatibility are checked.  \Cref{thm:support-local-normalization-transfer} shows that the lower bound follows automatically once such a normalizer is available.  \Cref{cor:carrier-independent-audit-upgrade} turns the same open problem into an extractor target: compute support, support-locality evidence, and budget from the proof-system artifact itself, without a carrier-specific manifest or hidden role advice.  \Cref{cor:prefix-typed-append-only-rup-cdcl} shows the current strongest RUP-style bridge: in an append-only trace, public typed label syntax replaces the proof-internal decoder.  \Cref{prop:prefix-typed-rup-support-witness-necessity} shows why the remaining support-local auxiliary witness cannot be removed by that syntax alone.  \Cref{cor:intrinsic-locality-prefix-typed-rup-cdcl} supplies the matching positive condition: fixed public completion-stable auxiliary rules make the witness an intrinsic proof-system invariant.  The open problem is now sharper: instantiate that intrinsic rule discipline in a natural solver/proof format, handle deletion-bearing or untyped traces, go beyond typed RUP-CDCL to a stronger or less local proof system, or identify a non-plain LRAT/DRAT/CDCL invariant that supplies the missing semantic support without merely assuming the provenance object.
\end{openproblem}

\begin{openproblem}[Beyond local semantic-relation certificates]\label{op:nonlinear}
\Cref{prop:semantic-relation-support} gives the local certificate cost of one nonlinear semantic relation, \Cref{prop:acyclic-or-tree-residual} gives the exact eliminated-frontier certificate for acyclic OR self-reduction trees, \Cref{prop:or-dag-frontier-residual} gives the corresponding result for finite acyclic OR DAGs after sound frontier quotienting, \Cref{prop:acyclic-boolean-frontier-residual} handles exposed finite acyclic Boolean frontier circuits by ordinary pointwise certificate complexity, \Cref{prop:visible-adaptive-frontier-residual} handles exposed finite adaptive decision trees, and \Cref{prop:visible-adaptive-branching-frontier-residual} handles exposed finite acyclic adaptive branching programs.  \Cref{prop:hidden-adaptive-transcript-insufficiency} shows that hidden adaptive summaries do not give such a theorem for free, while \Cref{prop:controller-complete-adaptive-audit-normalizer} records the sufficient audit when the hidden controller is locally certified against an exposed acyclic controller.  Develop a residual-completion theorem for whole systems of nonlinear semantic constraints beyond exposed acyclic or visible-adaptive frontiers, beginning with bounded SAT self-reduction
\[
  Y(u)=Y(u|_{x_i=0})\vee Y(u|_{x_i=1}).
\]
The equality quotient gives an affine completion space.  Self-reduction produces a constrained Boolean variety.  The local certificate costs in \Cref{prop:semantic-relation-support} identify which labels an audit must expose one relation at a time, the OR-frontier theorems solve the collapsed pure-OR tree and shared-DAG relations exactly, the Boolean-frontier theorem reduces exposed acyclic semantics to standard certificate complexity, and the visible-adaptive theorems show that logged adaptive control is harmless only when the whole decision tree or branching program is audited.  The hidden-transcript separator shows that a small reported transcript without this completeness certificate can hide all-frontier OR support; the controller-complete audit theorem identifies the positive certificate that would remove this ambiguity on a completion cylinder.  The remaining problem is the global analogue of residual distance for hidden-adaptive, cyclic, learned-clause, or proof-system-hidden SAT self-reduction systems.
\end{openproblem}

\begin{openproblem}[An explicit hard SAT-labeled carrier]\label{op:source}
Construct an explicit batched generator for which the actual SAT parity carrier has an unrestricted circuit lower bound above the substitution threshold.  Any such result must pass the formula-diversity, row-robustness, and quotient-distance tests proved here, as well as the reverse short-SAT burden in \Cref{prop:reverse}.
\end{openproblem}

\section{Conclusion}

Label-induced source functions provide a simple laboratory for separating three kinds of local information.  Distinct-formula certificates are limited by batched static lookup and therefore require near-threshold footprints of both polarities.  Occurrence certificates are limited by sparse-row patching and therefore must force many rows; the aggregator's minimum forcing numbers determine the query cost, with parity requiring complete rows on both sides.  Equality constraints replace private occurrences by an affine quotient, but completion soundness still demands large residual Hamming distance from the constant functions.

For polylogarithmic SAT parity carriers, these results sit exactly at a magnification threshold strong enough to imply $\SAT\notin\Ppoly$ if the source lower bound could be proved.  The paper does not supply that lower bound.  It factors out a smaller, standard question: the nondeterministic query complexity and completion geometry of the induced circuit-hardness predicate.  The displayed-atom, explicit-local resolution, support-local auxiliary-resolution, support-local extension-resolution, dependency-cone extension, canonical syntactic-cone, final backward-slice, slice-explicit ER, displayed acyclic ER, expansion-bounded compressed ER, macro-compressed ER, support-declared macro ER, support-audited LRAT, provenance-typed LRAT, namespace-typed LRAT, strict-prefix Qrep replay, trace-complete DPLL decision-log, antecedent-complete CDCL learned-clause, CDCL antecedent-erasure trichotomy, RUP-complete CDCL support normalizer, canonical-namespace RUP-CDCL standardizer, carrier-independent audit upgrade criterion, self-auditing RUP-CDCL instantiation, append-only self-auditing RUP-CDCL instantiation, prefix-typed append-only RUP-CDCL instantiation, prefix-typed RUP support-witness necessity, intrinsic-locality prefix-typed RUP-CDCL instantiation, support-audit necessity, support-erasure indistinguishability, plain LRAT extraction trichotomy, semantic-relation support, acyclic OR-frontier, quotiented OR-DAG frontier, acyclic Boolean-frontier, visible adaptive frontier, visible adaptive branching-program frontier, hidden adaptive transcript insufficiency, controller-complete adaptive audit, cone-budget, bounded-depth, and normalization-transfer theorems show how this cohered package becomes a proof-size statement once local semantic atoms are faithfully extracted, explicitly used, support-locally represented, syntactically dependency-audited, expanded with controlled overhead, compressed only through visible macro reuse, summarized by a sound support declaration, certified by a checker-visible support audit, certified by typed block provenance or a trusted label-variable namespace plus LRAT antecedent closure, replayed from a strict-prefix quotient dictionary, logged as a trace-complete DPLL controller with local witnesses, logged as an antecedent-complete CDCL learned-clause DAG, accepted as a RUP-complete CDCL/DRUP trace whose propagation certificates expose parsed formula-label support over the live database, standardized by a canonical label-variable interval manifest and live-database log, internalized as a self-auditing trace whose proof object carries its decoder and support audit, simplified in the deletion-free case to an append-only audit with no retained deleted-clause support records, simplified further in the typed case to a public-label namespace with no proof-internal decoder, rejected when public label typing is asked to certify auxiliary support-locality by itself, accepted again when fixed public completion-stable auxiliary rules make that locality intrinsic to the proof syntax, lifted to a carrier-independent proof-system payoff by an intrinsic audit extractor, separated from unaudited clausal refutation soundness, protected against provenance-erased trace ambiguity, forced through the plain-trace trichotomy, checked against local nonlinear semantic relations, collapsed to a pure OR frontier with exact residual certificate cost, exposed as an acyclic Boolean frontier whose support cost is ordinary certificate complexity, logged as a visible adaptive frontier whose transcript path is a certificate, logged as a visible adaptive branching program whose computation path is a certificate, rejected as an incomplete hidden adaptive transcript that can mask all-frontier support, or accepted as an audited hidden-adaptive controller only when the proof object pays for controller completeness and local equivalence.  The next decisive step is to instantiate the intrinsic rule discipline in a natural solver/proof format, handle deletion-bearing or untyped traces, or supply an explicit SAT-labeled source lower bound.

\begin{thebibliography}{99}

\bibitem{AKS83}
M.~Ajtai, J.~Koml\'os, and E.~Szemer\'edi.
\newblock An $O(n\log n)$ sorting network.
\newblock In \emph{Proceedings of the 15th Annual ACM Symposium on Theory of Computing}, pages 1--9, 1983.

\bibitem{AKPS24}
N.~Arteche, E.~Khaniki, J.~Pich, and R.~Santhanam.
\newblock From proof complexity to circuit complexity via interactive protocols.
\newblock \emph{arXiv preprint arXiv:2405.02232}, 2024.

\bibitem{AM25}
A.~Atserias and M.~M\"uller.
\newblock Simple general magnification of circuit lower bounds.
\newblock \emph{arXiv preprint arXiv:2503.24061}, 2025.

\bibitem{AnthonyBCS92}
M.~Anthony, G.~R. Brightwell, D.~A. Cohen, and J.~Shawe-Taylor.
\newblock On exact specification by examples.
\newblock In \emph{Proceedings of the Fifth Annual Workshop on Computational Learning Theory (COLT '92)}, pages 311--318. ACM, 1992.
\newblock doi:10.1145/130385.130420.

\bibitem{Batcher68}
K.~E.~Batcher.
\newblock Sorting networks and their applications.
\newblock In \emph{Proceedings of the 1968 Spring Joint Computer Conference}, volume~32 of \emph{AFIPS Conference Proceedings}, pages 307--314, 1968.

\bibitem{BuhrmanDeWolf02}
H.~Buhrman and R.~de Wolf.
\newblock Complexity measures and decision tree complexity: a survey.
\newblock \emph{Theoretical Computer Science}, 288(1):21--43, 2002.

\bibitem{ChapmanWilliams15}
B.~Chapman and R.~Williams.
\newblock The circuit-input game, natural proofs, and testing circuits with data.
\newblock In \emph{Proceedings of the 6th Innovations in Theoretical Computer Science Conference (ITCS 2015)}, pages 263--270. ACM, 2015.
\newblock doi:10.1145/2688073.2688115.

\bibitem{CHOPRS22}
L.~Chen, S.~Hirahara, I.~C.~Oliveira, J.~Pich, N.~Rajgopal, and R.~Santhanam.
\newblock Beyond natural proofs: hardness magnification and locality.
\newblock \emph{Journal of the ACM}, 69(4):25:1--25:49, 2022.

\bibitem{CIKK16}
M.~L.~Carmosino, R.~Impagliazzo, V.~Kabanets, and A.~Kolokolova.
\newblock Learning algorithms from natural proofs.
\newblock In \emph{Proceedings of the 31st Conference on Computational Complexity}, volume~50 of LIPIcs, pages 10:1--10:24, 2016.

\bibitem{DoliwaFSZ14}
T.~Doliwa, G.~Fan, H.~U. Simon, and S.~Zilles.
\newblock Recursive teaching dimension, VC-dimension and sample compression.
\newblock \emph{Journal of Machine Learning Research}, 15:3107--3131, 2014.
\newblock \url{https://jmlr.org/papers/v15/doliwa14a.html}.

\bibitem{Hirahara22PartialMCSP}
S.~Hirahara.
\newblock NP-hardness of learning programs and partial MCSP.
\newblock In \emph{Proceedings of the 63rd IEEE Symposium on Foundations of Computer Science}, pages 968--979, 2022.
\newblock doi:10.1109/FOCS54457.2022.00095.

\bibitem{Krajicek04}
J.~Kraj\'i\v{c}ek.
\newblock Implicit proofs.
\newblock \emph{Journal of Symbolic Logic}, 69(2):387--397, 2004.

\bibitem{MMW19}
D.~M.~McKay, C.~D.~Murray, and R.~R.~Williams.
\newblock Weak lower bounds on resource-bounded compression imply strong separations of complexity classes.
\newblock In \emph{Proceedings of the 51st Annual ACM Symposium on Theory of Computing}, pages 1215--1225, 2019.

\bibitem{Nisan91}
N.~Nisan.
\newblock CREW PRAMs and decision trees.
\newblock \emph{SIAM Journal on Computing}, 20(6):999--1007, 1991.

\bibitem{OPS21}
I.~C.~Oliveira, J.~Pich, and R.~Santhanam.
\newblock Hardness magnification near state-of-the-art lower bounds.
\newblock \emph{Theory of Computing}, 17(11):1--38, 2021.

\bibitem{OS18}
I.~C.~Oliveira and R.~Santhanam.
\newblock Hardness magnification for natural problems.
\newblock In \emph{Proceedings of the 59th IEEE Symposium on Foundations of Computer Science}, pages 65--76, 2018.

\bibitem{RR97}
A.~A.~Razborov and S.~Rudich.
\newblock Natural proofs.
\newblock \emph{Journal of Computer and System Sciences}, 55(1):24--35, 1997.

\end{thebibliography}

\end{document}
